% Typeset source of the paper. Generated from the assembled draft.
\documentclass[11pt,a4paper]{article}

% ---------- language and type ----------
\usepackage{polyglossia}
\setdefaultlanguage[variant=british]{english}
\usepackage{fontspec}
\usepackage{unicode-math}
\setmainfont{EBGaramond}[
  Extension = .otf, UprightFont = *-Regular, ItalicFont = *-Italic,
  BoldFont = *-SemiBold, BoldItalicFont = *-SemiBoldItalic,
  Numbers = {Lining,Proportional}, Ligatures = TeX]
\setmonofont{Inconsolatazi4}[Extension = .otf, UprightFont = *-Regular, BoldFont = *-Bold, Scale = MatchLowercase, HyphenChar = None]
\setmathfont{Garamond-Math.otf}[StylisticSet={7,9,10}]
\newfontfamily\titleface{Cormorant Garamond SemiBold}[ItalicFont = {Cormorant Garamond Medium Italic}, Numbers = Lining, Ligatures = TeX]
\newfontfamily\labelface{EBGaramond}[Extension = .otf, UprightFont = *-Medium, Numbers = {Lining,Proportional}]
\newfontfamily\groupface{EBGaramond}[Extension = .otf, UprightFont = *-Medium, ItalicFont = *-MediumItalic, Numbers = {Lining,Proportional}]
\newfontfamily\thmheadface{EBGaramond}[Extension = .otf, UprightFont = *-SemiBold]
\newcommand{\pagelabel}[1]{{\labelface\spacedsc{#1}}}
\newcommand{\spacedsc}[1]{{\scshape\addfontfeatures{LetterSpace=8}\MakeLowercase{#1}}}
\DeclareRobustCommand{\scn}[1]{{\upshape\addfontfeatures{Numbers={Lining,Proportional}}\scalebox{0.84}{#1}}}
\DeclareRobustCommand{\thd}[1]{{\scshape\addfontfeatures{LetterSpace=1.5}\MakeLowercase{#1}}}
\DeclareRobustCommand{\acr}[1]{{\addfontfeatures{Letters=UppercaseSmallCaps,LetterSpace=2}#1}}
\frenchspacing
\displaywidowpenalty=10000

% ---------- body size and leading: 12 on 15.2 ----------
\makeatletter
\renewcommand\normalsize{%
  \@setfontsize\normalsize{12}{15.2}%
  \abovedisplayskip 10\p@ \@plus2\p@ \@minus4\p@
  \abovedisplayshortskip 4\p@ \@plus3\p@
  \belowdisplayshortskip 7\p@ \@plus3\p@ \@minus3\p@
  \belowdisplayskip \abovedisplayskip
  \let\@listi\@listI}
\renewcommand\small{\@setfontsize\small{10.4}{13.4}}
\renewcommand\footnotesize{\@setfontsize\footnotesize{9.4}{12}}
\makeatother
\normalsize

% ---------- page ----------
\usepackage[a4paper,textwidth=132mm,top=33mm,bottom=39mm,headsep=9mm,footskip=14mm,headheight=14pt]{geometry}

\usepackage{amsmath,amsthm}
\usepackage{array}
\usepackage{tabularx}
\newcolumntype{L}[1]{>{\raggedright\arraybackslash}p{#1}}
\newcolumntype{Y}{>{\raggedright\arraybackslash}X}
\usepackage{booktabs}
\usepackage{longtable}
\setlength{\heavyrulewidth}{0.6pt}
\setlength{\lightrulewidth}{0.3pt}
\setlength{\extrarowheight}{2pt}
\usepackage{float}
\usepackage[section]{placeins}
\usepackage{xcolor}
\definecolor{oxford}{RGB}{0,33,71}
\definecolor{oxblood}{RGB}{116,30,44}
\definecolor{racing}{RGB}{24,74,52}
\usepackage{caption}
\DeclareCaptionLabelFormat{sclabel}{{\labelface\scshape\MakeLowercase{#1}}~{\labelface\scn{#2}}}
\captionsetup{figurename=Fig.,labelformat=sclabel,labelsep=period,margin=9mm,font=small,skip=7pt}
\usepackage{enumitem}
\usepackage{etoolbox}

% ---------- headings: centred titling face; italic subheadings and paragraph leads; no bold ----------
\usepackage{titlesec}
\usepackage{titletoc}
\setcounter{tocdepth}{1}
\titlecontents{section}[1.9em]{\small}{\contentslabel{1.9em}}{}{\hfill\contentspage}[\vspace{0.5pt}]
\titleformat{\section}[block]{\centering\titleface\fontsize{17pt}{20pt}\selectfont}{\thesection}{0.7em}{}
\titlespacing*{\section}{0pt}{30pt plus 5pt minus 3pt}{13pt}
\titleformat{\subsection}[block]{\normalsize\itshape}{\upshape\thesubsection}{0.75em}{}
\titlespacing*{\subsection}{0pt}{17pt plus 3pt minus 2pt}{6pt}
\titleformat{\subsubsection}[runin]{\normalsize\itshape}{\upshape\thesubsubsection}{0.6em}{}[.]
\titlespacing*{\subsubsection}{0pt}{9pt plus 2pt minus 1pt}{0.6em}
\titleformat{\paragraph}[runin]{\normalsize\itshape}{}{0pt}{}
\titlespacing*{\paragraph}{0pt}{9pt plus 2pt minus 1pt}{0.6em}

\usepackage{graphicx}
\usepackage{tikz}
\usetikzlibrary{positioning,arrows,calc,shadows}
\usepackage{pgfplots}
\pgfplotsset{compat=1.18}
\usepackage{lettrine}
\setcounter{DefaultLines}{2}
\renewcommand{\LettrineTextFont}{\scshape}
\usepackage{indentfirst}

\widowpenalty=10000
\clubpenalty=10000
\usepackage{needspace}
\setlength{\intextsep}{11pt plus 3pt minus 2pt}
\setlength{\floatsep}{11pt plus 3pt minus 2pt}
\setlength{\textfloatsep}{16pt plus 4pt minus 2pt}
\makeatletter\setlength{\@fptop}{0pt}\setlength{\@fpsep}{18pt plus 4pt}\setlength{\@fpbot}{0pt plus 1fil}\makeatother
\renewcommand{\topfraction}{0.85}\renewcommand{\bottomfraction}{0.7}\renewcommand{\textfraction}{0.12}\renewcommand{\floatpagefraction}{0.75}

% ---------- tables: small, lining tabular figures ----------
\renewcommand{\arraystretch}{1.08}
\AtBeginEnvironment{tabular}{\footnotesize\addfontfeatures{Numbers={Lining,Tabular}}}
\AtBeginEnvironment{tabularx}{\footnotesize\addfontfeatures{Numbers={Lining,Tabular}}}
\AtBeginEnvironment{longtable}{\footnotesize\addfontfeatures{Numbers={Lining,Tabular}}}

\usepackage{microtype}
\usepackage[superscript]{cite}

\makeatletter
\renewcommand\@biblabel[1]{#1.}
\renewenvironment{thebibliography}[1]
  {\section*{\refname}\small
   \list{\@biblabel{\@arabic\c@enumiv}}%
        {\settowidth\labelwidth{\@biblabel{#1}}%
         \leftmargin\labelwidth \advance\leftmargin\labelsep
         \setlength\itemsep{3.5pt plus 1pt}\setlength\parsep{0pt}%
         \usecounter{enumiv}\let\p@enumiv\@empty
         \renewcommand\theenumiv{\@arabic\c@enumiv}}%
   \sloppy\clubpenalty4000\widowpenalty4000\sfcode`\.\@m}
  {\def\@noitemerr{\@latex@warning{Empty `thebibliography' environment}}\endlist}
\makeatother

% ---------- abstract and front-matter notes: centred small-capital label over an inset block ----------
\newenvironment{insetblock}[1]{%
  \begin{center}{\footnotesize\pagelabel{#1}}\end{center}\vspace{-5pt}%
  \fontsize{10pt}{13.4pt}\selectfont
  \begin{list}{}{%
    \setlength{\leftmargin}{9mm}\setlength{\rightmargin}{9mm}%
    \setlength{\listparindent}{0pt}\setlength{\parsep}{6pt plus 2pt minus 1pt}%
    \setlength{\topsep}{0pt}\setlength{\partopsep}{0pt}%
  }\item[]\ignorespaces
}{\end{list}}
\renewenvironment{abstract}{\begin{insetblock}{Abstract}}{\end{insetblock}}

\usepackage{hyperref}
% Forbid mid-URL line breaks (a split URL extracts as a broken reference).
\def\UrlBreakPenalty{10000}
\def\UrlBigBreakPenalty{10000}
\pdfstringdefDisableCommands{\def\acr#1{#1}\def\spacedsc#1{#1}\def\thd#1{#1}\def\scn#1{#1}}
\hypersetup{
  colorlinks=true,
  allcolors=black,
  urlcolor=oxford,
  bookmarksnumbered=true,
  bookmarksopen=true,
  bookmarksopenlevel=1,
  pdfdisplaydoctitle=true,
  pdflang=en-GB,
  pdfauthor={K.R. Ryan},
  pdftitle={The Hidden Microstructure of Shared Balance Concentrated Liquidity: A Master Equation for the Dust Ledger and Propagation of Chaos},
  pdfsubject={Concentrated liquidity dust ledger dynamics, master equation and propagation of chaos},
  pdfkeywords={concentrated liquidity, geometric residual, propagation of chaos, interacting particle systems, multi-pool rebalancing, stochastic flow networks, dust ledger, Uniswap V3, Aerodrome, automated market maker, DeFi mechanism design}
}
\urlstyle{same}
\usepackage[bottom]{footmisc}
\renewcommand{\footnotelayout}{\fontsize{8.6pt}{11.4pt}\selectfont}
\setlength{\footnotesep}{9pt}
% Footnotes are marked by symbols, since citations are superscript numerals.
\renewcommand{\thefootnote}{\fnsymbol{footnote}}
\setlength{\skip\footins}{16pt plus 4pt minus 2pt}

% ---------- running head: spaced small capitals, centred; folio at the foot ----------
\usepackage{fancyhdr}
\pagestyle{fancy}
\fancyhf{}
\fancyhead[C]{\footnotesize\spacedsc{The Hidden Microstructure}}
\fancyfoot[C]{\small\thepage}
\renewcommand{\headrulewidth}{0pt}
\fancypagestyle{plain}{\fancyhf{}\fancyfoot[C]{\small\thepage}\renewcommand{\headrulewidth}{0pt}}

% ---------- theorems: small-capital label, name in roman, statement in italic ----------
% Each kind of statement keeps its own counter.
\newtheoremstyle{siphon}{11pt plus 2pt minus 1pt}{10pt plus 2pt minus 1pt}{\itshape}{}{}{}{\newline}%
  {{\thmheadface\scshape\MakeLowercase{\thmname{#1}}\thmnumber{ \scn{#2}}}\thmnote{\upshape\enspace\textperiodcentered\enspace #3}\vadjust{\vskip 3pt}}
\theoremstyle{siphon}
\newtheorem{theorem}{Theorem}
\newtheorem{proposition}{Proposition}
\newtheorem{lemma}{Lemma}
\newtheorem{corollary}{Corollary}
\renewcommand\qedsymbol{\raisebox{0.17ex}{\scalebox{0.68}{\ensuremath{\mdlgblksquare}}}}
\providecommand{\proofqed}{%
  \nobreak\hfil\penalty50\hskip2em\hbox{}\nobreak\hfil%
  \hbox{\qedsymbol}%
  {\parfillskip=0pt \finalhyphendemerits=0 \par}}

\usepackage{mdframed}
\newenvironment{codeblock}{%
  \begin{mdframed}[%
    backgroundcolor=white,%
    linewidth=0.5pt,linecolor=black!30,topline=false,bottomline=false,rightline=false,%
    innertopmargin=6pt,%
    innerbottommargin=6pt,%
    innerleftmargin=12pt,%
    innerrightmargin=12pt,%
    leftmargin=1.25em,%
    rightmargin=0pt,%
    skipabove=6pt,%
    skipbelow=6pt%
  ]\footnotesize%
}{%
  \end{mdframed}%
}

\setlength{\parindent}{1.25em}
\setlength{\parskip}{0pt}
\tolerance=1400
\emergencystretch=1.5em

\begin{document}
\thispagestyle{plain}

\begin{center}
\vspace*{11mm}
{\fontsize{23pt}{28pt}\selectfont\titleface The Hidden Microstructure of\\
Shared Balance Concentrated Liquidity}\\[12pt]
{\fontsize{15.5pt}{20pt}\selectfont\titleface\itshape A Master Equation for the Dust Ledger\\
and Propagation of Chaos}\\[19pt]
\rule{16mm}{0.4pt}\\[15pt]
{\fontsize{12pt}{14pt}\selectfont\spacedsc{K.\,R. Ryan}}\\[7pt]
{\fontsize{9pt}{12pt}\selectfont\itshape May 2026, revised November 2026}
\end{center}
\vspace{-5pt}
\centerline{\scalebox{1.15}{$\therefore$}}
\vspace{7pt}

\begingroup\renewcommand\thefootnote{}\footnotetext{\hypersetup{allcolors=black}Independent researcher \enspace\textperiodcentered\enspace {\scshape orcid}\enspace{\addfontfeatures{Numbers={OldStyle,Proportional}}\href{https://orcid.org/0009-0004-6295-7040}{0009-0004-6295-7040}}}\endgroup
\begingroup\renewcommand\thefootnote{}\footnotetext{This version revises those of May and July 2026. Its empirical sections are rebuilt from on-chain events and several of its statements are restated. A note in the repository lists what differs.}\endgroup

\begin{abstract}
A concentrated liquidity manager that retains rebalance leftovers in a ledger shared across pools couples its positions' mint inputs. This paper derives the ledger's exact update (the Master Equation), including standing balances and swaps, identifies which token remains, and establishes conservation and transfer identities across pools. It extends the geometric residual of \emph{The Geometric Siphon}\cite{ryan2026siphon}. Chain data from one operator on Aerodrome Slipstream (Base) cover $152{,}719$ rebalances across four deployed versions during the periods studied. The update closes to the wei for $150{,}357$; the remaining discrepancies are consistent with rewards or ledger movements that the events do not record. In a model of $N$ pools sharing one token, with direction-dependent random swap errors, the stationary balances of the other tokens of any fixed number of pools become independent at rate $1/N$ as $N$ grows (propagation of chaos). Each balance has an atom at zero with mass given in closed form. The exact contract simulation and chain record support the rule for which token remains where the model's kernel applies. In the contract simulation, a common price move that aligns pool directions prevents decoupling. On chain, the zero indicators of ledger entries are dependent in 14 of 59 pairs that never share a pool. Replays trace the excess of positive dependence to the shared balance at same-range rebalances, which the model omits. The shared ledger makes each leftover part of a portfolio state that can carry value and dependence between positions.
\end{abstract}

\begin{center}
\fontsize{9pt}{12.5pt}\selectfont
\pagelabel{Keywords}\enspace geometric residual \textperiodcentered\ interacting particle systems \textperiodcentered\ multi-pool rebalancing \\
dust ledger \textperiodcentered\ Uniswap V3 \textperiodcentered\ DeFi mechanism design\\[2pt]
\pagelabel{JEL}\enspace C63, C65, D47, G14
\end{center}

\vspace{12pt}
\begin{center}
\begin{minipage}{84mm}
\centerline{\footnotesize\pagelabel{Contents}}\vspace{5pt}
\makeatletter\@starttoc{toc}\makeatother
\end{minipage}
\end{center}
\vspace{6pt}

\section{Introduction}\label{sec:intro}

\lettrine[findent=2pt,nindent=0pt]{A}{ concentrated liquidity} (\acr{CL}) position holds liquidity within a price range. An active manager rebalances it by withdrawing its liquidity, optionally swapping tokens, and minting again into the same or a different range.\footnote{The deployed V9 implementation rebalances when the price reaches a trigger line set inside either range bound by a configured fraction of the range width. The trigger geometry is measured in a follow-on paper.\cite{ryan2026om} The analysis here conditions on realised rebalance events and does not depend on the trigger location.} The V3 mint geometry fixes the required token ratio at a given price inside the new range. Unless the supplied amounts match that ratio, the mint leaves part of one token unused. The standard \texttt{NonfungiblePositionManager} (\acr{NFPM}) draws only the amounts required by the mint. The architecture studied here retains the leftover in a depositor-keyed balance shared across positions, making it available to later mints. This shared balance is the \emph{dust ledger}.

\emph{The Geometric Siphon}\cite{ryan2026siphon}, referred to below as the prior paper, gives the swap-free geometric residual in closed form and approximates the leftover after the manager's swap to first order. This paper derives the leftover exactly in the presence of a standing balance and follows its redistribution through the shared ledger. The retained balance is part of the portfolio state. One position's rebalance changes the inputs available to another.

The \emph{leftover} of a rebalance is the amount of one token that its mint could not use. A \emph{ledger entry} is the amount of one token held in the shared balance. In real arithmetic, after a rebalance at most one of the rebalanced pool's two entries is non-zero. The \emph{execution residual} is the swap amount's proportional deviation from the amount required to eliminate the leftover. For a centred range with no standing balance, its sign determines which token remains (Corollary~\ref{cor:swapcorrected}). A position's \emph{equity} is its deployed value plus its entitlement to the ledger. A \emph{transfer} is the part of the minted value that exceeds the position's equity after deducting swap cost, funded by value owed to other positions (Fig.~\ref{fig:mechanism}).

The empirical study covers four successive versions of a Position Manager built on Aerodrome Slipstream (Base), \acr{PM}~V6 to V9, deployed during March 2026. All four implement the same depositor-keyed shared dust ledger. The per-event identity of Theorem~\ref{thm:binding} closes to the wei on $150{,}357$ of the $152{,}719$ rebalances they made in the windows studied. The remaining discrepancies are consistent with rewards or ledger movements that the event data do not record. The four contracts share a single operator.

Sections~\ref{sec:setup} to~\ref{sec:multipool} derive the ledger's per-event update, the binding-side criterion and the conservation and transfer identities. Section~\ref{sec:chaos} studies its long-run behaviour in a hub-spoke model, in which each pool pairs a common token with a distinct spoke token. Under the stated assumptions, the stationary balances of any fixed number of spokes become independent at rate $1/N$ as the number of pools $N$ grows, and each spoke balance is zero with a probability given in closed form. Section~\ref{sec:model_tests} tests the model's predictions through simulations of both the model and the deployed contract, the latter using exact contract arithmetic. Section~\ref{sec:empirical} checks the exact identities event by event against the chain record and uses replays to test which model predictions persist there. Section~\ref{sec:foundry} reports contract-level tests on two exchanges and machine-checked proofs.

The rule determining which token remains carries over to the contract and, on rebalances that change the range, to the chain. Independence does not carry over where a common price move aligns the pools' rebalance directions. On chain, for some pairs of ledger entries whose tokens never share a pool, whether one entry is zero is positively associated with whether the other is. Replays trace the excess of positive dependence to the shared balance at same-range rebalances, which the model does not describe. Section~\ref{sec:discussion} summarises the evidence and its implications.

\begin{figure}[!htb]
\centering
\definecolor{nodeteal}{HTML}{24528A}
\definecolor{nodetealfill}{HTML}{EEF2F7}
\definecolor{nodecrimson}{HTML}{9B2335}
\definecolor{nodecrimsonfill}{HTML}{F9EFEE}
\definecolor{nodegrey}{HTML}{5A5A5A}
\resizebox{\ifdim\width>\linewidth\linewidth\else\width\fi}{!}{%
\begin{tikzpicture}[
    >=latex,
    font=\footnotesize,
    pool/.style  = {rectangle, draw=nodeteal, line width=0.5pt, fill=nodetealfill,
                    minimum width=2.7cm, minimum height=1.25cm, align=center, inner sep=3pt, rounded corners=2pt},
    cell/.style  = {rectangle, draw=nodecrimson!80, line width=0.5pt, fill=nodecrimsonfill,
                    minimum width=1.1cm, minimum height=1.25cm, align=center, inner sep=2pt},
    arrow/.style = {->, draw=black!60, shorten >=3pt, shorten <=3pt, line width=0.5pt},
]
\node[pool] (p1) at (0, 0) {\acr{WETH}/\acr{AERO} \\[-1pt] {\scriptsize position rebalances}};
\node[cell] (c1) at (4.95, 0) {\acr{WETH} \\[-1pt] 0.190};
\node[cell] (c2) at (6.05, 0) {\acr{USDC} \\[-1pt] 0};
\node[cell] (c3) at (7.15, 0) {\acr{AERO} \\[-1pt] 0};
\node[pool] (p2) at (12.1, 0) {\acr{WETH}/\acr{USDC} \\[-1pt] {\scriptsize position rebalances}};
\node[anchor=south, font=\scriptsize] at (6.05, 0.72) {shared ledger, one entry for each token};
\node[anchor=north, font=\scriptsize, nodegrey] at (6.05, -0.74) {the entry stands until a pool holding the token rebalances};
\draw[arrow] (p1.east) -- node[above, font=\scriptsize] {first, leaves} node[below, font=\scriptsize] {0.190 \acr{WETH}} (c1.west);
\draw[arrow] (c3.east) -- node[above, font=\scriptsize] {next, its mint draws} node[below, font=\scriptsize] {the 0.190 \acr{WETH}} (p2.west);
\end{tikzpicture}}
\caption{A leftover passing between two pools through the shared ledger, using the amounts from the example in Section~\ref{sec:perevent}.}
\label{fig:mechanism}
\end{figure}
\section{Related work}\label{sec:related}

This paper builds on the geometric residual, single-pool convergence and zero-swap extinction results of \emph{The Geometric Siphon}\cite{ryan2026siphon} (Theorems~1 to~3) to study a shared ledger spanning multiple pools. Section~7.1 of that paper identifies the architectural requirement. Leftovers must be accessible to other positions through a shared balance. The Group~B portfolio reported in its \S5.4 provides a small example of cross-pool sharing. Both tokens of one pool also occur in other pools.

\emph{Convexity and Topology in Cross-Pool Dust Redistribution}\cite{pi2026convexity} gives an independent graph-theoretic treatment of cross-pool dust circulation. Its Proposition~1 gives a condition for circulation, and Proposition~2 relates a token's residence time to the number of pools that contain it. Connecting two token nodes whenever they share a pool gives the portfolio graph used in that work. Connectivity alone does not guarantee netting. The positions' token requirements must also offset, which depends on relative token returns, as the shared-token corollary in the prior paper shows. The residence-time identity is consistent with the update frequencies in Section~\ref{sec:chaos_statement}. When all pools rebalance at the same rate, a token shared by more pools has its ledger entry updated more often.

Cartea \emph{et al.}\cite{cartea2024} analyse liquidity provision using \emph{predictable loss} (\acr{PL}), which combines the convexity-related loss of an \acr{LP} position with the opportunity cost of locked capital. Milionis \emph{et al.}\cite{milionis2022b} introduce \emph{loss-versus-rebalancing} (\acr{LVR}) for general \acr{CFMM}s, defined as the return gap between an \acr{AMM} \acr{LP} and a continuously rebalancing benchmark under arbitrage against stale prices. These costs can accumulate while the manager leaves a position unchanged. Heimbach \emph{et al.}\cite{heimbach2022} analyse impermanent loss and the variation in V3 liquidity-provider returns across strategies. The present paper concerns redistribution through the shared ledger when the manager rebalances a position. The swap-free geometric residual vanishes when a position is withdrawn and reminted into the same range (Theorem~1 of the prior paper), even though \acr{LVR} may have accumulated while the position was open.

The dust ledger is constant between rebalance events and changes by discrete jumps. Section~\ref{sec:chaos} studies its stochastic dynamics using methods for interacting particle systems, surveyed by Sznitman\cite{sznitman1991}. Graham and M\'el\'eard\cite{graham1993} prove propagation of chaos for a fully connected loss network with alternate routing. Their analysis relates dependence between components to overlap in their interaction histories, a connection used in Appendix~\ref{app:chaosproof}. Chaintron and Diez\cite{chaintron2022} survey mean-field jump processes and interaction-graph methods. In the absence of swaps, fees and external transfers, each token's total across positions and the ledger is conserved by Theorem~\ref{thm:extinction}, analogous to conservation of customer numbers in a closed queueing network\cite{kelly1979}. To the author's knowledge, this jump-process treatment of shared dust ledgers has not previously appeared in the DeFi literature.

Here, \emph{Master Equation} names the ledger's per-event update rule. Once a stochastic law for the events and their inputs is specified, this rule determines a forward equation for the ledger's distribution. Sznitman\cite{sznitman1991} uses \emph{master equation} for the forward Markov equation of an interacting particle system. Mean-field game theory\cite{carmona2018,cardaliaguet2019} uses the term for an equation on the space of probability measures, a different object.

The V3 amount equations in Appendix~\ref{app:v3formulas} are taken from Adams \emph{et al.}\cite{adams2021v3}, equations~6.29--6.30.

\section{The system}\label{sec:setup}

The notation for an individual rebalance follows the prior paper\cite{ryan2026siphon}. The shared ledger is represented by a vector of token balances that changes over time. Appendix~\ref{app:notation} collects the notation. Token values in \acr{USD} are calculated at fixed reference prices $(p_1, \ldots, p_n)$. \acr{LP} fee receipts, gas payments and externally induced rebalance \acr{MEV} are assumed not to change the dust ledger directly.

\paragraph{Portfolio.} A portfolio consists of $m$ concentrated liquidity positions across $n$ tokens, held by one depositor through one Position Manager contract. For notation, take one position per pool. If several positions share a pool, position-level sums count them separately, using the same token pair and pool price. The portfolio defines a bipartite graph $G$ with token nodes $\{T_1, \ldots, T_n\}$ and pool nodes $\{P_1, \ldots, P_m\}$. Each pool node connects to its two token nodes and therefore has degree 2. The degree $\deg(T_i)$ counts the distinct pools containing token $T_i$, not the positions. For pool $P_k$, write its token pair as $(T_a, T_b)$, with $T_a$ denoting token0 and $T_b$ token1 in V3 ordering.

\paragraph{Dust ledger.} The dust ledger is the vector $D_{\symup{pool}}(t) \in \mathbb{R}^n_{\geq 0}$, whose entries give the amount held of each portfolio token at time $t$. It represents what the prior paper\cite{ryan2026siphon} calls the dust pool. The ledger is constant between rebalance events. At each rebalance, the event inputs determine the new balances of the pool's two tokens. All other entries remain unchanged. For a pool $P_k$ with tokens $(T_a, T_b)$, write the two shared balances available to it as
\[
(x_{\symup{dust}}, y_{\symup{dust}}) = (D_{\symup{pool},T_a}(t), D_{\symup{pool},T_b}(t)).
\]

\paragraph{Rebalance events.} Write $\tau_1, \tau_2, \ldots$ for the rebalance times. A rebalance at time $\tau$ is described by the pool $P_k$, the old range $[s_a, s_b]$ and the new range $[s_a', s_b']$. Its other inputs are the sqrt price $s(\tau)$, the pre-event position value $V(\tau^-)$ and the signed swap amounts $(\sigma_x, \sigma_y) \in \mathbb{R}^2$ defined in Section~\ref{sec:swap}. When the swap moves the pool price, the withdrawal and mint use different prices, as described in that section. The per-event identities take these inputs as given and make no assumption about when rebalances occur. In the stochastic model of Section~\ref{sec:chaos}, rebalance timing is specified externally rather than determined by the ledger balances.

\section{The Master Equation}\label{sec:masterequation}

Theorem~\ref{thm:binding} gives the exact update of the ledger at one rebalance. It follows from the V3 amount equations and requires no stochastic assumptions.

\subsection{Per-event jump}\label{sec:perevent}

The mint chooses the greatest liquidity that the available amounts of both tokens can support. This is a one-dimensional linear programme, whose closed-form solution gives the ledger update.

Consider a pool $P_k$ with $T_a$ as token0 and $T_b$ as token1. At sqrt price $s$, its position is rebalanced from $[s_a, s_b]$ into $[s_a', s_b']$, with $s$ strictly inside the new range. With liquidity $L$, the old position returns the amounts
\begin{equation}\label{eq:oldamount}
x_w = L\,a_x(s, s_a, s_b), \qquad y_w = L\,a_y(s, s_a, s_b),
\end{equation}
where $a_x, a_y$ are the V3 amount functions (Appendix~\ref{app:v3formulas}). With a standing balance $(x_{\symup{dust}}, y_{\symup{dust}})$ of the pool's two tokens and a signed swap correction $(\sigma_x, \sigma_y)$, the mint is offered the inputs
\begin{equation}\label{eq:avail}
\hat x = x_w + x_{\symup{dust}} - \sigma_x, \qquad \hat y = y_w + y_{\symup{dust}} - \sigma_y.
\end{equation}
Write $g_x = a_x(s, s_a', s_b')$ and $g_y = a_y(s, s_a', s_b')$ for the token amounts required per unit of liquidity in the new range.

\needspace{4\baselineskip}
\begin{theorem}[Master Equation jump]\label{thm:binding}
For a rebalance of the pool $P_k$ at sqrt price $s$ into a new range $[s_a', s_b']$ with $s$ in its interior, and mint inputs $(\hat x, \hat y)$ as in~\eqref{eq:avail}, both non-negative, the new liquidity is
\begin{equation}\label{eq:lnew}
L_{\symup{new}}  =  \min\biggl(\frac{\hat x}{g_x},\; \frac{\hat y}{g_y}\biggr),
\end{equation}
and the new dust-ledger entries for $T_a, T_b$ are
\begin{equation}\label{eq:newdust}
x_{\symup{dust}}'  =  \hat x - L_{\symup{new}}\,g_x, \qquad y_{\symup{dust}}'  =  \hat y - L_{\symup{new}}\,g_y,
\end{equation}
both non-negative. At least one of them is zero, and exactly one is zero unless $\hat x\,g_y = \hat y\,g_x$.
\end{theorem}

\noindent Each token amount supports a maximum liquidity. The smaller of the two limits the mint, and the token imposing that limit is the binding token. Any unused amount of the other token becomes its new ledger entry. The prior paper calls the binding token the missing token and the other the surplus token. Any non-zero entry is the \emph{leftover} of the rebalance. The inputs for which nothing remains form the ratio-preserving locus $\{\hat x / \hat y = g_x / g_y\}$, of measure zero in input-ratio space.

\begin{proof}[Proof sketch]
The mint maximises $L_{\symup{new}}$ subject to the two budget constraints $L_{\symup{new}}\,a_x(s, s_a', s_b') \leq \hat x$ and $L_{\symup{new}}\,a_y(s, s_a', s_b') \leq \hat y$. Since the objective increases with $L_{\symup{new}}$, the optimum is the smaller of the two bounds, giving~\eqref{eq:lnew}. The leftovers in~\eqref{eq:newdust} are the unused amounts under the two constraints. The full derivation, including the swap-corrected case, is given in Appendix~\ref{app:masterproof}.
\end{proof}

Each rebalance withdraws the old position's liquidity and combines the withdrawn amounts with the shared balances of its two tokens. After any swap correction (Section~\ref{sec:swap}), the mint uses the available amounts and credits the leftover to the shared ledger. Entries for all other tokens remain unchanged.
\begin{equation}\label{eq:master}
D_{\symup{pool},i}(\tau^+)  =  \begin{cases}
x_{\symup{dust}}' & \text{if } T_i = T_a, \\[2pt]
y_{\symup{dust}}' & \text{if } T_i = T_b, \\[2pt]
D_{\symup{pool},i}(\tau^-) & \text{otherwise.}
\end{cases}
\end{equation}
All terms in \eqref{eq:master} are observable per event from on-chain data and contract state. With an empty standing balance and no swap, Theorem~\ref{thm:binding} reduces to the geometric-residual result in Theorem~1 of the prior paper\cite{ryan2026siphon}. The leftover has value $|\Delta R|$ in that paper's notation and is held in the surplus token. The ranges that leave nothing are its ratio-preserving ranges. Repeating the swap-free update and discarding the balance after each event gives the setting of its zero-swap extinction theorem.

\paragraph{Example.}\label{par:example}
Consider a \acr{WETH}/\acr{USDC} position with \acr{WETH} as token0 and a range centred on 2{,}000 \acr{USDC} per \acr{WETH}. Its lower and upper bounds are obtained by dividing and multiplying the centre price by $1.05$. The price rises to the upper bound of 2{,}100, where the position holds only \acr{USDC}. Withdrawing its liquidity returns 10{,}000 \acr{USDC}. The new range is centred on 2{,}100 with the same factor, and at its centre it takes the two tokens in equal value.

With no swap the mint is offered no \acr{WETH}. It mints nothing, and the whole 10{,}000 \acr{USDC} is the leftover.

Assume the swap executes at 2{,}100 \acr{USDC} per \acr{WETH}, with no fee or price impact. Swapping 5{,}000 \acr{USDC} for 2.381 \acr{WETH} supplies the exact ratio required by the mint, leaving nothing unused. Swap error is measured relative to this amount.

Suppose that the swap is $4\%$ too large, 5{,}200 \acr{USDC} for 2.476 \acr{WETH}. \acr{USDC} is now the binding token, the mint deploys 9{,}600 \acr{USDC} of value, and 0.190 \acr{WETH}, worth 400 \acr{USDC}, is left. The ledger entry of \acr{WETH} is then 0.190 and that of \acr{USDC} is zero.

If the swap is $4\%$ too small, 4{,}800 \acr{USDC} for 2.286 \acr{WETH}, then \acr{WETH} is the binding token, 400 \acr{USDC} is left, and the entry of \acr{WETH} is zero. In both cases, the leftover is worth $4\%$ of the withdrawn value. An oversized swap leaves the token bought; an undersized swap leaves the token sold. Figure~\ref{fig:example_bars} compares these two cases with the one that follows, a rebalance that draws on a standing balance.

The shared ledger enters through the standing balance. Suppose a \acr{WETH}/\acr{AERO} position in the same portfolio left 0.190 \acr{WETH} in the ledger at its last rebalance. The swap that was $4\%$ too small now leaves nothing. The mint draws the 0.190 \acr{WETH}, deploys 10{,}400 \acr{USDC} of value, and both entries are zero. The position has minted 400 \acr{USDC} more than it withdrew, out of what another position left behind. One position has therefore supplied part of another's mint through the shared ledger. Section~\ref{sec:extinction} formalises this transfer.

\begin{figure}[!htb]
\centering
\definecolor{nodeteal}{HTML}{24528A}
\definecolor{nodetealfill}{HTML}{EEF2F7}
\definecolor{nodecrimson}{HTML}{9B2335}
\definecolor{nodecrimsonfill}{HTML}{F9EFEE}
\definecolor{nodegrey}{HTML}{5A5A5A}
\resizebox{\ifdim\width>\linewidth\linewidth\else\width\fi}{!}{%
\begin{tikzpicture}[font=\footnotesize, >=latex]
\node[anchor=south west, inner sep=0pt] at (2.10, 1.00) {withdrawn};
\node[anchor=south west, inner sep=0pt] at (6.30, 1.00) {after the swap};
\node[anchor=south west, inner sep=0pt] at (10.50, 1.00) {minted, and left over};
\node[anchor=east, font=\scriptsize] at (14.05, 1.62) {\acr{WETH}};
\filldraw[fill=nodetealfill, draw=nodeteal, line width=0.5pt] (12.90, 1.53) rectangle ++(0.25, 0.18);
\node[anchor=east, font=\scriptsize] at (12.60, 1.62) {\acr{USDC}};
\filldraw[fill=nodecrimsonfill, draw=nodecrimson, line width=0.5pt] (11.50, 1.53) rectangle ++(0.25, 0.18);
\node[anchor=east, align=right] at (1.90, 0.28) {swap 4\% \\ too large};
\filldraw[fill=nodecrimsonfill, draw=nodecrimson, line width=0.5pt] (2.100, 0.000) rectangle (5.400, 0.550);
\node[font=\scriptsize] at (3.750, 0.275) {10{,}000};
\draw[->, nodegrey, line width=0.5pt] (5.56, 0.28) -- (6.20, 0.28);
\filldraw[fill=nodecrimsonfill, draw=nodecrimson, line width=0.5pt] (6.300, 0.000) rectangle (7.884, 0.550);
\node[font=\scriptsize] at (7.092, 0.275) {4{,}800};
\filldraw[fill=nodetealfill, draw=nodeteal, line width=0.5pt] (7.884, 0.000) rectangle (9.600, 0.550);
\node[font=\scriptsize] at (8.742, 0.275) {5{,}200};
\draw[densely dashed, black, line width=0.5pt] (7.950, -0.080) -- (7.950, 0.630);
\node[anchor=north west, inner sep=0pt, font=\scriptsize] at (6.30, -0.14) {sells 5{,}200 \acr{USDC} for 2.476 \acr{WETH}};
\draw[->, nodegrey, line width=0.5pt] (9.70, 0.28) -- (10.40, 0.28);
\filldraw[fill=nodecrimsonfill, draw=nodecrimson, line width=0.5pt] (10.500, 0.000) rectangle (12.084, 0.550);
\node[font=\scriptsize] at (11.292, 0.275) {4{,}800};
\filldraw[fill=nodetealfill, draw=nodeteal, line width=0.5pt] (12.084, 0.000) rectangle (13.668, 0.550);
\node[font=\scriptsize] at (12.876, 0.275) {4{,}800};
\filldraw[fill=nodeteal, draw=nodeteal, line width=0.5pt] (13.788, 0.000) rectangle (13.920, 0.550);
\node[anchor=north east, align=right, font=\scriptsize\itshape] at (13.99, -0.08) {leftover 0.190 \acr{WETH} \\[-1pt] the token bought};
\node[anchor=east, align=right] at (1.90, -1.48) {swap 4\% \\ too small};
\filldraw[fill=nodecrimsonfill, draw=nodecrimson, line width=0.5pt] (2.100, -1.750) rectangle (5.400, -1.200);
\node[font=\scriptsize] at (3.750, -1.475) {10{,}000};
\draw[->, nodegrey, line width=0.5pt] (5.56, -1.48) -- (6.20, -1.48);
\filldraw[fill=nodecrimsonfill, draw=nodecrimson, line width=0.5pt] (6.300, -1.750) rectangle (8.016, -1.200);
\node[font=\scriptsize] at (7.158, -1.475) {5{,}200};
\filldraw[fill=nodetealfill, draw=nodeteal, line width=0.5pt] (8.016, -1.750) rectangle (9.600, -1.200);
\node[font=\scriptsize] at (8.808, -1.475) {4{,}800};
\draw[densely dashed, black, line width=0.5pt] (7.950, -1.830) -- (7.950, -1.120);
\node[anchor=north west, inner sep=0pt, font=\scriptsize] at (6.30, -1.89) {sells 4{,}800 \acr{USDC} for 2.286 \acr{WETH}};
\draw[->, nodegrey, line width=0.5pt] (9.70, -1.48) -- (10.40, -1.48);
\filldraw[fill=nodecrimsonfill, draw=nodecrimson, line width=0.5pt] (10.500, -1.750) rectangle (12.084, -1.200);
\node[font=\scriptsize] at (11.292, -1.475) {4{,}800};
\filldraw[fill=nodetealfill, draw=nodeteal, line width=0.5pt] (12.084, -1.750) rectangle (13.668, -1.200);
\node[font=\scriptsize] at (12.876, -1.475) {4{,}800};
\filldraw[fill=nodecrimson, draw=nodecrimson, line width=0.5pt] (13.788, -1.750) rectangle (13.920, -1.200);
\node[anchor=north east, align=right, font=\scriptsize\itshape] at (13.99, -1.83) {leftover 400 \acr{USDC} \\[-1pt] the token sold};
\node[anchor=east, align=right] at (1.90, -3.23) {too small, with \\ 0.190 \acr{WETH} \\ in the ledger};
\filldraw[fill=nodecrimsonfill, draw=nodecrimson, line width=0.5pt] (2.100, -3.500) rectangle (5.400, -2.950);
\node[font=\scriptsize] at (3.750, -3.225) {10{,}000};
\filldraw[fill=nodeteal, draw=nodeteal, line width=0.5pt] (5.460, -3.500) rectangle (5.592, -2.950);
\draw[->, nodegrey, line width=0.5pt] (5.69, -3.23) -- (6.20, -3.23);
\filldraw[fill=nodecrimsonfill, draw=nodecrimson, line width=0.5pt] (6.300, -3.500) rectangle (8.016, -2.950);
\node[font=\scriptsize] at (7.158, -3.225) {5{,}200};
\filldraw[fill=nodetealfill, draw=nodeteal, line width=0.5pt] (8.016, -3.500) rectangle (9.600, -2.950);
\node[font=\scriptsize] at (8.808, -3.225) {4{,}800};
\filldraw[fill=nodeteal, draw=nodeteal, line width=0.5pt] (9.600, -3.500) rectangle (9.732, -2.950);
\draw[densely dashed, black, line width=0.5pt] (7.950, -3.580) -- (7.950, -2.870);
\node[anchor=north west, inner sep=0pt, font=\scriptsize] at (6.30, -3.64) {the same swap; the mint draws the 0.190 \acr{WETH}};
\draw[->, nodegrey, line width=0.5pt] (9.83, -3.23) -- (10.40, -3.23);
\filldraw[fill=nodecrimsonfill, draw=nodecrimson, line width=0.5pt] (10.500, -3.500) rectangle (12.216, -2.950);
\node[font=\scriptsize] at (11.358, -3.225) {5{,}200};
\filldraw[fill=nodetealfill, draw=nodeteal, line width=0.5pt] (12.216, -3.500) rectangle (13.932, -2.950);
\node[font=\scriptsize] at (13.074, -3.225) {5{,}200};
\node[anchor=north east, align=right, font=\scriptsize\itshape] at (14.25, -3.58) {nothing left \\[-1pt] mints 10{,}400};
\node[anchor=south west, inner sep=1pt, font=\scriptsize] at (7.990, 0.61) {exact swap};
\end{tikzpicture}}
\caption{The rebalance example, shown to scale in value at 2{,}100 \acr{USDC} per \acr{WETH}. Dashed lines mark the exact swap amounts; solid slivers show the leftovers.}
\label{fig:example_bars}
\end{figure}
\subsection{Swap correction}\label{sec:swap}

In its automatic mode, \acr{PM}~V9 swaps tokens within the rebalance transaction when needed, to bring the supplied amounts towards the ratio required by the new range. It trades through the pool being rebalanced, using the contract's own token balance. Across the four contract versions, swaps accompany 96\% of the rebalances that change the range and 59\% of all rebalances. Rebalances without a swap mostly retain the same range (Appendix~\ref{app:methodology}).

The swap is recorded in a Pool.Swap event. The signed amounts $(\sigma_x, \sigma_y)$ are positive for tokens the contract pays and negative for tokens it receives. They enter equation~\eqref{eq:avail} directly, so Theorem~\ref{thm:binding} applies both with and without a swap.

The theorem uses a single price $s$. When the swap moves the pool price, the withdrawal amounts $(x_w, y_w)$ are evaluated at the pre-swap price using~\eqref{eq:oldamount}, while the new range's liquidity requirements $g_x,g_y$ are evaluated at the post-swap price used by the mint (Appendix~\ref{app:methodology}).

The withdrawn amounts $(x_w, y_w)$ include principal and any uncollected trading fees. The gauge-staked positions studied here have no such uncollected fees. The pair $(\sigma_x, \sigma_y)$ gives the net amounts the contract pays away between withdrawal and mint, whether the swap uses the pool itself or an external router.

If the swap carries the price to or beyond a bound of the new range, or the new range is placed with the price outside it, the mint is single-sided and the entire amount of the other token is credited to the ledger (Appendix~\ref{app:masterproof}). If the swap is skipped and the price is inside the new range, the theorem applies with $\sigma_x = \sigma_y = 0$. Integer rounding can leave a small balance of the binding token, bounded by the token amount required for one unit of liquidity.

\section{Multi-pool consequences}\label{sec:multipool}

The Master Equation determines which token a rebalance leaves and how much remains. Across a portfolio, it also gives conservation identities and an accounting of transfers between positions.

\subsection{Binding-side sign rule}\label{sec:connectorrule}

Once the mint inputs are known, Theorem~\ref{thm:binding} determines the token left and its amount. Expressing the inputs as value shares separates the effects of the price's displacement from the old range, the standing balance and the swap. Throughout this subsection values are in units of $T_b$ at the mint price $s$. For mint inputs $(\hat x, \hat y)$ write $\hat W = s^2\hat x + \hat y$ for their value and define
\[
\hat\omega_a = \frac{s^2\hat x}{\hat W}, \qquad \omega'_a = \frac{s^2 g_x}{s^2 g_x + g_y}
\]
as the value shares of $T_a$ supplied to the mint and required by the new range, respectively.

The proposition concerns a single rebalance. Cross-pool sharing enters through the balances available to that mint. The prior paper's propositions on netting and neutral accounting describe how requirements across positions can offset and how value is attributed between them.

\needspace{4\baselineskip}
\begin{proposition}[Binding-side criterion]\label{thm:connector}
In the setting of Theorem~\ref{thm:binding}, let the mint inputs $(\hat x, \hat y)$ be non-negative and not both zero. The token left, and what it is worth, are
\[
\begin{array}{@{}l@{\qquad}l@{\qquad}l@{}}
\text{if } \hat\omega_a < \omega'_a & \text{the leftover is in } T_b & \text{and is worth } \hat W\,\dfrac{\omega'_a - \hat\omega_a}{\omega'_a}, \\[6pt]
\text{if } \hat\omega_a > \omega'_a & \text{the leftover is in } T_a & \text{and is worth } \hat W\,\dfrac{\hat\omega_a - \omega'_a}{1 - \omega'_a}, \\[6pt]
\text{if } \hat\omega_a = \omega'_a & \text{there is no leftover.} &
\end{array}
\]
\end{proposition}

\noindent The token supplied in a smaller value share than the new range requires is used up, leaving the other token. This restates Theorem~\ref{thm:binding} in value shares. With an empty balance and no swap, it recovers the prior paper's geometric residual\cite{ryan2026siphon}.

The two corollaries concern ranges centred in tick space. The old range is centred on $s_0$ and the new range on $s$, so that $s_a s_b = s_0^2$ and $s_a' s_b' = s^2$. Define the following four quantities, with the first two following the prior paper\cite{ryan2026siphon},
\[
\begin{array}{@{}l@{\qquad}l@{}}
\ell = \ln(s/s_0) & \text{the logarithmic displacement,} \\[3pt]
h = \ln(s_b/s_0) & \text{the logarithmic half-width of the old range,} \\[3pt]
V = s^2 x_w + y_w & \text{the withdrawn value,} \\[3pt]
D = s^2 x_{\symup{dust}} + y_{\symup{dust}} & \text{the value of the standing balances,}
\end{array}
\]
and the function
\[
\alpha(\ell) = \frac{\sinh|\ell|}{e^{h} - \cosh\ell} \quad \text{for } |\ell| \leq h, \qquad \alpha(\ell) = 1 \quad \text{for } |\ell| \geq h .
\]
The quantity $\alpha(\ell)$ is the fraction of the withdrawn value left unused when those tokens alone are offered to a new range centred on the current price, without a swap. It rises from zero at the centre of the old range to one at its bounds.

\needspace{4\baselineskip}
\begin{corollary}[Swap-free rule]\label{cor:swapfree}
In the setting of Proposition~\ref{thm:connector}, let the ranges be centred in tick space, with $\ell$, $h$, $V$, $D$ and $\alpha(\ell)$ as above, and let there be no swap.
\begin{enumerate}[leftmargin=*,label=(\roman*)]
\item With no standing balance of $T_a$ or $T_b$, the leftover is in $T_b$ when $s > s_0$, in $T_a$ when $s < s_0$, and absent when $s = s_0$. It is the fraction $\alpha(\ell)$ of the withdrawn value $V$, which is the whole of it when the price is at or beyond a bound of the old range.
\item With a standing balance worth $D$, if $D < \alpha(\ell)\,V$ the leftover is in the token named in (i), whatever the composition of the balance. A balance worth $\alpha(\ell)\,V$ held wholly in the token that the mint of part~(i) uses up removes the leftover, and a larger one moves the leftover to that token.
\end{enumerate}
\end{corollary}

\noindent Without a swap or standing balance, the price move determines which token remains. It is the token representing the larger share of the withdrawn value. A balance supplied by another position can reverse this outcome only if its value exceeds $\alpha(\ell)\,V$.

\needspace{4\baselineskip}
\begin{corollary}[Swap-corrected rule]\label{cor:swapcorrected}
In the setting of Proposition~\ref{thm:connector}, let the ranges be centred in tick space with $s \neq s_0$, with no standing balance, and let the swap exchange, at the price $s$, the value
\[
Q = \tfrac{1}{2}\,\alpha(\ell)\,V\,(1 + \varepsilon), \qquad -1 < \varepsilon \leq 1/\alpha(\ell),
\]
of the token left in Corollary~\ref{cor:swapfree}(i) for the same value of the other, the upper bound being the swap that sells the whole of the former. Then the leftover is worth $\alpha(\ell)\,V\,|\varepsilon|$, and it is in
\[
\begin{cases}
\text{the token named in Corollary~\ref{cor:swapfree}(i)} & \text{if } \varepsilon < 0, \\
\text{the other token, the one the swap bought,} & \text{if } \varepsilon > 0.
\end{cases}
\]
\end{corollary}

\noindent At its centre, the new range requires equal token values. Swapping half the value that would otherwise remain, $\tfrac{1}{2}\,\alpha(\ell)\,V$, therefore eliminates the leftover. The \emph{execution residual} $\varepsilon$ measures the swap's proportional deviation from this exact amount\cite{ryan2026om}. An undersized swap leaves the token named by the swap-free rule; an oversized swap leaves the token bought.

In the example of Section~\ref{sec:perevent}, the price is at the bound of the old range, so $\alpha = 1$ and $V$ is 10{,}000 \acr{USDC}. The errors $\varepsilon = \pm 0.04$ give a leftover worth 400 \acr{USDC} in one token or the other. The example with a standing balance falls under Proposition~\ref{thm:connector}, which allows both a swap and a balance.

Corollary~\ref{cor:swapcorrected} corresponds to the prior paper's swap-size corollary with the price-impact term omitted. Its swap-scale parameter equals $1 + \varepsilon$\cite{ryan2026siphon}. The proofs are in Appendix~\ref{app:connectorproof}.

Most range changes on the deployed manager include a swap, so the swap-free rule alone does not determine which token remains. The contract sizes the swap using the price corresponding to the \emph{floor tick}, obtained by rounding down to a whole tick. Execution uses the pool's exact price, and the sizing rule makes no allowance for the swap's own price impact.

A follow-on paper\cite{ryan2026om} shows that the resulting execution residual is positive on most upward moves and of either sign on downward moves. The leftover of a swap-corrected rebalance is usually the token bought, reversing the swap-free rule. The leftover is in the token bought on $74\%$ of V9's swap-corrected range changes (Section~\ref{sec:framework_predictions}), and on $61\%$ of the V9 swaps and $84\%$ of the V6 to V8 swaps scored in the prior paper\cite{ryan2026siphon}.

\subsection{Conservation and transfers}\label{sec:extinction}

Under a shared ledger the leftover of one pool can be drawn into the mint of another. The prior paper's Theorem~3\cite{ryan2026siphon} considers a position rebalanced without a swap whenever the price exits its range. If its leftover is not returned, the position decays geometrically. Once the leftover is kept in a ledger shared across pools, a rebalance without a swap leaves each token's total over the positions and the ledger unchanged. The identity counts both deployed tokens and tokens held in the ledger. It does not imply that the amount deployed remains unchanged.

For a position $k$ and a token $T_i$ write $a_{k,i}(s)$ for the amount of $T_i$ the position holds when the sqrt price of its pool $P_k$ is $s$. It is $L^{(k)} a_x(s, s_a^{(k)}, s_b^{(k)})$ when $T_i$ is token0 of $P_k$, the same with $a_y$ when it is token1, and zero when $T_i \notin P_k$.

\needspace{4\baselineskip}
\begin{theorem}[Multi-pool mass conservation without a swap]\label{thm:extinction}
Let a portfolio have $m \geq 1$ positions sharing a depositor-keyed dust mapping. For rebalances without a swap, and in the absence of \acr{LP} fees and external token transfers, the per-token total
\begin{equation}\label{eq:masscons}
M_i(t)  :=  \sum_k a_{k,i}\bigl(s_k(\tau)\bigr)  +  D_{\symup{pool},i}(t)
\end{equation}
is unchanged by every rebalance event, each position being evaluated at the sqrt price $s_k(\tau)$ of its own pool at the event, up to the integer rounding of the pool's amounts. The fixed-price \acr{USD} aggregate $V_{\symup{total}}(t) := \sum_i M_i(t)\,p_i$ inherits the same event-level invariance.
\end{theorem}

\noindent Under the theorem's assumptions, a rebalance moves tokens only between the position and the ledger. Each token's total therefore remains unchanged, up to rounding. A swap changes the two token totals by the amounts exchanged. At fixed reference prices, the change in total value is the net value paid away by the swap, with the opposite sign (Appendix~\ref{app:extinctionproof}).

With one position, the identity states that the position and its ledger together hold unchanged token totals across each event. This is the counterpart in tokens of the prior paper's retained-excursion proposition, which is stated in value.

\paragraph{Transfers between positions.}\label{rem:partition}
Theorem~\ref{thm:extinction} tracks token totals without assigning them to positions. A mint can use shared balances regardless of which position supplied them. To attribute value between positions, assign each position $k$ an \emph{entitlement} $G_k \geq 0$ to the ledger, following the prior paper\cite{ryan2026siphon}. The entitlements sum to the ledger's value $\sum_i D_{\symup{pool},i}\,p_i$. With $V_k$ the position's deployed value, its \emph{equity} $U_k = V_k + G_k$ is the sum of its deployed value and its claim on the ledger.

For a rebalance of position $j$, let $\bar\sigma_j = p_a\sigma_x + p_b\sigma_y$ be the net value paid away through the swap, with $p_a$ and $p_b$ the reference prices of the pool's two tokens. The equity available to its mint is then $U_j - \bar\sigma_j$, and the \emph{transfer} at the event is the part of the minted value above that equity,
\[
\Xi_j = \bigl(V_j^{\symup{new}} - U_j + \bar\sigma_j\bigr)^{+}.
\]
If $\Xi_j = 0$ the position has minted within its own equity. Its entitlement becomes $U_j - \bar\sigma_j - V_j^{\symup{new}}$ and all other entitlements remain unchanged.

If $\Xi_j > 0$ the position has used $\Xi_j$ of value owed to other positions. Its entitlement becomes zero and the other entitlements are reduced proportionally, by $\Xi_j$ in total. The mint's budget bounds $\Xi_j$ by their total, since a mint takes no more than the position's withdrawal and the ledger, less what its swap paid away.

Write $\xi_k$ for the net transfer to position $k$ at the event. It is $\Xi_j$ for the rebalancing position and $-\Xi_j\,G_k/\sum_{i \neq j} G_i$ for every other position. All net transfers are zero when $\Xi_j = 0$.

At the fixed reference prices of Section~\ref{sec:setup}, entitlements remain unchanged between events. This is the prior paper's transfer accounting under its ownerless rule, evaluated at fixed prices\cite{ryan2026siphon}. When transaction prices are used instead, entitlements are revalued between events. Section~\ref{sec:framework_predictions} applies proportional revaluation, preserving each position's share of the ledger's value.

A position is a \emph{net receiver} over a window if its cumulative net transfer is positive and a \emph{net giver} if it is negative. These correspond to the prior paper's net absorber and net donor. In the example of Section~\ref{sec:perevent}, the receiving position drew 0.190 \acr{WETH} from the ledger. It had no entitlement of its own and incurred no swap cost. It therefore received a transfer worth 400 \acr{USDC} from the position that had supplied that balance.

\needspace{4\baselineskip}
\begin{lemma}[Transfers sum to zero]\label{lem:partition}
Let $m \geq 1$ positions share a ledger, in the absence of \acr{LP} fees and external token transfers, and let $\xi_k$ be the net transfers at a rebalance of position $j$. Then:
\begin{enumerate}[leftmargin=*,label=(\roman*)]
\item the entitlements still sum to the value of the ledger after the event, and $\sum_k \xi_k = 0$, with or without a swap;
\item over any window the cumulative net transfers of the positions sum to zero, so they are either all zero or include at least one strictly positive and at least one strictly negative value;
\item if $m = 1$ the transfer is zero at every event;
\item the value change of the rebalancing position is $\Delta V_j = \xi_j - \bar\sigma_j - (G_j^{\symup{new}} - G_j)$.
\end{enumerate}
\end{lemma}

\noindent Transfers redistribute value between positions without changing their combined value. A transfer need not equal the change in a position's deployed value at its own rebalance, $\Delta V_k = L_{\symup{new}}^{(k)}\,v_k^{\symup{new}}(s) - L^{(k)}\,v_k^{\symup{old}}(s)$, where $v_k^{\symup{old}}$ and $v_k^{\symup{new}}$ are the position values per unit of liquidity over the old and new ranges. By part~(iv), the change in deployed value equals the transfer received, minus the swap cost and the change in the position's entitlement. The proof is in Appendix~\ref{app:extinctionproof}, and Section~\ref{sec:framework_predictions} presents the V9 evidence.

\section{A model of the ledger and propagation of chaos}\label{sec:chaos}

Sections~\ref{sec:masterequation} and~\ref{sec:multipool} give exact event-level identities under their stated conditions. This section studies the ledger's long-run distribution in a model with one hub token shared by $N$ pools. Each pool pairs the hub with a distinct spoke token. At each rebalance, the model draws a direction and an execution residual at random.

After each rebalance of its own pool, a spoke balance is zero with a probability given in closed form (Proposition~\ref{thm:chaos2B}). In the stationary distribution, any fixed number of spoke balances become independent as $N$ grows, with a discrepancy of order $1/N$ (Proposition~\ref{thm:chaos2A}). This asymptotic independence is called propagation of chaos\cite{sznitman1991}. Both results extend to unequal positions under a condition on their sizes (Proposition~\ref{cor:hetero}). With a common environment, the same asymptotic independence holds after conditioning on the environment, but may fail without that conditioning (Proposition~\ref{prop:commondir}).

The exact identities of Sections~\ref{sec:masterequation} and~\ref{sec:multipool} do not depend on this model. Sections~\ref{sec:model_tests} and~\ref{sec:model_chain} examine how far the model describes the deployed contract. Readers mainly interested in the simulations and empirical results can read the assumptions and the comparison with the deployed contract, then continue to Section~\ref{sec:model_tests}. Figure~\ref{fig:map} shows how the statements depend on one another.

\paragraph{Terms.} The \emph{law} of a quantity is its probability distribution. A \emph{stationary law} is a distribution of balances that remains unchanged under the model's updates, even though individual balances continue to change. The \emph{kernel} is the rule that takes the balances before an event to their law after it. A \emph{mark} is what is drawn at an event, here a direction and an execution residual.

An \emph{atom} of a law is a value that has positive probability. The law of a spoke balance has an atom at zero, whose size is the probability that the balance is exactly zero. The Wasserstein distance $W_1$ measures how far two distributions are apart. It is the smallest expected distance between balances paired according to those distributions.

The model initially holds each position's liquidity fixed. In the contract, every mint resets liquidity according to the standing balance absorbed and the leftover retained. Section~\ref{sec:possize} allows position size to change.

\begin{figure}[!htb]
\centering
\definecolor{nodeteal}{HTML}{24528A}
\definecolor{nodetealfill}{HTML}{EEF2F7}
\definecolor{nodecrimson}{HTML}{9B2335}
\definecolor{nodecrimsonfill}{HTML}{F9EFEE}
\definecolor{nodegray}{HTML}{5A5A5A}
\definecolor{nodegrayfill}{HTML}{F4F4F4}
\resizebox{\ifdim\width>\linewidth\linewidth\else\width\fi}{!}{%
\begin{tikzpicture}[
    >=latex, font=\footnotesize,
    exact/.style = {rectangle, draw=nodeteal, line width=0.5pt, fill=nodetealfill, text width=2.95cm, minimum height=0.95cm, align=flush center, inner sep=2pt, rounded corners=2pt},
    model/.style = {rectangle, draw=nodecrimson!80, line width=0.5pt, fill=nodecrimsonfill, text width=2.95cm, minimum height=0.95cm, align=flush center, inner sep=2pt, rounded corners=2pt},
    kern/.style  = {rectangle, draw=nodegray, line width=0.5pt, fill=nodegrayfill, text width=2.95cm, minimum height=0.95cm, align=flush center, inner sep=2pt, rounded corners=2pt},
    dep/.style   = {->, draw=black!60, shorten >=2pt, shorten <=2pt, line width=0.55pt},
]
\node[exact] (t1) at (0, 0)      {Theorem~\ref{thm:binding} \\[-1pt] {\scriptsize the jump of the ledger}};
\node[exact] (p1) at (3.75, 0)   {Proposition~\ref{thm:connector} \\[-1pt] {\scriptsize the token left}};
\node[exact] (c12) at (7.5, 0)   {Corollaries~\ref{cor:swapfree},~\ref{cor:swapcorrected} \\[-1pt] {\scriptsize without and with a swap}};
\node[exact] (t2) at (11.25, 0)  {Theorem~\ref{thm:extinction}, Lemma~\ref{lem:partition} \\[-1pt] {\scriptsize conservation, transfers}};
\node[kern]  (k)  at (0, -1.55)   {the kernel~\eqref{eq:Khom} \\[-1pt] {\scriptsize one event of the model}};
\node[model] (l2) at (3.75, -1.55) {Lemma~\ref{lem:stationary} \\[-1pt] {\scriptsize the stationary law}};
\node[model] (p2) at (7.5, -1.55)  {Proposition~\ref{thm:chaos2A} \\[-1pt] {\scriptsize independence at rate $1/N$}};
\node[model] (p45) at (11.25, -1.55) {Propositions~\ref{cor:hetero},~\ref{prop:commondir} \\[-1pt] {\scriptsize unequal spokes, a common direction}};
\node[model] (p3) at (3.75, -3.1) {Proposition~\ref{thm:chaos2B} \\[-1pt] {\scriptsize the atom at zero}};
\node[model] (l3) at (11.25, -3.1) {Lemma~\ref{lem:indicators} \\[-1pt] {\scriptsize zero indicators}};
\draw[dep] (t1) -- (p1);
\draw[dep] (p1) -- (c12);
\draw[dep] (t1.north) -- ++(0,0.3) -| (t2.north);
\draw[dep] (t1) -- (k);
\draw[dep] (k) -- (l2);
\draw[dep] (l2) -- (p2);
\draw[dep] ([xshift=0.95cm]k.north) -- ++(0,0.2) -| ([xshift=-0.7cm]p2.north);
\draw[dep] (p2) -- (p45);
\draw[dep] (l2) -- (p3);
\draw[dep] (k.south) |- (p3.west);
\draw[dep] ([xshift=1.0cm]p3.north) -- ++(0,0.3) -| (p45.south);
\draw[dep, dashed, shorten <=0pt] (0,-3.1) -- (0,-3.875) -| (l3.south);
\end{tikzpicture}}
\caption{Dependencies between the statements. The top row contains the exact identities; the remaining rows concern the stochastic model. The dashed link to Lemma~\ref{lem:indicators} indicates that it uses only the kernel properties specified in its hypotheses.}
\label{fig:map}
\end{figure}
\paragraph{Assumptions.}
Two quantities are drawn at each event of a pool, a \emph{direction} $\eta \in \{+1, -1\}$ and an \emph{execution residual} $\varepsilon$. The residual is that of Corollary~\ref{cor:swapcorrected}, the amount of the swap relative to the exact amount, less one. The sign of the direction depends on whether it is defined for a pool or for a particular ledger entry.

\emph{Pool direction.} In the frame of the pool, $\eta$ is the side of the range at which the event fires. It is $+1$ at the upper bound, where the withdrawn position holds token1 only and the corrective swap buys token0. It is $-1$ at the lower bound, where the swap buys token1.

For an event inside the range, as on chain and in the simulation of the contract, $+1$ denotes a swap that buys token0. If no swap occurs, $\eta$ is the direction of the swap that would correct the mismatch. The model and its statements use the pool direction, with the hub at token0.

\emph{Entry-relative direction.} For comparisons involving an individual ledger entry, as in Sections~\ref{sec:model_contract} and~\ref{sec:model_chain}, direction is positive when the corrective swap sells that entry's token. This agrees with pool direction for token1 and reverses it for token0. The sign of the product $\eta \varepsilon$ is reversed with it.

The model exchanges tokens at the pool price, so specifying the amount sold also determines the amount bought. Values are expressed in common units in which that price is one. To leading order in range width, the leftover equals the absolute execution residual multiplied by the magnitude of the imbalance the swap must remove, $A_\eta$ of~\eqref{eq:imbalance} (Appendix~\ref{app:chaosproof}). The statements retain a positive scale parameter $c_0$ multiplying this product. For the kernel derived here, $c_0 = 1$. With $K$ the bound on the residual in (A1$'$) write
\[
\kappa := c_0 K, \qquad \kappa_1 := c_0\,\mathbb{E}|\varepsilon| .
\]
\begin{enumerate}[leftmargin=*,label=(A\arabic*$'$)]
\item Residual. At each event a direction $\eta$ is drawn with probabilities $(p_+, p_-)$, and then a residual $\varepsilon$ from a law $\nu_\eta$ on $[-K, K]$, independently of the state and of every other event, with $\kappa < 1/2$. The two laws need not coincide, and no symmetry, zero mean or Gaussianity is assumed.
\setcounter{enumi}{1}
\item[(A2)] Boundary trigger. Every range has the same width $w = s_b - s_a$, and the event fires at a bound of the range, so that the withdrawn position holds one token only.
\item[(A3)] Homogeneous exchangeable spokes. The parameters $(L, w)$, $(p_+, p_-)$ and $(\nu_+, \nu_-)$ are shared by the pools. The pool that fires at each event is drawn uniformly from the $N$ pools, independently of the earlier events, of the marks and of the state. The hub token is at the same position in every pool. The liquidity $L$ is held fixed. Proposition~\ref{cor:hetero} treats unequal spokes.
\end{enumerate}
These three assumptions define the model. Table~\ref{tab:constants} collects the main constants of its statements, where $u^+ = \max(u, 0)$ and $u^- = \max(-u, 0)$.

\begin{table}[!htb]
\centering
\caption{The constants of the model.}
\label{tab:constants}
\begin{tabular}{@{}llp{8.2cm}@{}}
\toprule
\thd{Constant} & \thd{Definition} & \thd{Role}  \\
\midrule
$c_0$ & one, as derived & Scale of the leftover, relative to the residual times the imbalance. \\
$\kappa$ & $c_0 K$ & Bound on the leftover, as a share of the imbalance~\eqref{eq:imbalance}. Below one half, which keeps the imbalance positive. \\
$\kappa_1$ & $c_0\,\mathbb{E}|\varepsilon|$ & Mean leftover, as a share of the imbalance. Sets the contraction rate of Lemma~\ref{lem:stationary}. \\
$b^*$ & $\kappa Lw/(1 - \kappa)$ & Bound on every balance. \\
$q$ & $\mathbb{P}(\eta \varepsilon \geq 0)$ & Mass at zero of a spoke's balance (Proposition~\ref{thm:chaos2B}). \\
$\varsigma$ & $c_0^{1/2}\,\mathbb{E}\bigl[|\varepsilon|^{1/2}\bigr]$ & Less than one. Enters the constant of Proposition~\ref{thm:chaos2A}. \\
$\gamma_\pm$ & $c_0\,\mathbb{E}\bigl[(\eta\varepsilon)^\pm\bigr]$ & Mean leftover reaching the hub ($+$) and the spoke ($-$), as shares of the imbalance. \\
\bottomrule
\end{tabular}
\end{table}

\paragraph{The model and the deployed contract.}
The model differs from the deployed manager in five ways.

First, the contract computes the residual from the event inputs rather than sampling it. The within-tick price offset provides the noise-like component (Appendix~\ref{app:chaosproof} gives the residual's first-order form and its size on chain). The residual is $-1$ when the contract skips the swap, which it does when the mismatch is below a threshold or the price lies outside the new range. Its law differs by direction, which is why (A1$'$) carries two laws. It also depends on position size and the standing balance, which (A1$'$) excludes, and its tail can exceed the bound $\kappa < 1/2$.

Second, the manager fires inside the range, so the withdrawn position holds both tokens. The swap must remove the position's token mismatch, which is smaller than its whole value.

Third, liquidity changes at every event (Section~\ref{sec:possize}).

Fourth, a move in the price of a shared token can align rebalance directions across pools holding that token. Proposition~\ref{prop:commondir} allows this dependence through a common environment.

Fifth, about two fifths of the manager's events are \emph{same-range rebalances}, which withdraw a position and mint it again into the same range. The contract skips the swap on most of them. The prior paper finds that the ledger grows on range changes and shrinks on same-range rebalances\cite{ryan2026siphon}. The model has no such event.

Table~\ref{tab:residual_chain} in Appendix~\ref{app:chaosproof} gives the residual measured on chain, and Sections~\ref{sec:model_tests} and~\ref{sec:model_chain} set the model against the simulated contract and the chain record.

\subsection{Hub-spoke topology and statement of the limit}\label{sec:chaos_statement}

Consider the hub-spoke portfolio, with one hub token $T_0$ shared by $N$ pools and each pool pairing it with its own spoke token $T_k$ (Fig.~\ref{fig:hubspoke}). Write $H := D_{\symup{pool},T_0}$ and $B_k := D_{\symup{pool},T_k}$ for the ledger entries of the hub and of spoke $k$, throughout this section and Appendix~\ref{app:chaosproof}.

Every event updates the hub balance, while only an event of pool $k$ updates $B_k$. With equal firing rates, the hub is therefore updated $N$ times as often as any individual spoke. This difference in update frequency motivates the asymptotic independence proved below.

The update frequencies are consistent with Proposition~2 of the note on convexity and topology\cite{pi2026convexity}, in which residence time is inversely proportional to degree. The Group~B portfolio in \S5.4 of the prior paper\cite{ryan2026siphon} contains a three-spoke instance of this structure, formed by its three pools holding \acr{WETH}. A fourth pool is attached to one of the spoke tokens.

\begin{figure}[!htb]
\centering
\definecolor{nodeteal}{HTML}{24528A}
\definecolor{nodetealfill}{HTML}{EEF2F7}
\definecolor{nodecrimson}{HTML}{9B2335}
\definecolor{nodecrimsonfill}{HTML}{F9EFEE}
\definecolor{nodegray}{HTML}{5A5A5A}
\definecolor{nodegrayfill}{HTML}{F4F4F4}
\begin{tikzpicture}[
    >=latex,
    font=\small,
    hub/.style       = {circle, draw=nodeteal, line width=0.5pt,
                        fill=nodetealfill, minimum size=1.7cm, align=center, inner sep=1pt},
    spoke/.style     = {circle, draw=nodecrimson!80, line width=0.5pt,
                        fill=nodecrimsonfill, minimum size=1.5cm, align=center, inner sep=1pt},
    pool/.style      = {rectangle, draw=nodegray, line width=0.5pt,
                        fill=nodegrayfill, minimum width=0.95cm, minimum height=0.65cm, align=center,
                        font=\footnotesize, inner sep=2pt},
    edge/.style      = {draw=black!55, line width=0.5pt},
    fired/.style     = {draw=black, line width=1.4pt},
]
\node[hub] (hub) at (0,0) {$T_0$ \\[-1pt] {\footnotesize entry $H$}};
\node[pool] (p1) at (-3.2, 1.6) {$P_1$};
\node[pool] (p2) at (-1.6, 2.4) {$P_2$};
\node[pool] (p3) at ( 1.6, 2.4) {$P_3$};
\node[pool] (pN) at ( 3.2, 1.6) {$P_N$};
\node[spoke] (s1) at (-4.4, 3.2) {$T_1$ \\[-1pt] {\footnotesize $B_1$}};
\node[spoke] (s2) at (-2.0, 4.0) {$T_2$ \\[-1pt] {\footnotesize $B_2$}};
\node[spoke] (s3) at ( 2.0, 4.0) {$T_3$ \\[-1pt] {\footnotesize $B_3$}};
\node[spoke] (sN) at ( 4.4, 3.2) {$T_N$ \\[-1pt] {\footnotesize $B_N$}};
\node[font=\Large] at (0, 2.8) {$\cdots$};
\foreach \p/\s in {p1/s1, p2/s2, pN/sN} {
    \draw[edge] (hub) -- (\p);
    \draw[edge] (\p) -- (\s);
}
\draw[fired] (hub) -- (p3);
\draw[fired] (p3) -- (s3);
\node[anchor=west, font=\footnotesize, align=left] at (2.9, 0.15) {an event of $P_3$ rewrites $H$ and $B_3$ \\ and leaves one of them at zero};
\end{tikzpicture}
\caption{The hub-spoke portfolio. A rebalance of pool $k$ updates the hub balance $H$ and its spoke balance $B_k$. All other entries remain unchanged.}
\label{fig:hubspoke}
\end{figure}

Let the hub token be token0 of every pool. By (A2), at an event of pool $P_k$ in direction $\eta = +1$ the withdrawn position holds only the spoke token and the corrective swap buys hub token, and in direction $\eta = -1$ it holds only hub token and the swap buys spoke token. Amounts are valued in units in which the pool price is one, and to leading order in $w$ a position withdrawn at either bound is then worth $Lw$. Write
\begin{equation}\label{eq:imbalance}
A_\eta  :=  Lw + \eta\,(B_k - H)
\end{equation}
for the imbalance the swap is asked to remove, so that a residual $\varepsilon > 0$ is a swap that bought too much. By Theorem~\ref{thm:binding} the event then rewrites the two entries, to leading order in $w$, as
\begin{equation}\label{eq:Khom}
(H',\, B_k')  = 
\begin{cases}
\bigl(c_0\,|A_\eta|\,|\varepsilon|,\ 0\bigr) & \text{if } \eta A_\eta \varepsilon > 0, \text{ the leftover going to the hub,} \\[3pt]
\bigl(0,\ c_0\,|A_\eta|\,|\varepsilon|\bigr) & \text{otherwise, the leftover going to the spoke,}
\end{cases}
\end{equation}
where $c_0 = 1$ is the constant introduced before (A1$'$), and every other entry is unchanged. This rule is the kernel of the model, written $\widetilde{\mathcal{K}}_w$ in Appendix~\ref{app:chaosproof}.

Table~\ref{tab:kernel_sides} shows which entry is zero after each combination of direction and residual sign, when the imbalance is positive. Again, an oversized swap leaves the token bought and an undersized swap leaves the token sold. The spoke's entry is zero exactly when $\eta \varepsilon \geq 0$. Both entries are zero when $\varepsilon = 0$.

\begin{table}[!htb]
\centering
\caption{The entry an event sets to zero, with the hub at token0.}
\label{tab:kernel_sides}
\begin{tabular}{@{}llllll@{}}
\toprule
\thd{Direction} & \thd{Fires at} & \thd{Swap buys} & \thd{Residual} & \thd{Zero entry} & \thd{Leftover on}  \\
\midrule
$\eta = +1$ & upper bound & hub token & $\varepsilon > 0$ & spoke & hub \\
$\eta = +1$ & upper bound & hub token & $\varepsilon < 0$ & hub & spoke \\
$\eta = -1$ & lower bound & spoke token & $\varepsilon > 0$ & hub & spoke \\
$\eta = -1$ & lower bound & spoke token & $\varepsilon < 0$ & spoke & hub \\
\bottomrule
\end{tabular}
\end{table}

In the example of Section~\ref{sec:perevent}, the hub is \acr{WETH} and the spoke is \acr{USDC}. The event fires at the upper bound and the swap buys \acr{WETH}, so the first two rows apply, with $\varepsilon = +0.04$ and $-0.04$. Valuing both tokens in \acr{USDC} makes the pool price one in the kernel's units. With $c_0 = 1$ and an empty ledger, the leftover is worth $0.04 \times 10{,}000$ \acr{USDC} in each case.

In the entry-relative frame, an entry is zero when $\eta \varepsilon \geq 0$. The swap sells \acr{USDC}, so the \acr{USDC} entry is zero when the swap is oversized. It buys \acr{WETH}, so the \acr{WETH} entry is zero when the swap is undersized. If the hub is token1 of the pool, the same kernel holds with $\eta$ replaced by $-\eta$.

Under (A1$'$), with
\[
b^*  :=  \frac{\kappa Lw}{1 - \kappa}  <  Lw ,
\]
the set on which $H \leq b^*$ and every $B_j \leq b^*$ is invariant, meaning that balances which start in it stay in it, and $A_\eta \geq Lw - b^* > 0$ on it (Appendix~\ref{app:chaosproof}). The analysis below restricts the balances to this invariant set. Write $X^N = (H, B_1, \ldots, B_N)$ for the balances, and for a law $\pi^N$ of $X^N$ write $\pi^N_{(m)}$ for the law of $m$ of the spoke balances. A family of laws is tight when none of its mass escapes to infinity.

\needspace{4\baselineskip}
\begin{lemma}[Stationary law]\label{lem:stationary}
Under (A1$'$)--(A3), the chain $X^N$ admits a unique stationary law $\pi^N$ on $[0, b^*]^{N+1}$. From any initial law on that set the law of $X^N$ converges to $\pi^N$ in $W_1$, the distance contracting by the factor $1 - (1 - \kappa_1)/N$ at each event, and the family $\{\pi^N_{(m)}\}_{N \geq m}$ is tight on $\mathbb{R}_{\geq 0}^m$ for every fixed $m$.
\end{lemma}

\noindent From any initial distribution within the invariant set, the distribution of balances converges to the same stationary law. When two copies use the same sequence of events and marks, their expected distance contracts at each event by at least the factor stated in the lemma. Two statements follow, one on the joint law of the spokes and one on the law of a single spoke. In both, $\mu$ and $\pi_H$ are the limiting laws of a spoke and the hub, defined by the fixed-point system in Appendix~\ref{app:chaosproof}.

\needspace{4\baselineskip}
\begin{proposition}[Propagation of chaos at rate $1/N$]\label{thm:chaos2A}
Assume (A1$'$)--(A3). Let $\pi^N$ be the unique stationary law of $X^N = (H, B_1, \ldots, B_N)$ on $[0, b^*]^{N+1}$ and $(\pi_H, \mu)$ the unique solution of the fixed-point system~\eqref{eq:McKVfixed}. For $1 \leq m \leq N$ let $\pi^N_{(m)}$ be the law of $m$ spoke balances under $\pi^N$ and $\pi^N_{(0,m)}$ the joint law of the hub balance and $m$ spoke balances. Then
\[
W_1\bigl(\pi^N_{(m)},\,\mu^{\otimes m}\bigr)  \leq  W_1\bigl(\pi^N_{(0,m)},\,\pi_H \otimes \mu^{\otimes m}\bigr)  \leq  \frac{m\,(m + 5)}{2 (1 - \varsigma)^4}\cdot\frac{b^*}{N} ,
\]
where $W_1$ is the Wasserstein distance for the sum of absolute differences and $\varsigma < 1$ is the constant of Table~\ref{tab:constants}. In particular $\pi^N_{(m)} \to \mu^{\otimes m}$ weakly for every fixed $m$.
\end{proposition}

\noindent In stationarity, any fixed collection of spokes approaches independent copies of the limiting spoke balance as the number of pools grows. The bound quantifies the remaining dependence and the difference from the limiting marginal law.

\paragraph{Remark (moments, and when the rate is attained).}
The means $\bar h$ and $\bar b$ of the hub and spoke balances under $\pi^N$ are those of $\pi_H$ and $\mu$ at every $N$, and
\[
\operatorname{Cov}_{\pi^N}(H, B_k)  =  -\,\frac{\bar h\,\bar b}{N}\cdot\frac{1 - \chi_-}{1 - \chi_- + (1 - 1/N)\,\chi_+} ,
\]
where $\chi_\pm := c_0\,\mathbb{E}\bigl[(\eta\varepsilon)^\pm\,\eta\bigr]$, with $u^+ = \max(u, 0)$ and $u^- = \max(-u, 0)$. When each of the two tokens receives a leftover with positive probability, $\bar h\,\bar b > 0$. The discrepancy is then of order $1/N$, rather than merely bounded above at that order, both for the hub-spoke pair and for the law of a single spoke. When leftovers fall on one token only, $\pi^N = \pi_H \otimes \mu^{\otimes N}$ exactly. These are proved with the proposition in Appendix~\ref{app:chaosproof}.

\needspace{4\baselineskip}
\begin{proposition}[Spoke marginal, exact atom and slack-side approximation]\label{thm:chaos2B}
Under (A1$'$)--(A3), for every $N \geq 1$ the stationary law $\mu^N = \pi^N_{(1)}$ of a spoke balance has the binary atom-mixture form
\begin{equation}\label{eq:chaoslimit}
\mu^N(db)  =  q \delta_0(db)  +  (1 - q)\,\rho^N_+(db),
\end{equation}
with atom
\[
q  :=  \mathbb{P}(\eta \varepsilon \geq 0)  =  p_+ \nu_+[0, K] + p_- \nu_-[-K, 0]
\]
and $\rho^N_+$ a law on $(0, b^*]$. The atom is exact and does not depend on $N$. The slack-side law satisfies
\[
W_1\bigl(\rho^N_+,\,\rho_+\bigr)  \leq  c_0 b^*\;\mathbb{E}\bigl[\,|\varepsilon| \mid \eta \varepsilon < 0\,\bigr],
\]
where $\rho_+$ is the law of $c_0 Lw\,|\varepsilon|$ conditional on $\eta \varepsilon < 0$. The limit $\mu$ of Proposition~\ref{thm:chaos2A} has the same form, with the same $q$ and the same bound.
\end{proposition}

\noindent After an event of its own pool, a spoke balance is zero with probability $q$. In stationarity, conditional on being positive, it has the slack-side distribution described in the proposition.

The probability of zero is exact for every $N$. After an event of its own pool, the spoke balance is zero exactly when $\eta \varepsilon \geq 0$, and it does not change until that pool's next event. The slack-side distribution is approximated by the law of $c_0 Lw\,|\varepsilon|$ conditional on $\eta \varepsilon < 0$, with the stated error bound. The approximation is controlled by the ratio $b^*/Lw$ of balance to position size. Neither conclusion requires the asymptotic-independence result.

With one residual law for both directions, no probability mass at zero and equally likely directions, $q = 1/2$ regardless of the residual law. With equally likely directions and no probability mass at zero, a value of $q$ different from one half therefore indicates that the residual's sign depends on direction.

\paragraph{Remark (the Gaussian case).}
Let $\nu_+ = \nu_-$ be the mean-zero Gaussian law of standard deviation $\sigma_{\symup{slip}}$, conditioned to $[-K, K]$, and let $\sigma_{\symup{hn}} = c_0 Lw\,\sigma_{\symup{slip}}$. A law symmetric about zero with no atom there gives $q = 1/2$ exactly.

Under the conditioned Gaussian law $|\varepsilon|$ is half-normal with scale $\sigma_{\symup{slip}}$ conditioned to $[0, K]$, so $\rho_+$ is the half-normal law $\mathrm{HN}(\sigma_{\symup{hn}})$ conditioned to $[0, c_0 LwK]$. Up to a relative error of order $e^{-K^2/(2\sigma_{\symup{slip}}^2)}$ from the conditioning, the leading-order law of a spoke balance is $\tfrac{1}{2}\delta_0 + \tfrac{1}{2}\,\mathrm{HN}(\sigma_{\symup{hn}})$ and the mean of $\rho_+$ is $\sigma_{\symup{hn}}\sqrt{2/\pi}$. 

The stationary slack-side mean differs from the mean of $\rho_+$ by at most the bound of Proposition~\ref{thm:chaos2B}.

Proposition~\ref{thm:chaos2A} concerns the joint law, and the Remark that follows it gives its rate. Both propositions and the Gaussian case are simulated in Section~\ref{sec:chaos_numerical} and set against the contract in Section~\ref{sec:model_contract}, and the chain evidence is in Section~\ref{sec:model_chain}.

The hub has a different limiting distribution from the spokes. Its balance is updated at every event of every pool. Its probability of a zero balance is $\mathbb{P}(\eta \varepsilon \leq 0)$, since each event sets at least one of the two balances to zero. Its law $\pi_H$ is one factor of the product limit in Proposition~\ref{thm:chaos2A}. The proofs are in Appendix~\ref{app:chaosproof}.
\subsection{Unequal spokes and a common direction}\label{sec:chaos_hetero}

Proposition~\ref{thm:chaos2A} assumes equal spokes and directions that are independent from one event to the next. The next two results allow unequal spokes and a common environment.

Let spoke $k$ have its own size $L_k w_k$, its own firing rate, and its own direction probabilities and residual laws on $[-K_k, K_k]$. The marks are independent of the state and of one another as in (A1$'$). The boundary trigger of (A2) holds with one width $w_k$ for the ranges of pool $k$. The hub is at the token0 position of every pool, one constant $c_0$ serves all pools, and every $L_k$ is held fixed. The imbalance at an event of pool $k$ is $A_\eta = L_k w_k + \eta\,(B_k - H)$, and the kernel is~\eqref{eq:Khom} with it. The statement uses the following, for pool $k$ and for a set $T$ of spokes.
\begin{center}
\begin{tabular}{@{}l@{\qquad}p{0.6\linewidth}@{}}
$\kappa_k := c_0 K_k < 1$ & bounds the leftover of pool $k$, as a share of its imbalance, \\[2pt]
$q_k := \mathbb{P}(\eta \varepsilon \geq 0)$ & which the proposition shows to be the mass at zero of the balance of spoke $k$, \\[2pt]
$\varsigma^*$ & the largest over the pools of $c_0^{1/2}\,\mathbb{E}\bigl[|\varepsilon|^{1/2}\bigr]$, \\[2pt]
$\varpi_k$ & the probability that an event is of pool $k$, its share of the pools' firing rates, positive for every pool, \\[2pt]
$\varpi^* := \max_k \varpi_k$, \ $\varpi_T := \sum_{k \in T} \varpi_k$ & the largest such share, and the share of the pools in $T$.
\end{tabular}
\end{center}
The condition on the sizes is
\begin{equation}\label{eq:heterobox}
b^*  :=  \max_k \frac{\kappa_k\,L_k w_k}{1 - \kappa_k}  <  \min_k L_k w_k ,
\end{equation}
which keeps the largest possible standing balance below the size of every position, ensuring that the imbalance remains positive. This condition also forces every $\kappa_k$ below one half. For laws $\pi_H, \mu_1, \ldots, \mu_N$ on $[0, b^*]$ let $\pi_H^{(k)}$ and $\mu_k'$ be the laws of $H'$ and $B_k'$ after an event of pool $k$ from $(H, B_k) \sim \pi_H \otimes \mu_k$. The fixed-point system is
\begin{equation}\label{eq:McKVhetero}
\mu_k  =  \mu_k' \quad (k = 1, \ldots, N), \qquad \pi_H  =  \sum_{k} \varpi_k\,\pi_H^{(k)} ,
\end{equation}
in which each spoke has its own law and the hub law is a mixture over the pool that fired last.

\needspace{4\baselineskip}
\begin{proposition}[Unequal spokes]\label{cor:hetero}
For the hub-spoke chain with unequal spokes as above, suppose~\eqref{eq:heterobox}. Then the chain has a unique stationary law $\pi^N$ on $[0, b^*]^{N+1}$, the fixed-point system~\eqref{eq:McKVhetero} has a unique solution $(\pi_H, \mu_1, \ldots, \mu_N)$, and for every set $T$ of $m$ spokes the joint law $\pi^N_{(0,T)}$ of the hub balance and the balances of the spokes in $T$ satisfies
\[
W_1\Bigl(\pi^N_{(0,T)},\ \pi_H \otimes \bigotimes_{k \in T} \mu_k\Bigr)  \leq  \frac{m\,\bigl(\varpi^* + 2 \varpi_T\bigr)}{(1 - \varsigma^*)^4}\;b^* .
\]
Each spoke law has the form
\[
\mu_k  =  q_k\,\delta_0 + (1 - q_k)\,\rho_{+,k} ,
\]
where $q_k$ is also the exact mass at zero of the balance of spoke $k$ under $\pi^N$, and $\rho_{+,k}$ is a law on $(0, b^*]$ within $c_0 b^*\,\mathbb{E}\bigl[\,|\varepsilon| \mid \eta \varepsilon < 0\,\bigr]$ in $W_1$ of the law of $c_0 L_k w_k\,|\varepsilon|$ given $\eta \varepsilon < 0$, both taken for pool $k$. The means of all balances under $\pi^N$ are those of $\pi_H$ and the $\mu_k$.
\end{proposition}

\noindent Spokes of different sizes, rates and residual laws still decouple, each with its own atom, under the stated size condition and provided no individual pool retains a non-vanishing share of all events as the number of pools grows. With rates of the same order, $\varpi_k \leq c/N$, the bound is at most $3 c\,m^2\,b^*/\bigl((1 - \varsigma^*)^4 N\bigr)$, and Proposition~\ref{thm:chaos2A} is the case of equal spokes up to the constant. No closed form for the joint stationary law is claimed. Section~7 of the supplement proves the result by adapting the proof of Proposition~\ref{thm:chaos2A}.

The atom of a spoke survives without~\eqref{eq:heterobox} wherever that spoke's own size exceeds the bound,
\begin{equation}\label{eq:heterocond}
L_k w_k  >  b^* ,
\end{equation}
whatever the sizes of the others (Section~7 of the supplement).

\paragraph{A common direction.}
A shared market condition can affect several pools at once. The model represents this by an environment state that changes the direction probabilities and residual distributions.

Take equal spokes under (A2) and (A3), with $b^* = \kappa Lw/(1 - \kappa)$ as in Section~\ref{sec:chaos_statement}. Let the direction probabilities and residual laws of (A1$'$) depend on a state $e$ common to the pools, with $c_0|\varepsilon| \leq \kappa < 1/2$ in every state. The state may fix the direction itself, or shift the direction probabilities and residual laws.

Let $(e_t)$ be any sequence of states, one for each event. Conditional on the full sequence, events are independent as before, with the mark at time $t$ drawn under $e_t$. For laws $\pi_H$ and $\mu$ on $[0, b^*]$, let $\pi_H'$ and $\mu'$ be the laws of $H'$ and $B_k'$ after an event from $(H, B_k) \sim \pi_H \otimes \mu$ under state $e_t$. The reference laws are defined by the recursion
\begin{equation}\label{eq:MFflow}
\pi_H^t  =  \bigl(\pi_H^{t-1}\bigr)' , \qquad \mu^t  =  \Bigl(1 - \frac{1}{N}\Bigr)\,\mu^{t-1} + \frac{1}{N}\,\bigl(\mu^{t-1}\bigr)' ,
\end{equation}
the primes being taken for the pair $(\pi_H^{t-1}, \mu^{t-1})$. The hub updates at every event. A particular spoke updates with probability $1/N$, so its distribution is a mixture of its unchanged and updated laws. The recursion has exactly one solution defined for all $t$, which for a constant sequence is the $\pi_H$ and $\mu$ of Proposition~\ref{thm:chaos2A}. It contains $N$, so unless the sequence is constant the reference law depends on $N$, and the proposition is a comparison at each $N$.

\needspace{4\baselineskip}
\begin{proposition}[A common direction]\label{prop:commondir}
For a given sequence $(e_t)$ let $\rho_t$ be the law of $X^N$ at time $t$ for the chain run from the infinite past, and let $(\pi_H^t, \mu^t)$ be the solution of the recursion~\eqref{eq:MFflow}. Then for every sequence, every $t$ and $1 \leq m \leq N$,
\[
W_1\bigl(\rho_{t,(0,m)},\ \pi_H^t \otimes (\mu^t)^{\otimes m}\bigr)  \leq  \frac{m\,(5m + 13)}{(1 - \varsigma^*)^5}\cdot\frac{b^*}{N} ,
\]
where $\rho_{t,(0,m)}$ is the joint law of the hub balance and $m$ spoke balances and $\varsigma^*$ is the supremum over the states of $c_0^{1/2}\,\mathbb{E}\bigl[|\varepsilon|^{1/2}\bigr]$. The means of the balances under $\rho_t$ are those of $\pi_H^t$ and $\mu^t$. If the sequence is random, with any law, and does not depend on the events, the same holds conditionally on it, and for two spokes $i \neq j$
\[
\Bigl|\,\operatorname{Cov}\bigl(B_i, B_j\bigr) - \operatorname{Var}\bigl(\bar b_t\bigr)\Bigr|  \leq  \frac{46}{(1 - \varsigma^*)^5}\cdot\frac{(b^*)^2}{N} , \qquad \bar b_t := \int b\,\mu^t(db) ,
\]
with $\bar b_t$ a function of the sequence.
\end{proposition}

\noindent Conditional on the full environment sequence, the spokes decouple at the rate of Proposition~\ref{thm:chaos2A}. Without conditioning, their covariance is the variance of the conditional mean $\bar b_t$, up to a term of order $1/N$. The size of that variance depends on the environment.

When the environment is independent from event to event, the variance is at most of order $1/N$.

When the environment is drawn once and never changes, the variance does not depend on $N$. The spokes then do not decouple if the mean spoke balance differs between environment states. For environments between these two cases, the size of the variance is not determined here. The proof and the reasons for the two cases are in Section~7 of the supplement.

\paragraph{Zero indicators.}
Closeness in $W_1$ does not guarantee similar probabilities of an exactly zero balance. A balance that is always slightly positive can be close to zero in distance while never equalling zero. Zero indicators therefore require a separate argument.

An entry remains unchanged until a pool containing its token rebalances. Its zero indicator therefore records the outcome of the most recent event involving that token. The following identity applies on any portfolio graph under the stated assumptions.

Number the events by $t$, with pool $k_t$ firing in environment state $e_t$. For a token $a$, write $t_a \leq t$ for the time of the last event involving it, assumed to exist. Write $I_a(t)$ for the indicator that its ledger entry is zero after event $t$. This is the outcome recorded at event $t_a$.

\needspace{4\baselineskip}
\begin{lemma}[Zero indicators]\label{lem:indicators}
Suppose that, given the environment and the order in which the pools fire, the outcomes of different events are independent, an event of pool $k$ in state $e$ leaving the entry of its token $a$ at zero with probability $q^a(k, e)$. Let $a$ and $b$ be tokens with no pool in common, and write $Q_a(t) := q^a(k_{t_a}, e_{t_a})$. Then $I_a(t)$ and $I_b(t)$ are independent given the environment and the order of events, and
\[
\operatorname{Cov}\bigl(I_a(t),\,I_b(t)\bigr)  =  \operatorname{Cov}\bigl(Q_a(t),\,Q_b(t)\bigr) .
\]
If moreover the environment is stationary, the firing pools of different events are independent of one another and of the environment, and $a$ and $b$ are spokes whose pools each fire with probability $1/N$ and which have the same probability $q_t$ of a zero entry in the state of time $t$, a stationary sequence with autocovariance $C$, then
\[
\operatorname{Cov}\bigl(I_a(t),\,I_b(t)\bigr)  =  \sum_{j \geq 1} \frac{1}{N}\Bigl(1 - \frac{1}{N}\Bigr)^{j-1} C(j) .
\]
\end{lemma}

\noindent Under the lemma's conditional-independence assumptions, the dependence between two zero indicators equals the dependence between the zero probabilities at their most recent events. The kernel~\eqref{eq:Khom} satisfies the hypothesis wherever the imbalance is positive at every event, which holds under (A1$'$)--(A3) and under~\eqref{eq:heterobox}, provided neither the firing nor the marks react to the ledger. The side of the leftover is then a function of the event's own mark, and $q^a$ is the probability that $\eta \varepsilon \geq 0$ for a spoke.
With independent directions the indicators of two spokes are therefore independent at every $N$. For a regime that is a stationary Markov chain on two states, with $q_t$ of correlation $\varrho$ between consecutive events, $C(j) = \operatorname{Var}(q)\,\varrho^{\,j}$ and the series is
\begin{equation}\label{eq:indicatorregime}
\operatorname{Cov}\bigl(I_a(t),\,I_b(t)\bigr)  =  \operatorname{Var}(q)\;\frac{\varrho}{N\,(1 - \varrho) + \varrho} ,
\end{equation}
which tends to $\operatorname{Var}(q)\,\tau_{\mathrm{r}}/(\tau_{\mathrm{r}} + 2)$ when each state is left with probability $1/(N\tau_{\mathrm{r}})$ at each event, so that $\varrho = 1 - 2/(N\tau_{\mathrm{r}})$. A regime that persists for a number of events of the order of $N$ therefore leaves a covariance that does not decay with $N$, whenever the probability of a zero entry differs between its two states. The proof of the lemma is in Section~7 of the supplement.

Recorded features such as pool, direction and position-size category can be used to group events. The frequency of zero balances within each group estimates its zero probability. Under the lemma's assumptions, these probabilities determine the dependence expected between two indicators through the groups of their most recent events.

Provided event timing does not react to the ledger, additional dependence may reflect an unrecorded common influence or an outcome that depends on the standing balance. The latter makes the outcome depend on earlier events and breaks the lemma's hypothesis. The observed discrepancy alone cannot distinguish these explanations.

The second set of replays in Section~\ref{sec:model_chain} isolates the contribution of shared balances to the excess of positive dependence by removing the balance passed between pools. The contribution of common prices at the mint is bounded there, but not excluded.

\subsection{Position size}\label{sec:possize}

The chain of Section~\ref{sec:chaos_statement} holds $L$ fixed. In the contract, each mint determines the position's new liquidity. Take the kernel~\eqref{eq:Khom} with $c_0 = 1$ and all values in common units. At a rebalance of pool $P_k$ in direction $\eta$ the imbalance is $A_\eta = L_k w + \eta\,(B_k - H)$, with the pool's own size in place of the common one of~\eqref{eq:imbalance}. The mint inputs sum to $L_k w + B_k + H$ and differ by $|A_\eta|\,|\varepsilon|$, so the new liquidity satisfies
\begin{equation}\label{eq:Lupdate}
L_k'\,w = L_k w + B_k + H - |A_\eta|\,|\varepsilon| ,
\end{equation}
and the leftover $|A_\eta|\,|\varepsilon|$ goes to the hub or to the spoke as in~\eqref{eq:Khom}. Value is then conserved at every event. The change in deployed value $L_k w$ is not zero on average. Its mean given the state and the direction is $H + B_k - |A_\eta|\,\mathbb{E}\bigl[\,|\varepsilon| \mid \eta\,\bigr]$, which is negative at an empty balance and positive when the absorbed balance exceeds the expected leftover.

Call $Z_k = L_k w + B_k$ the \emph{own value} of position $k$, its deployed value together with the balance of its spoke token, which no other pool can draw on. It leaves out the hub balance, which every pool can draw on. This own value is distinct from the entitlement-based equity defined in Section~\ref{sec:extinction}.

When each pool fires with probability $1/N$ and the imbalances of the two positions are positive in both directions, the difference between the own values of two positions shrinks in expectation by the factor $1 - \gamma_+/N$ at each event, with an additional term depending on the difference between their spoke balances, with $\gamma_+$ of Table~\ref{tab:constants} the mean leftover that reaches the hub. 

If leftovers never reach the hub, every leftover returns to its own position, and once the hub balance is empty the own value of each position is constant. 

This is the counterpart, for pools that share one token, of the single-pool convergence result of the prior paper\cite{ryan2026siphon}. There the donated fraction is an assumption, and here it follows from the kernel. The statement and its proof are in the supplement (Proposition~S1).

\section{The model in simulation}\label{sec:model_tests}

The simulations first check the model's quantitative statements and measure how conservative its bounds are (Section~\ref{sec:chaos_numerical}). They then test which statements persist under the deployed contract's arithmetic and rebalance policy (Section~\ref{sec:model_contract}).

\subsection{The model chain}\label{sec:chaos_numerical}

The simulation implements the chain defined in Section~\ref{sec:chaos_statement}, with $c_0 = 1$ and units chosen so that $Lw = 1$. At each event, it chooses a pool uniformly, draws a direction and a residual independently of the past, and applies the kernel~\eqref{eq:Khom}. The runs cover $N = 2, 4, \ldots, 1024$. For each combination of residual law and $N$, it runs for between $2 \times 10^9$ and $2 \times 10^{10}$ events after a burn-in set by the contraction bound of Lemma~\ref{lem:stationary}. The floating-point results agree with an exact rational implementation of the kernel to within $2 \times 10^{-15}$. Both implementations take the same branch at every event compared.

Four residual laws satisfy (A1$'$). Under the \emph{Gaussian} law, the residual in both directions is Gaussian with standard deviation $0.05$, conditioned to $[-0.45, 0.45]$. Under the \emph{uniform} law, it is uniform on $[-0.45, 0.45]$ in both directions. Under the \emph{skewed} law, $\nu_+$ is uniform on $[-0.02, 0.30]$ and $\nu_-$ is uniform on $[-0.25, 0.25]$. Under these three laws, the two directions are equally likely.

The \emph{measured} law takes $\nu_+$ and $\nu_-$ from the execution residuals of the $72{,}028$ V9 rebalances with a measured swap in the fully covered window. The residuals are grouped by the token the swap buys, and the direction probabilities are those of the data. The residuals are clipped at $|\varepsilon| = 0.45$, which affects $4.35\%$ of events.

A fifth law, the \emph{unclipped} law, uses the measured residuals without clipping. The swaps the contract skipped are added to $\nu_-$ as an atom at $\varepsilon = -1$. This law lies outside (A1$'$) and shows what happens when the bounded-residual assumption is violated.

\emph{The simulated zero share agrees with the exact atom.} The \emph{zero share} of a spoke is the fraction of its events after which its balance is zero. Under the four laws that satisfy (A1$'$), it equals $q$ to within Monte Carlo error at every $N$. In every simulated event under these four laws, the spoke balance is zero exactly when $\eta\varepsilon \geq 0$ (Table~\ref{tab:chaosverif}). Under the measured law, the atom is $q = 0.762$, while the share of events with a non-negative residual is $0.713$. The atom therefore depends on the direction as well as on the sign of the residual. Under the unclipped law, the imbalance changes sign on up to $1.3\%$ of events, and the simulated zero share falls short of $q$ by approximately the same amount.

\emph{The slack side is close to its leading form.} The simulated slack-side law is within a few per cent of $\rho_+$ in $W_1$, relative to the mean of $\rho_+$. On the same scale, the bound of Proposition~\ref{thm:chaos2B} is $43$ to $82\%$. Under the Gaussian law, the slack-side mean equals $\sigma_{\symup{hn}}\sqrt{2/\pi}$ of the Remark after Proposition~\ref{thm:chaos2B} to six digits. The variance exceeds the half-normal's by $0.5\%$, and the quantiles are within $0.3\%$ of the half-normal's (supplement, Table~S1).

\emph{The spokes decouple at order $1/N$.} Under the skewed and measured laws, the simulated value of $N$ times the correlation between two spoke balances is constant from $N = 2$ to $1024$ (Fig.~\ref{fig:decoupling}; supplement, Table~S2). The simulated scaled covariance $N\,\operatorname{Cov}(B_1, B_2)$ agrees with the exact value derived in the proof of Proposition~\ref{thm:chaos2A} within Monte Carlo error at every $N$. The mean balances equal $\bar h$ and $\bar b$ to within $6 \times 10^{-6}$.

Under the Gaussian and uniform laws, the same mean-zero residual law is used for both directions. The balances of two spokes are then uncorrelated at every $N$, as the exact covariance requires. Their dependence appears in the squared balances, again at order $1/N$ (supplement, Section~1 and Table~S2). Under all four laws, the zero indicators of two spokes are uncorrelated at every $N$, as Lemma~\ref{lem:indicators} requires.

\emph{Convergence is at least as fast as the contraction bound of Lemma~\ref{lem:stationary}.} Two copies started from different balances and driven by the same events converge towards one another at least at the rate given by Lemma~\ref{lem:stationary}. The simulated rate is between one and 1.4 times the rate predicted by the bound. The Gaussian and uniform laws have the same $q = 1/2$ but different values of $1 - \kappa_1$, and their rates differ accordingly.

Table~\ref{tab:chaosverif} summarises each law across all values of $N$. Its first two columns report the atom $q$ and the largest deviation of the simulated zero share from $q$. The next two report the largest $W_1$ distance between the slack-side law and $\rho_+$, and the bound of Proposition~\ref{thm:chaos2B}; both are divided by the mean of $\rho_+$. The final two columns compare the contraction bound $1 - \kappa_1$ with the observed contraction rate multiplied by $N$, for $N \geq 16$.

\begin{table}[!htb]
\centering
\caption{The simulated chain against Proposition~\ref{thm:chaos2B} and Lemma~\ref{lem:stationary}, by residual law.}
\label{tab:chaosverif}
\begin{tabular}{@{}lrrrrrr@{}}
\toprule
\thd{Law} & $q$ & \thd{largest} $|\hat q - q|$ & $W_1$\thd{/mean} & \thd{bound} & $1 - \kappa_1$ & $N \times$ \thd{rate}  \\
\midrule
Gaussian & 0.5000 & $2.8 \times 10^{-5}$ & 0.0015 & 0.818 & 0.960 & 0.98--1.08 \\
Uniform & 0.5000 & $4.5 \times 10^{-6}$ & 0.0673 & 0.818 & 0.775 & 0.87--0.95 \\
Skewed & 0.7188 & $5.3 \times 10^{-6}$ & 0.0650 & 0.429 & 0.867 & 0.92--1.11 \\
Measured & 0.7621 & $8.5 \times 10^{-6}$ & 0.0539 & 0.818 & 0.900 & 0.96--1.12 \\
Unclipped & 0.7716 & $1.1 \times 10^{-2}$ & 0.0572 & -- & 0.861 & 0.96--1.03 \\
\bottomrule
\end{tabular}
\end{table}


\paragraph{Unequal spokes, position size and a common direction.}
Three further runs use the same simulation (supplement, Table~S3).

The first run gives the spokes unequal sizes, $L_k w_k = 10^{U_k}$ with $U_k$ uniform on $[-1/2, 1/2]$. The sizes are drawn once for each $N$, and the run uses the skewed and measured laws. These sizes do not satisfy the sufficient condition~\eqref{eq:heterobox}, so Proposition~\ref{cor:hetero} makes no prediction for this run. The runs suggest that the atom and the rate $1/N$ can persist beyond its condition.

On every spoke that satisfies~\eqref{eq:heterocond}, the atom equals $q$ within Monte Carlo error. At every event with a positive imbalance, the identity behind the atom holds. On the other spokes, the atom falls short of $q$ by close to the share of that spoke's events at which the imbalance is not positive, to within $0.01$ under these two laws. A spoke that fails~\eqref{eq:heterocond} but meets no event with a non-positive imbalance keeps the exact atom. Condition~\eqref{eq:heterocond} is therefore sufficient but not necessary. The correlation between two standardised spoke balances still falls with $N$, with an exponent of $-0.97$.

The second run lets the mint set the size, as in~\eqref{eq:Lupdate}. Two groups of spokes start at own values $1.6$ and $0.4$. The run uses the four laws and a one-sided law with $\eta\varepsilon \leq 0$ always ($\nu_+$ uniform on $[-0.30, 0]$, $\nu_-$ uniform on $[0, 0.25]$). The difference in mean own value between the two groups contracts at the rate predicted in Section~\ref{sec:possize} (supplement, Proposition~S1(iv)). The first-order rate $\gamma_+/N$ overstates the contraction by up to $14\%$ under the uniform law, whose spoke balances are largest. The rate computed with the simulated balance gap agrees within $1.5\%$ under every law for $N \geq 16$. Under the one-sided law, no leftover reaches the hub and the gap does not change. Total value is constant up to rounding in every chain.

The third run adds a regime $e \in \{+1, -1\}$ common to all pools, which serves as the environment of Proposition~\ref{prop:commondir}. The regime switches with probability $1/(N\tau_{\mathrm{r}})$ at each event, and each event's direction equals $e$ with probability $(1 + r)/2$. The runs take $r \in \{0, 0.3, 0.6, 0.9\}$ and $\tau_{\mathrm{r}} \in \{1, 10\}$ under the skewed law. For $r > 0$, the common regime produces a dependence that does not decay with $N$.

Write $q_\pm := \nu_\pm\{\eta\varepsilon \geq 0\}$. Across the two regime states, the probability of a zero entry has variance $\tfrac14\,(q_+ - q_-)^2\,r^2$. The covariance of the zero indicators of two spokes is given by~\eqref{eq:indicatorregime}, with the limit $\tfrac14\,(q_+ - q_-)^2\,r^2\,\tau_{\mathrm{r}}/(\tau_{\mathrm{r}} + 2)$. The simulation matches this covariance at every $N$, $r$ and $\tau_{\mathrm{r}}$. The regime makes two spokes' zero indicators dependent only when the two states give different zero probabilities, so that $q_+ \neq q_-$.

The simulated correlation between spoke balances becomes positive, reaching $0.006$ to $0.15$ at $N = 1024$, against $-2 \times 10^{-5}$ when directions are independent. Proposition~\ref{prop:commondir} predicts that the covariance equals the variance of the conditional mean of a spoke balance given the regime's path, up to a term of order $1/N$. That variance was not computed separately. Proposition~\ref{thm:chaos2A} proves the $1/N$ rate when directions are independent across events. This simulation shows that the rate need not hold when directions are dependent across events.

\subsection{The model against the contract}\label{sec:model_contract}

The contract comparison uses an exact integer simulation of the deployed manager and pool, extended from the prior paper's shared-balance simulation\cite{ryan2026siphon} to $N = 63$ spokes. The model samples residuals, holds liquidity fixed and draws directions independently across pools. The contract computes residuals from pool state, resets liquidity at every mint and exposes pools to common token-price moves. The simulation reproduces these features. It computes residuals as the contract does, rebalances at an in-range trigger and re-mints at every event. A move in the hub's price affects all pools at once.

Token prices follow random walks. Each pool's price equals its fair price, the other liquidity in each pool is full range, and positions earn no fees. The \emph{liquidity share} of a position is its liquidity as a share of the pool's other liquidity. It starts at $3 \times 10^{-4}$, $3 \times 10^{-3}$, $0.03$ and $0.1$. Each value is run with the hub at token0 and at token1, and with the hub's price either fixed or moving as much as each spoke's. The integer state of the extended simulation agrees bit for bit with a reference implementation. With the extensions switched off, it also agrees with the prior paper's simulation.

\emph{The contract follows the model's side rule.} After a pool's event, its spoke-token balance is zero exactly when $\eta\varepsilon \geq 0$ in the entry-relative frame of the spoke entry. This holds with the hub at either token position and at every $N$ tested. In these simulations, the standing balance never changed which token was left (Table~\ref{tab:contract_side}). Token position enters only through the law of $\varepsilon$ given $\eta$.

At liquidity shares of $3 \times 10^{-4}$ and $3 \times 10^{-3}$, a swap that buys token1 is undersized on nearly every event, and a swap that buys token0 is oversized. The spoke entry is therefore zero after almost every event when the spoke is token1 and after almost none when it is token0. Once the liquidity share exceeds the floor-tick term (Appendix~\ref{app:chaosproof}), every swap is oversized and the direction alone determines which token is left. The zero share is then close to one half. It is not exactly one half, because the swap buys token0 on about $52\%$ of events. The first-order zero shares are computed with a uniform offset and equally likely directions. They differ from the realised zero shares by $0.02$ to $0.025$ in absolute share at every tested liquidity share above $3 \times 10^{-4}$.

\begin{table}[!htb]
\centering
\caption{Zero side and slack side under the contract, $N = 7$, hub's price fixed.}
\label{tab:contract_side}
\resizebox{\ifdim\width>\linewidth\linewidth\else\width\fi}{!}{%
\begin{tabular}{@{}llcccc@{}}
\toprule
 & & \multicolumn{2}{c}{\thd{Spoke entry zero}} & \multicolumn{2}{c}{\thd{Slack side}}  \\
\cmidrule(lr){3-4}\cmidrule(l){5-6}
\thd{Liquidity share} & \thd{Spoke at} & \thd{Zero share} & \thd{First order} & \thd{Median rel.\ error} & \thd{KS distance}  \\
\midrule
$3 \times 10^{-4}$ & token1 & $0.998$ & $0.999$ & $0.010$ & $0.020$ \\
$3 \times 10^{-4}$ & token0 & $0.002$ & $0.001$ & $0.010$ & $0.115$ \\
$3 \times 10^{-3}$ & token1 & $0.952$ & $0.931$ & $0.010$ & $0.108$ \\
$3 \times 10^{-3}$ & token0 & $0.047$ & $0.069$ & $0.010$ & $0.130$ \\
$0.03$ & token1 & $0.525$ & $0.500$ & $0.019$ & $0.265$ \\
$0.03$ & token0 & $0.477$ & $0.500$ & $0.019$ & $0.418$ \\
$0.1$ & token1 & $0.524$ & $0.500$ & $0.062$ & $0.458$ \\
$0.1$ & token0 & $0.475$ & $0.500$ & $0.065$ & $0.435$ \\
\bottomrule
\end{tabular}}
\end{table}

\emph{With the hub price fixed, dependence between spoke balances falls approximately as $1/N$ where it can be detected.} Where the liquidity share is $0.03$ or more, the correlation between two spokes' balances falls with $N$ at a fitted exponent between $-1.11$ and $-0.98$. Each $95\%$ interval lies inside $[-1.3, -0.95]$. In Table~\ref{tab:contract_q1}, $\bar c$ is the mean correlation between two spoke balances, and $N\bar c$ is this correlation multiplied by $N$, at $N = 63$. The model column gives the same quantity for the model chain using the contract's first-order residual law. The correlations have the model's sign but not its magnitude. $N$ times the correlation ranges from about two thirds to about seven times the model chain's value, with the largest ratio at $N = 2$.

At the two smaller liquidity shares, three of the four series show no detectable dependence. The fourth, a liquidity share of $3 \times 10^{-4}$ with the hub at token0, is undecided. As in the model, the zero indicators of two spokes show no dependence in seven of the eight series. The remaining series is the undecided one above.

\begin{table}[!htb]
\centering
\caption{Dependence between two spokes under the contract.}
\label{tab:contract_q1}
\resizebox{\ifdim\width>\linewidth\linewidth\else\width\fi}{!}{%
\begin{tabular}{@{}llcccccc@{}}
\toprule
 & & \multicolumn{4}{c}{\thd{Hub's price fixed}} & \multicolumn{2}{c}{\thd{Hub's price moving,} $N = 63$}  \\
\cmidrule(lr){3-6}\cmidrule(l){7-8}
\thd{Liquidity share} & \thd{Hub} & \thd{Exponent} & $95\%$ \thd{interval} & $N\bar c$ & \thd{Model} & \thd{Corr., balances} & \thd{Corr., zeros}  \\
\midrule
$0.03$ & token0 & $-0.98$ & $[-1.02, -0.95]$ & $-0.062$ & $-0.016$ & $0.190$ & $0.252$ \\
$0.03$ & token1 & $-1.11$ & $[-1.30, -0.95]$ & $-0.006$ & $-0.010$ & $0.225$ & $0.252$ \\
$0.1$ & token0 & $-0.99$ & $[-1.01, -0.97]$ & $-0.110$ & $-0.054$ & $0.217$ & $0.252$ \\
$0.1$ & token1 & $-0.99$ & $[-1.02, -0.97]$ & $-0.083$ & $-0.043$ & $0.213$ & $0.251$ \\
\bottomrule
\end{tabular}}
\end{table}

\emph{With a common price move they do not decouple} (Fig.~\ref{fig:decoupling}). When the hub's price moves, the correlation between two spokes' balances is about $0.2$ and the correlation between their zero indicators is $0.25$. Both hold at every $N$ from $2$ to $63$, wherever the liquidity share is $0.03$ or more. The last two columns of Table~\ref{tab:contract_q1} give the values at $N = 63$.

For a liquidity share of $0.03$ with the hub at token0, the simulation was run with a shared ledger and with a separate balance for each position. The separate-balance portfolios give a correlation of $0.19$ to $0.20$ between balances and the same $0.25$ between indicators. The dependence therefore comes from the price move and not from the ledger. A move in the hub's price gives every pool's swap the same direction, and the direction determines which token is left. Proposition~\ref{prop:commondir} and Lemma~\ref{lem:indicators} describe such a dependence in the model. Their hypotheses were not checked for the simulated contract. The size of the dependence reported here is measured in the contract simulation and is not derived from those statements.

\emph{The leading-order size of the slack side is accurate at liquidity shares up to $0.03$ and less so at $0.1$.} The non-zero balance after an event is compared with the leading-order prediction, $|\varepsilon|$ times the position's own value mismatch, divided by the binding token's required share (as summarised by the median relative error in Table~\ref{tab:contract_side}). At liquidity shares up to $0.03$, with the hub's price fixed or moving, the median absolute relative error is $0.009$ to $0.021$. At most four events in $260{,}000$ differ from the leading-order size by more than $10\%$. At $0.1$, the median is $0.06$ to $0.07$, and $10$ to $24\%$ of events differ by more than $10\%$.

\emph{The half-normal law of the Gaussian case does not describe the contract.} The Kolmogorov--Smirnov (\acr{KS}) distance between the non-zero balance and a half-normal of the same second moment is $0.10$ to $0.49$, except where the spoke is almost never slack (the last column of Table~\ref{tab:contract_side}). Applying the model's leading-order approximation to the same events produces a similar mismatch. The misfit therefore lies in the residual's law and not in the kernel. Under the contract, $\varepsilon$ is not a centred Gaussian.

The transfers of Section~\ref{sec:extinction} sum to zero by construction, and in the simulation they do so at every event up to rounding. A separate test examined how position sizes settle across pools of different depth. It did not meet the criterion declared for it in advance, so no statement on sizes across pools is made here.
\begin{figure}[!htb]
\centering
\definecolor{nodeteal}{HTML}{24528A}
\definecolor{nodecrimson}{HTML}{9B2335}
\definecolor{nodegrey}{HTML}{5A5A5A}
\begin{tikzpicture}
\begin{loglogaxis}[width=10.2cm, height=6.6cm, xmin=1.6, xmax=1500, ymin=3e-6, ymax=0.6,
    xlabel={number of pools $N$}, ylabel={correlation of two spoke balances, magnitude},
    xtick={2,8,32,128,512}, xticklabels={2,8,32,128,512}, tick label style={font=\footnotesize}, label style={font=\footnotesize},
    axis lines=left, tick align=outside, clip=false]
\addplot[nodegrey, densely dashed, line width=0.5pt, domain=2:1024, samples=2] {0.0062/x};
\addplot[nodegrey, densely dashed, line width=0.5pt, domain=2:1024, samples=2] {0.0197/x};
\addplot[nodegrey, densely dashed, line width=0.5pt, domain=2:63, samples=2] {0.059/x};
\addplot[only marks, mark=*, mark size=1.5pt, nodeteal] coordinates {(2,0.0100432) (4,0.00496671) (8,0.00248244) (16,0.00123612) (32,0.000621151) (64,0.000307615) (128,0.000154615) (256,7.7671e-05) (512,3.75936e-05) (1024,1.92787e-05)};
\addplot[only marks, mark=o, mark size=1.5pt, nodeteal] coordinates {(2,0.0031384) (4,0.00157539) (8,0.000780525) (16,0.000385636) (32,0.000192594) (64,9.67277e-05) (128,4.9651e-05) (256,2.38277e-05) (512,1.24219e-05) (1024,5.31323e-06)};
\addplot[only marks, mark=square*, mark size=1.5pt, nodecrimson] coordinates {(2,0.0281133) (3,0.0196981) (5,0.0110429) (7,0.00861001) (15,0.00362892) (31,0.00183425) (63,0.000983385)};
\addplot[mark=square, mark size=1.5pt, nodecrimson, line width=0.5pt] coordinates {(2,0.168398) (3,0.178895) (5,0.183526) (7,0.184825) (15,0.187847) (31,0.189831) (63,0.190474)};
\node[anchor=west, font=\scriptsize, nodecrimson] at (axis cs:75,0.19) {contract, common price move};
\node[anchor=west, font=\scriptsize, nodecrimson] at (axis cs:75,0.0011) {contract, hub price fixed};
\node[anchor=south west, font=\scriptsize, nodeteal] at (axis cs:150,0.00016) {model, skewed law};
\node[anchor=north east, font=\scriptsize, nodeteal] at (axis cs:230,0.000022) {model, measured law};
\node[anchor=south west, font=\scriptsize, nodegrey] at (axis cs:2.2,0.000006) {dashed lines fall as $1/N$};
\end{loglogaxis}
\end{tikzpicture}
\caption{Correlation between two spoke balances against the number of pools, in magnitude. It is negative and falls approximately as $1/N$ except under a common price move. Contract results use liquidity share $0.03$ with the hub at token0.}
\label{fig:decoupling}
\end{figure}
\section{Empirical validation}\label{sec:empirical}

This section tests the framework against four versions of the Position Manager deployed on Aerodrome Slipstream on Base, from \acr{PM}~V6, where the leftover was first observed, to \acr{PM}~V9. Table~\ref{tab:cross_version} gives their rebalance dates. Section~\ref{sec:cross_version} describes the changes between versions; Appendix~\ref{app:methodology} describes the contracts and data sources.

Except for the items identified in Section~\ref{sec:cross_version}, token amounts in Sections~\ref{sec:cross_version} to~\ref{sec:framework_predictions} and the Section~\ref{sec:model_chain} comparisons are reconstructed from chain events. Dollar values in Section~\ref{sec:framework_predictions} use token prices recorded by the operator's engine. Pool events give the withdrawal, swap, mint and fee amounts for each rebalance. The contract's \texttt{DustCredited} and \texttt{DustWithdrawn} events reconstruct the standing balance in chain order. This replay matches all 42 \texttt{DustWithdrawn} events and agrees with archived reads of the contract's \texttt{dustBalance} mapping at the preceding block for 600 sampled token balances.

The V9 analyses use three windows, which differ in what the event data cover.
\begin{center}
\begin{tabular}{@{}>{\raggedright\arraybackslash}p{0.36\linewidth}@{\quad}l@{\quad}>{\raggedright\arraybackslash}p{0.36\linewidth}@{}}
\toprule
\thd{Analysis} & \thd{V}\scn{9}\thd{ window (}\scn{2026}\thd{)} & \thd{Coverage}  \\
\midrule
Theorem~\ref{thm:binding} (Tables~\ref{tab:cross_version} to~\ref{tab:modelift}) and transfers (Table~\ref{tab:donor_absorber}) & 21 Mar to 8 May & every rebalance to block 45{,}746{,}143, the original dataset cutoff \\
Whole life (Table~\ref{tab:cross_version}) and the sign rule (Table~\ref{tab:connector}) & 21 Mar to 14 Oct & the rebalances whose terms are available (Section~\ref{sec:cross_version}), and the range changes whose token flows are on chain \\
The model against the record (Tables~\ref{tab:zero_share} to~\ref{tab:same_range_v9}) & 21 Mar to 14 May & the events of every pool \\
\bottomrule
\end{tabular}
\end{center}

\subsection{Cross-version replication: V6 to V9}\label{sec:cross_version}

For each contract, Table~\ref{tab:cross_version} counts the rebalances it made and those on which the conservation form of Theorem~\ref{thm:binding} holds exactly in integer token units. The conservation form is $x_{\symup{dust}}' = x_w + x_{\symup{dust}} - \sigma_x - x_{\symup{mint}}$, together with its counterpart in $y$. The column of terms available counts the rebalances whose burn, swap and mint amounts are known. What is observed and what is inferred differs between the rows.

\begin{table}[!htb]
\centering
\caption{Conservation form of Theorem~\ref{thm:binding} on every rebalance, V6 to V9.}
\label{tab:cross_version}
\begin{tabular}{@{}llrrr@{}}
\toprule
\thd{Version} & \thd{Dates (}\scn{2026}\thd{)} & \thd{Rebalances} & \thd{Terms available} & \thd{Close to the wei}  \\
\midrule
\acr{PM}~V6 & 3--5 Mar & 738 & 738 & 653 \\
\acr{PM}~V7 & 6--12 Mar & 3{,}636 & 3{,}636 & 3{,}102 \\
\acr{PM}~V8 & 13--20 Mar & 22{,}006 & 22{,}006 & 20{,}263 \\
\acr{PM}~V9 & 21 Mar--8 May & 126{,}339 & 126{,}339 & 126{,}339 \\
\midrule
Total & & 152{,}719 & 152{,}719 & 150{,}357 \\
\midrule
\acr{PM}~V9, whole life & 21 Mar--14 Oct & 280{,}097 & 142{,}814 & 142{,}814 \\
\bottomrule
\end{tabular}
\end{table}

In the V9 window to 8~May, every term is read from chain events, including the fee leg. The fee leg is zero on every event, as the gauge-staked positions imply. Every rebalance in this window closes to the wei on both tokens, whether or not it changed the range and whether or not it swapped (Table~\ref{tab:results}).

For V6 to V8, the burn, swap and mint amounts are observed. The event data contain no Collect events, so the fee leg is inferred from the conservation identity. An event counts as closing only when the inferred fee leg is exactly zero.

Of the 2{,}362 events that do not close, 1{,}040 are in the \acr{WETH}/\acr{AERO} pool. There, V7 and V8 route the gauge's \acr{AERO} reward through the balance the rebalance draws on. On each of these events, the reward the identity implies is non-negative and the other token closes exactly. The other 1{,}322 have the pattern of ledger movements that the event data do not record for those contracts, namely deposits into, withdrawals from, or reassignments of the balance. The available event data cannot distinguish these discrepancies from such unrecorded movements.

After 14 May 2026, the event data hold principal and mint amounts for only 17 of V9's 70 pools, and 137{,}283 rebalances of the whole life lack a term. Of the 142{,}814 that have all their terms, 7{,}276 fall after 14 May and have a fee leg that is implied by the identity and not observed, and 504 have a swap leg taken from the operator's log. Every one of the 142{,}814 closes. The rows of Table~\ref{tab:cross_version} to 8 May are read from chain events alone.

Each new version changed substantial portions of the contract code. The identity holds exactly on every V9 event whose terms are observed. The events of the earlier versions that do not close are consistent with the reward flows and unrecorded ledger movements described above, which were not reconstructed. The four contracts share one operator (Section~\ref{sec:limitations}).

\subsection{V9 by portfolio and kind of rebalance}\label{sec:results}

Table~\ref{tab:results} breaks the V9 row of Table~\ref{tab:cross_version} down by portfolio and kind of rebalance. Portfolios A to J are the ten portfolios of the depositor studied here that rebalanced in that window. The five portfolios of other depositors on the same contract are grouped in one row. Because a pool can belong to several portfolios, the pool counts do not sum to the total.

\begin{table}[!htb]
\centering
\caption{V9 rebalances by portfolio and kind.}
\label{tab:results}
\begin{tabular}{@{}lrrrrrr@{}}
\toprule
 & & \multicolumn{2}{c}{\thd{Range change}} & \multicolumn{2}{c}{\thd{Same range}} &  \\
\cmidrule(lr){3-4}\cmidrule(lr){5-6}
\thd{Portfolio} & \thd{Rebalances} & \thd{Swap} & \thd{No swap} & \thd{Swap} & \thd{No swap} & \thd{Pools}  \\
\midrule
A & 13{,}617 & 12{,}515 & 764 & 65 & 273 & 6 \\
B & 20{,}121 & 13{,}764 & 473 & 1{,}029 & 4{,}855 & 10 \\
C & 13{,}263 & 3{,}862 & 77 & 531 & 8{,}793 & 9 \\
D & 1{,}251 & 807 & 18 & 53 & 373 & 1 \\
E & 6{,}843 & 2{,}241 & 64 & 66 & 4{,}472 & 5 \\
F & 9{,}072 & 4{,}024 & 319 & 339 & 4{,}390 & 6 \\
G & 8{,}771 & 2{,}372 & 48 & 440 & 5{,}911 & 7 \\
H & 3{,}678 & 1{,}626 & 43 & 314 & 1{,}695 & 6 \\
I & 1{,}315 & 680 & 12 & 176 & 447 & 3 \\
J & 1{,}009 & 269 & 5 & 79 & 656 & 1 \\
Other depositors & 47{,}399 & 26{,}620 & 1{,}201 & 1{,}478 & 18{,}100 & 24 \\
\midrule
Total & 126{,}339 & 68{,}780 & 3{,}024 & 4{,}570 & 49{,}965 & 52 \\
\bottomrule
\end{tabular}
\end{table}

Of the $71{,}804$ rebalances that change the range, 96\% carry a swap, compared with 8\% of same-range rebalances. Every one closes (Table~\ref{tab:cross_version}). These results cover three distinct cases: a range change with a swap ($68{,}780$ events), a range change without one ($3{,}024$), and a same-range rebalance ($54{,}535$). The swap-free range-change case adds a standing balance to the prior paper's setting. In a same-range rebalance the swap is usually skipped, and the mint draws on the withdrawal and the standing balance alone. On the 252 rebalances that credit exactly nothing, the mint used the whole of both tokens.

\paragraph{What the leftover depends on.}\label{sec:lift}
When every input is read from chain, the forward and conservation forms of Theorem~\ref{thm:binding} coincide. This agreement does not show which inputs explain the leftover. Table~\ref{tab:modelift} asks how much of the record is needed to predict it, adding inputs to the mint rule in four steps.
\begin{center}
\begin{tabular}{@{}l@{\quad}>{\raggedright\arraybackslash}p{0.82\linewidth}@{}}
(i) & the withdrawn amounts and the new range alone, at the pre-swap price, with no swap and an empty balance \\[1pt]
(ii) & the standing balance added \\[1pt]
(iii) & the swap's realised amounts added, the mint still solved at the pre-swap price \\[1pt]
(iv) & the mint solved at the post-swap price recorded by the swap event
\end{tabular}
\end{center}
The leftover of step~(i) is the prior paper's residual, of which Corollary~\ref{cor:swapfree}(i) is the centred case. Step~(ii) adds the standing balance, as in Corollary~\ref{cor:swapfree}(ii). Each step is compared with the chain result in two ways. First, does it identify the token holding the leftover? The chain outcome is classified by which token has greater value at the pre-swap price. Second, what is the ratio of the predicted value to the actual value? The 252 events with no leftover are excluded.

\begin{table}[!htb]
\centering
\caption{Leftover predicted from part of the chain record.}
\label{tab:modelift}
\begin{tabular}{@{}lrrrrrr@{}}
\toprule
 & \multicolumn{4}{c}{\thd{Right token (\%)}} & \multicolumn{2}{c}{\thd{Swap, range change}}  \\
\cmidrule(lr){2-5}\cmidrule(lr){6-7}
 & \multicolumn{2}{c}{\thd{Range change}} & \multicolumn{2}{c}{\thd{Same range}} & \thd{Median} & \thd{Within}  \\
\thd{Step} & \thd{Swap} & \thd{No swap} & \thd{Swap} & \thd{No swap} & \thd{ratio} & 10\% (\%)  \\
\midrule
(i) Withdrawal & 27.3 & 65.6 & 41.4 & 71.8 & 17.3 & 0.6 \\
(ii) and balance & 26.8 & 94.9 & 32.7 & 97.3 & 17.1 & 0.5 \\
(iii) and swap & 72.0 & 94.9 & 69.7 & 97.3 & 0.74 & 21.2 \\
(iv) and post-swap price & 100.0 & 94.9 & 100.0 & 97.3 & 1.00 & 100.0 \\
\midrule
Events & 68{,}780 & 2{,}983 & 4{,}570 & 49{,}754 & & \\
\bottomrule
\end{tabular}
\end{table}


When the swap itself is computed from the state of the pool before the transaction, as the contract computes it, the mint does forecast the leftover. The prior paper tests that form on about $140{,}100$ V9 events\cite{ryan2026siphon}. The steps here ask a narrower question, namely what the record's own quantities determine. At the post-swap price, step~(iv) names the leftover's token on all but one of the $73{,}350$ swap-corrected rebalances and matches its value to within 10\% on all but three. On swap-free rebalances, its few misses are leftovers at the scale of the pool's integer rounding (Appendix~\ref{app:methodology}).

Even given the swap's realised amounts, the post-swap price is needed to reconstruct the leftover of a swap-corrected rebalance. Without that price, the record says little about it. Step~(iii) has the swap's realised amounts and lacks only the price move that the swap itself caused. Yet it names the token on 72.0\% of swap-corrected range changes, against 73.9\% for a guess that knows only each pool's mix of leftover tokens. Its predicted value has a median of 0.74 of the actual value, with an interquartile range of 0.18 to 1.04 (supplement, Figure~S1, where both are fractions of the withdrawn value). In no portfolio does the share at step~(iii) exceed the pool-mix guess by more than a point (supplement, Table~S4).

The withdrawal alone does worse than the guess, at 27.3\%. On most of these events, the leftover is in the token the swap bought, as Corollary~\ref{cor:swapcorrected} describes for $\varepsilon > 0$. The withdrawal alone also overstates the value by a median factor of 17, since the swap puts nearly all of the geometric residual into the mint. On swap-free rebalances, the missing input is the standing balance. Adding it raises the share of range changes with the right token from 65.6\% to 94.9\%, as Corollary~\ref{cor:swapfree}(ii) describes.

\subsection{The sign rule and transfers on chain}\label{sec:framework_predictions}

Proposition~\ref{thm:connector} and Lemma~\ref{lem:partition} are consequences of Theorem~\ref{thm:binding} that hold event by event. Checking them on the V9 record requires no stochastic assumptions.

\paragraph{The token left (Proposition~\ref{thm:connector} and its corollaries).}
The three statements are checked on every V9 range change whose token flows are on chain, 82{,}728 events from 21~March to 14~October 2026. Of these, 79{,}189 carry an in-pool swap, and 79{,}185 of those leave a token. A further 3{,}460 carry no swap. The remaining 79 swap through an external router only and enter the criterion alone. The last two columns of Table~\ref{tab:connector} give the share of events on which the swap-free rule names the token left and the share on which the leftover is in the token the swap bought.

\emph{The criterion} is Theorem~\ref{thm:binding} and names the token left on all but 80 of the events. Each of the 80 is a rebalance whose pre-swap price could not be reconstructed from its Swap event, or a tie at the scale of the mint's rounding, and one is a swap-free event whose price is inferred from its own mint.

\emph{The swap-free rule} (Corollary~\ref{cor:swapfree}) concerns the 3{,}249 swap-free range changes that left a token. It names the token on $57\%$ of them, $78\%$ after a downward displacement and $43\%$ after an upward one. The disagreements have two sources, each a case in which a condition of the corollary does not hold.

The first source is the placement of the new range. Below the bound of part~(ii), the standing balance is worth less than $\alpha(\ell)\,V$. There, the criterion with a tick-centred new range agrees with the rule on every event, as part~(ii) requires. Each of the 1{,}063 disagreements below the bound is therefore decided by where the manager placed the new range. Among these disagreements, the new range's centre is a median 2.3 tick spacings from the price, and within half a spacing on one in ten. They are therefore placement and not the rounding of a centred range.

The second source is the standing balance. Above the bound, the balance reverses the rule on 259 of its 336 disagreements, alone or together with the placement.

These events are also a selected set. Almost all swap-free range changes (3{,}316 of 3{,}460) are automatic rebalances on which the contract skipped the swap. They are therefore selected on the fit between the new range and the withdrawn position, and what that fit leaves decides the side.

\emph{The swap-corrected rule} (Corollary~\ref{cor:swapcorrected}) concerns the 79{,}185 swap-corrected range changes that left a token. The leftover is the token the swap bought on $98.0\%$ of those that bought token0 and on $48.3\%$ of those that bought token1. On all but $0.02\%$ of the events whose residual is measured, it is the bought token exactly when $\varepsilon > 0$. The swap-free rule names the token on $27.5\%$ of the 79{,}185. It does so on fewer than half the events in 78 of the 81 portfolio-pool groups with 30 events or more. The rank correlation between displacement from the old centre and signed leftover is negative in 79 of the 81 groups (median $-0.61$).
\begin{table}[!htb]
\centering
\caption{The swap-free and swap-corrected rules on V9 range changes.}
\label{tab:connector}
\begin{tabular}{@{}lrcc@{}}
\toprule
\thd{Range changes} & \thd{Events} & \thd{Swap-free rule agrees} & \thd{In token bought}  \\
\midrule
Swap-free, all                         & 3{,}249  & 0.57 & n/a \\
\quad balance below $\alpha(\ell)V$    & 2{,}560  & 0.58 & n/a \\
\quad balance at or above $\alpha(\ell)V$ & 689   & 0.51 & n/a \\
Swap-corrected, swap bought token0     & 40{,}238 & 0.04 & 0.98 \\
Swap-corrected, swap bought token1     & 38{,}947 & 0.52 & 0.48 \\
Swap-corrected, all                    & 79{,}185 & 0.27 & 0.74 \\
\bottomrule
\end{tabular}
\end{table}
\paragraph{Transfers between positions (Lemma~\ref{lem:partition}).}
The transfer defined in Section~\ref{sec:extinction} is computed by carrying an entitlement for every position through the replayed V9 ledger, as in the prior paper\cite{ryan2026siphon}. At each event, a position is assigned to the portfolio whose ledger it draws on. Ten positions changed portfolio during the deployment, so each appears once in each portfolio. Table~\ref{tab:donor_absorber} reports the results over the paper's V9 window. Entitlements use transaction prices and are revalued proportionally between events, following the prior paper.

By construction, net transfers sum to zero within every portfolio at every event. The computed sums do so to the rounding of floating-point arithmetic. Every portfolio with more than one position has both a net receiver and a net giver (Lemma~\ref{lem:partition}(ii)). Portfolio~J, with one position, has no transfer (part~(iii)). Portfolio~D holds one pool and three positions in it, and its positions transfer between themselves. Across portfolios A to J, net receivers gained \$50{,}235 in transfers from other positions over the window. Gross transfers were \$1.40 million, most of it offset by what the same positions gave up to others.

Over all of V9's portfolios, the net transfer in the window was \$58{,}952. Over V9's whole life, for which the ledger's own events suffice, it was \$122{,}940. The prior paper divides ten positions between two arrangements for emissions. With each of the ten counted as two, one under each arrangement, its first definition gives \$122{,}940, the whole-life figure to the dollar. With each counted once, it gives \$121{,}530\cite{ryan2026siphon}. This agreement suggests, but does not establish, that these are the ten positions that changed portfolio here.

For positions with a net transfer above \$100, the division into receivers and givers does not depend on the valuation. At fixed reference prices, which is this paper's own accounting, every such position keeps its sign (64 positions across the V9 portfolios). The total over all V9 portfolios is then \$62{,}053, against \$58{,}952 at transaction prices. Under proportional revaluation, an entitlement is a claim on the ledger's value and not on particular tokens, and $3.7\%$ of the value transferred passed between positions whose pools share no token. In the table, a net receiver gained more than \$1 from other positions over the window and a net giver lost more than \$1.

A position's change in value at its own rebalances is not its transfer (Lemma~\ref{lem:partition}(iv)). It also includes swap cost, operator deposits and recycled emissions. Earlier versions of this paper printed this change, taken from the operator's log for the events they sampled, as their table of donors and absorbers. Those totals appear in the final column. They sum to $+\$129{,}743$, even though transfers sum to zero within each portfolio. At the \$100 threshold, on portfolio-pool units, the earlier labels differ from the transfer-based roles in nine of the ten portfolios.


\begin{table}[!htb]
\centering
\caption{Transfers between positions over the paper's V9 window (\acr{US} dollars).}
\label{tab:donor_absorber}
\resizebox{\ifdim\width>\linewidth\linewidth\else\width\fi}{!}{%
\begin{tabular}{@{}lcccrr@{}}
\toprule
\thd{Portfolio} & \thd{Positions} & \thd{Receivers} & \thd{Givers} & \thd{Transferred} & \thd{Earlier versions}  \\
\midrule
A & \phantom{0}7 & 4 & 3 & $5{,}275$  & $+17{,}783$ \\
B & 10 & 5 & 5 & $8{,}011$  & $+1{,}285$ \\
C & \phantom{0}9 & 5 & 4 & $3{,}790$  & $+13{,}003$ \\
D & \phantom{0}3 & 2 & 1 & $134$      & $-146$ \\
E & \phantom{0}5 & 2 & 3 & $9{,}902$  & $+2{,}497$ \\
F & \phantom{0}6 & 2 & 4 & $12{,}965$ & $+24{,}004$ \\
G & \phantom{0}7 & 3 & 4 & $5{,}542$  & $+847$ \\
H & \phantom{0}6 & 4 & 2 & $4{,}024$  & $-7{,}697$ \\
I & \phantom{0}3 & 2 & 1 & $592$      & $+14{,}724$ \\
J & \phantom{0}1 & 0 & 0 & $0$        & $+63{,}443$ \\
\bottomrule
\end{tabular}}
\end{table}
\subsection{The model against the chain record}\label{sec:model_chain}

The model is compared with the V9 ledger in two respects. The first is how often an entry is zero, which Proposition~\ref{thm:chaos2B} gives as an identity of the kernel. The second is the dependence between distinct entries, which by Proposition~\ref{thm:chaos2A} vanishes in the limit. The comparisons use portfolios A to J over the fully covered window, 21~March to 14~May 2026, in which the record holds every pool's events.

\paragraph{The atom on chain (Proposition~\ref{thm:chaos2B}).}
The identity behind the atom can be checked event by event on chain. The proposition's probability calculation also assumes that marks are independent of the state, which the chain record cannot establish. The proposition identifies an entry's zero share, the share of its own events after which it is zero, with the share on which $\eta \varepsilon \geq 0$. Here $\eta$ is the entry-relative direction and $\varepsilon$ is the measured execution residual, or $-1$ where the contract skipped the swap. Table~\ref{tab:zero_share} gives both shares, with each event counted once for each of its pool's two tokens.

On the range changes of the 42 entries of portfolios A to J with at least 100 of them, the two shares differ by at most $0.0034$. The residual is itself measured from the event's own amounts. This is therefore the check of Corollary~\ref{cor:swapcorrected} in Section~\ref{sec:framework_predictions} taken entry by entry, and not a second test. Over all own events, the median difference is zero and the largest is $0.36$. The entries that differ are those whose events are mostly same-range rebalances, which lie outside the kernel.

An entry counts as zero when its balance is within the integer rounding of the mint. The table uses the $78{,}770$ rebalances whose residual is measured or whose swap was skipped at a recoverable price. On $1{,}713$ same-range events, both entries are zero in that sense, which is why the two same-range shares sum to more than one.

The atom depends on the token's position in its pool and not on the sign of the residual alone. On swap-corrected range changes, a token at the token1 position is left at zero after $73\%$ of its events and a token at token0 after $27\%$. The residual is negative on $26\%$ of the same events. Across the 42 entries, the share with $\varepsilon < 0$ is unrelated to the observed zero share (correlation $-0.16$, median difference $0.41$).
\begin{table}[!htb]
\centering
\caption{Zero share of ledger entries after their own events, observed and predicted.}
\label{tab:zero_share}
\begin{tabular}{@{}llrrrr@{}}
\toprule
\thd{Rebalances} & \thd{Position} & \thd{Events} & \thd{Observed} & $\eta\varepsilon \geq 0$ & $\varepsilon < 0$  \\
\midrule
Range change, swap      & token0 & 41{,}820 & 0.266 & 0.266 & 0.259 \\
                        & token1 & 41{,}820 & 0.734 & 0.734 & 0.259 \\
Range change, no swap   & token0 & 1{,}750  & 0.313 & 0.313 & 1 \\
                        & token1 & 1{,}750  & 0.687 & 0.687 & 1 \\
Same range              & token0 & 35{,}200 & 0.286 & 0.331 & 0.951 \\
                        & token1 & 35{,}200 & 0.763 & 0.669 & 0.951 \\
\bottomrule
\end{tabular}
\end{table}
\paragraph{Dependence between ledger entries (Proposition~\ref{thm:chaos2A}).}
Proposition~\ref{thm:chaos2A} concerns the model, whose directions are independent across events. On chain, the question is whether two entries of one portfolio's ledger move together. Over the fully covered window, the ten portfolios hold 45 entries with at least 100 own events each. For every pair of entries of one portfolio, the two balances and the two zero indicators are sampled after each of the portfolio's rebalances. The null is 200 copies of the record in which one entry's sequence is shifted in time. The shift keeps each sequence as it is and breaks its alignment with the other. A pair counts as dependent when its statistic lies outside the central $95\%$ of the copies (Table~\ref{tab:chaos_v9}).

Pairs whose tokens share a pool are negatively dependent, since each rebalance of that pool sets at least one of its two entries to zero. The 59 pairs never in one pool lie in seven of the portfolios. Of these, 14 have zero indicators dependent beyond the null, where three would be expected and at most seven in $95\%$ of independent records. Of the 14, 13 are positive. The dependence sits in tokens that share a partner, 12 of 39 such pairs against 2 of 20 with none.

The indicators are then adjusted for direction. From each is subtracted its mean over the events that share the token's position and its role at its last rebalance (bought, sold or no swap). This conditions on a direction common to pools that hold a token whose price has moved. After the adjustment, 11 of the 59 pairs remain dependent.

The balances are less clear. Over the 59 pairs, the median magnitude of their correlation is $0.039$, against a null standard deviation of $0.042$. Yet 11 of the 59 lie outside the central $95\%$ of their null, 10 of them positive. The test on balances is the noisier of the two, and the indicators are the statistic read here.

The entries of a portfolio are therefore not independent on chain. The dependence shows in their zero indicators, is positive, runs through shared tokens, and is only partly carried by the direction of the swap. The model with independent directions does not produce it. In that model, the zero indicators of two spokes are independent at every $N$ (Lemma~\ref{lem:indicators}), and the balances are dependent only at order $1/N$. Under the four laws of Section~\ref{sec:chaos_numerical}, $N$ times their correlation lies between $-0.021$ and $0$ at every $N$ from $2$ to $1024$ (supplement, Table~S2). A direction common to the pools can produce positive dependence that does not fall with $N$. It does so in the regime simulated for the model (supplement, Table~S3) and under a common price move in the simulated contract (Table~\ref{tab:contract_q1}).
\begin{table}[!htb]
\centering
\caption{Dependence between two entries of one portfolio's ledger.}
\label{tab:chaos_v9}
\begin{tabular}{@{}lrrrcc@{}}
\toprule
 & & \multicolumn{2}{c}{\thd{Balances}} & \multicolumn{2}{c}{\thd{Zero indicators, dependent}}  \\
\thd{Pair} & \thd{Pairs} & \thd{Median} $|\rho|$ & \thd{Null s.d.} & \thd{Raw} & \thd{Direction adjusted}  \\
\midrule
Tokens of one pool         & 36 & 0.25  & 0.029 & 36 & 35 \\
Sharing a partner token    & 39 & 0.043 & 0.037 & 12 & \phantom{0}9 \\
No common partner          & 20 & 0.036 & 0.051 & \phantom{0}2 & \phantom{0}2 \\
Never in one pool, all     & 59 & 0.039 & 0.042 & 14 & 11 \\
\bottomrule
\end{tabular}
\end{table}
\paragraph{Where the dependence comes from.}
To locate the dependence, the record is replayed. A replay takes the recorded events of a portfolio in order and recomputes its ledger with Theorem~\ref{thm:binding}. With every input as recorded, it returns the chain ledger to the wei. For each pair, the statistic is the covariance of its two zero indicators, divided by the standard deviation that this covariance has when one entry's sequence is shifted in time. The statistic is averaged over the 59 pairs never in one pool. On chain, the average is $0.47$. When every outcome is drawn independently, the average has standard deviation $0.13$.

A first set of replays shuffles parts of the record among the events of one pool. It shuffles first the residuals of the range changes, then their directions, then the recorded outcomes of the same-range rebalances (supplement, Table~S8). At every step that keeps the same-range outcomes, the chain value lies inside the central $95\%$ of the replays. At every step that shuffles them, no replay reaches it. The excess of positive dependence therefore passes through the same-range rebalances. These replays do not divide it further by cause. A same-range rebalance usually repeats its pool's last outcome, as shown below, so a step that keeps those outcomes keeps the outcomes of the range changes with them.

\begin{table}[!htb]
\centering
\caption{Mean standardised covariance between entries never in one pool, same-range rebalances replayed.}
\label{tab:same_range_v9}
\begin{tabular}{@{}lrrrr@{}}
\toprule
 & \multicolumn{2}{c}{\scn{59}\thd{ pairs}} & \thd{Sharing a partner} & \thd{Dependent}  \\
\cmidrule(lr){2-3}
\thd{Replayed with} & \thd{Mean} & \thd{Replay s.d.} & \thd{Mean} & \thd{pairs}  \\
\midrule
Nothing changed (chain)              & 0.47    &      & 0.75    & 14 (13) \\
Each pool's own balances             & $-0.06$ &      & $-0.08$ & 7 (5) \\
Balance from another pool removed    & $-0.09$ &      & $-0.24$ & 10 (5) \\
That balance shuffled                & 0.03    & 0.06 & 0.06    & 12 (8) \\
\quad within the week                & 0.06    & 0.05 & 0.14    & 13 (8) \\
Price of the event shuffled          & 0.46    & 0.04 & 0.73    & 14 (12) \\
\quad within the week                & 0.44    & 0.04 & 0.70    & 14 (12) \\
\midrule
Independent outcomes                 & 0.00    & 0.13 & 0.00    & 3 (2) \\
\bottomrule
\end{tabular}

\vspace{4pt}
\parbox{0.86\linewidth}{\footnotesize The replay s.d.\ is the standard deviation across replays of the mean over the 59 pairs. The dependent pairs are those beyond the time-shift null of Table~\ref{tab:chaos_v9}, with the positive ones in brackets, and for a shuffled row they are means over replays.}
\end{table}
On a same-range rebalance, the contract skips the swap on nine events in ten. The withdrawn position then already fits its range at the price of the event, so by Theorem~\ref{thm:binding} the standing balances alone determine which token remains. A balance supplied by another pool can change that outcome, and so can carry dependence between pools through the shared token.

Where a rebalance meets only the balances left by its own pool's last event, it leaves the same token as that event. On the record, this holds on all but 86 of the $21{,}917$ such events.\footnote{Of the 86, 74 are single-sided mints outside the range, which use up the balance of the token they mint, and the rest differ at the pool's rounding.} A third of the same-range rebalances whose swap is skipped meet a balance left by another pool, and on $1{,}876$ of them that balance changes the token left. Only where a rebalance meets a balance left by another pool can one pool's last outcome reach another pool's mint. Apart from 232 events that are out of range or have balances of a few wei, the price of the event enters only where the rebalance meets two non-zero balances, on $1{,}399$ events.

A second set of replays passes these rebalances through Theorem~\ref{thm:binding} in exact integer arithmetic and keeps every other event as recorded (Fig.~\ref{fig:same_range}, Table~\ref{tab:same_range_v9}). When each pool keeps its own balances, the statistic falls from $0.47$ to $-0.06$. When the balance left by another pool is replaced by one shuffled among the same pool's other same-range rebalances, taken as a share of the position, the statistic falls to $0.03$ and no replay reaches the chain value. When the price of the event is shuffled instead and the ledger is kept, the statistic stays at $0.46$. The excess of positive dependence between entries never in one pool is therefore carried by the shared balance.

A balance left by another pool in a shared token can only bring an entry to zero, so the dependence it carries directly is positive. What the shuffle removes is positive in 29 of the 39 pairs that share a partner and negative in the other 10. The coupling is the one this paper describes, the leftover of one pool drawn into the mint of another. It arises here in a setting outside (A1$'$)--(A3), under which the side is set by $\eta\varepsilon$ alone.

\begin{figure}[!htb]
\centering
\definecolor{nodeteal}{HTML}{24528A}
\definecolor{nodetealfill}{HTML}{EEF2F7}
\definecolor{nodecrimson}{HTML}{9B2335}
\definecolor{nodecrimsonfill}{HTML}{F9EFEE}
\definecolor{nodegrey}{HTML}{5A5A5A}
\begin{tikzpicture}
\begin{axis}[width=11.4cm, height=5.4cm, xmin=-0.22, xmax=0.66, ymin=0, ymax=545,
    xlabel={mean dependence between entries never in one pool}, ylabel={replays},
    xtick={-0.2,0,0.2,0.4,0.6}, tick label style={font=\footnotesize}, label style={font=\footnotesize},
    axis lines=left, tick align=outside, clip=false, ytick={0,100,200,300}]
\addplot[const plot, draw=nodeteal, fill=nodetealfill, line width=0.5pt] coordinates {(-0.18,0) (-0.16,1) (-0.14,4) (-0.12,13) (-0.10,39) (-0.08,62) (-0.06,104) (-0.04,132) (-0.02,217) (-0.00,260) (0.02,272) (0.04,266) (0.06,237) (0.08,166) (0.10,110) (0.12,54) (0.14,41) (0.16,15) (0.18,5) (0.20,2) (0.22,0) (0.24,0)} \closedcycle;
\addplot[const plot, draw=nodecrimson, fill=nodecrimsonfill, line width=0.5pt] coordinates {(0.30,0) (0.32,1) (0.34,16) (0.36,44) (0.38,81) (0.40,158) (0.42,279) (0.44,384) (0.46,383) (0.48,306) (0.50,199) (0.52,89) (0.54,37) (0.56,14) (0.58,7) (0.60,2) (0.62,0) (0.64,0)} \closedcycle;
\draw[densely dashed, line width=0.5pt] (axis cs:0.473,0) -- (axis cs:0.473,468);
\node[anchor=south, font=\scriptsize] at (axis cs:0.473,468) {the chain, 0.47};
\node[anchor=south west, font=\scriptsize, nodeteal, align=left] at (axis cs:0.105,153) {balance from another \\[-2pt] pool shuffled};
\node[anchor=south east, font=\scriptsize, nodecrimson, align=right] at (axis cs:0.415,284) {price of the \\[-2pt] event shuffled};
\draw[densely dotted, nodeteal, line width=0.5pt] (axis cs:-0.063,0) -- (axis cs:-0.063,330);
\node[anchor=south, font=\scriptsize, nodeteal, align=center] at (axis cs:-0.098,330) {each pool keeps its \\[-2pt] own balances, $-0.06$};
\end{axis}
\end{tikzpicture}
\caption{The second set of replays against the chain. Each histogram is 2{,}000 replays.}
\label{fig:same_range}
\end{figure}

The statistic is a mean over pairs, and it falls to zero without every pair doing so. The 14 pairs that are dependent on chain average $2.6$ of their null standard deviations there, and still about $1.2$ under each of the three replays that cut the link. The other 45 average $-0.2$ on chain and between $-0.35$ and $-0.52$ after. Part of the $1.2$ is owed to the 14 having been chosen on their chain value. With the balance shuffled, 12 pairs remain dependent on average, 8 of them positive, where 3 would be expected. The shared balance thus accounts for the excess of positive over negative dependence and not for every dependent pair. What remains, of both signs, is not explained.

Two further limits apply. First, the shuffle of the price has little power, so a common price at the moment of the mint is not excluded. Its contribution is bounded. On simulated records whose price at these events was set by a variable common to the pools, the shuffle removed at most $0.14$ in the settings simulated, which is under three tenths of the chain value. Second, the range changes keep their recorded outcomes in these replays. This is right to the accuracy with which the residual sets their side, although in a ledger replayed differently their leftovers would differ in size.

\section{Computational verification}\label{sec:foundry}

This section reports contract-level tests against two exchanges' position managers and machine-checked proofs in real arithmetic.

\subsection{Contract-level tests}\label{sec:foundry_tests}

The evidence of Section~\ref{sec:empirical} is the record of one manager. The contract-level tests use the unmodified position managers of two exchanges, Aerodrome Slipstream and Uniswap V3, on a fork of Base. A reference contract adds to each a ledger keyed by depositor and shared across pools. It leaves every step of the mint, the burn and the swap to the exchange. Under controlled flows, the tests cover Theorem~\ref{thm:binding}, Proposition~\ref{thm:connector} with the swap-corrected rule, and Theorem~\ref{thm:extinction}. The two corollaries are tested separately, on an integer model of the pool. Appendix~\ref{app:foundry} maps the tests to the statements (Table~\ref{tab:foundry-mapping}), and the supplement gives each scenario with its captured figures.

Without a swap, the test records the predicted liquidity and ledger entries before the transaction, then compares them with the contract's outputs. With a swap, the test evaluates Theorem~\ref{thm:binding} using the two swap legs and the resulting price, both read from the pool's own event in that transaction. It obtains the liquidity and entries returned by the contract. The theorem is evaluated on the chain record in the same way (Section~\ref{sec:results}). On both exchanges, the agreement is to the base unit of each token. It holds without a swap, with a swap, over consecutive rebalances, and when the leftover of one pool is drawn into the mint of another. After every call, the ledger equals the tokens the contract holds.

Two further tests concern the token left. They start from a position that would leave its surplus token without a swap. A swap one tenth below the amount that would leave nothing leaves that token, and a swap one tenth above it leaves the token the swap bought. This holds on both exchanges. The leftover's value is within $0.3\%$ of the expression of Proposition~\ref{thm:connector}, evaluated with the swap taken at the pool's price. The difference is the pool's fee and the swap's price impact. The integer model of the pool needs no network. On it, the proposition is checked on 400 arbitrary mints, Corollary~\ref{cor:swapfree} at displacements across and beyond the old range, and Corollary~\ref{cor:swapcorrected} at displacements inside it and at its upper bound. Each check holds to one part in $10^{9}$ of the value involved.

For Theorem~\ref{thm:extinction}, two positions in two pools that share a token are rebalanced eight times without a swap on each exchange. No rebalance raises a token's total over the positions and the ledger. Each lowers it by at most one base unit, which the pool keeps when it rounds a mint up. A token outside the rebalanced pool does not move.

The two position managers are deployed and maintained separately and key their pools differently, by tick spacing on Slipstream and by fee on Uniswap V3. Agreement on both supports reading the jump as a property of the V3 mint under a shared ledger and not of one exchange.

\subsection{Machine-checked proofs}\label{sec:lean}

The formalisation is written in Lean~4 against mathlib, so the statements are proved in a form that a machine checks. It works in real arithmetic, so the rounding of the pool's integer arithmetic lies outside it. The formalisation fully proves the paper's exact statements. These are Theorem~\ref{thm:binding}, Proposition~\ref{thm:connector} with Corollaries~\ref{cor:swapfree} and~\ref{cor:swapcorrected}, Theorem~\ref{thm:extinction} and Lemma~\ref{lem:partition}.

For the model of Section~\ref{sec:chaos}, the formalisation derives the kernel~\eqref{eq:Khom} from Theorem~\ref{thm:binding} and checks what follows from it at a single event. This covers the steps on which Lemma~\ref{lem:stationary} rests, the identity behind the atom of Proposition~\ref{thm:chaos2B} with the value of $q$, the first step of the proof of Proposition~\ref{thm:chaos2A}, parts of the statement on position size (supplement, Proposition~S1), and the Gaussian case of Proposition~\ref{thm:chaos2B} for a residual that is not conditioned to $[-K, K]$. The statements about laws are not formalised. The stationary law, the bound of Proposition~\ref{thm:chaos2A}, Propositions~\ref{cor:hetero} and~\ref{prop:commondir} and Lemma~\ref{lem:indicators} rest on the proofs of Appendix~\ref{app:chaosproof} and of the supplement. The supplement sets out the coverage statement by statement.

\section{Discussion}\label{sec:discussion}

\subsection{What holds where}\label{sec:holds}

Table~\ref{tab:scorecard} summarises where the principal statements are tested or proved. The three kinds of evidence are the model chain in simulation (Section~\ref{sec:chaos_numerical}), the contract in exact simulation and in contract-level tests (Sections~\ref{sec:model_contract} and~\ref{sec:foundry}) and the chain record (Section~\ref{sec:empirical}). A dash marks a test that does not apply. The exact identities agree with the evidence in the settings shown, subject to the exceptions noted in the table and in Section~\ref{sec:empirical}. The model's statements are proved under its assumptions. What carries over to the contract and the chain is the rule for the side of the leftover, but not, in general, the independence of the balances.

\begin{table}[!htbp]
\centering
\caption{The statements of the paper and the reach of their evidence.}
\label{tab:scorecard}
\footnotesize
\begin{tabular}{@{}>{\raggedright\arraybackslash}p{\dimexpr0.24\linewidth-1.5\tabcolsep\relax}>{\raggedright\arraybackslash}p{\dimexpr0.225\linewidth-1.5\tabcolsep\relax}>{\raggedright\arraybackslash}p{\dimexpr0.255\linewidth-1.5\tabcolsep\relax}>{\raggedright\arraybackslash}p{\dimexpr0.28\linewidth-1.5\tabcolsep\relax}@{}}
\toprule
\thd{Statement} & \thd{Model chain} & \thd{Contract} & \thd{Chain record}  \\
\midrule
\multicolumn{4}{@{}l}{{\groupface\itshape{Exact identities}}} \\
\addlinespace[2pt]
Theorem~\ref{thm:binding}, the jump & -- & to the base unit on two exchanges, with and without a swap & to the wei on $150{,}357$ of $152{,}719$ rebalances \\
\addlinespace[2pt]
Proposition~\ref{thm:connector} and its corollaries, the token left & -- & side follows $\eta\varepsilon$ at every $N$ tested; on two exchanges a swap below the exact amount leaves the token sold and one above it the token bought & the criterion on all but 80 of $82{,}728$ range changes; the swap-free rule on $57\%$ of those without a swap, disagreeing where its hypotheses do not hold \\
\addlinespace[2pt]
Theorem~\ref{thm:extinction} and Lemma~\ref{lem:partition}, conservation and transfers & -- & conserved to the pool's rounding unit on two exchanges; transfers sum to zero to rounding & transfers of \$50{,}235 net in ten portfolios; the zero sum holds by construction \\
\addlinespace[4pt]
\multicolumn{4}{@{}l}{{\groupface\itshape{Proved for the model}}} \\
\addlinespace[2pt]
Proposition~\ref{thm:chaos2A}, decoupling at rate $1/N$ & order $1/N$, at the exact covariance & $1/N$ with the hub's price fixed, at liquidity shares of $0.03$ or more; no decoupling under a common price move & balances beyond the null in 11 of 59 pairs; zero indicators dependent in 14 of 59 \\
\addlinespace[2pt]
Proposition~\ref{thm:chaos2B}, the atom & exact at every $N$ simulated & zero share equal to the share with $\eta\varepsilon \geq 0$ at every $N$ tested & the same identity within $0.0034$ on range changes, which is the check of Corollary~\ref{cor:swapcorrected} taken by entry; larger differences on same-range rebalances \\
\addlinespace[2pt]
Half-normal slack side for a Gaussian residual (Remark after Proposition~\ref{thm:chaos2B}) & variance within $0.5\%$ & does not carry over & not tested \\
\addlinespace[2pt]
Proposition~\ref{cor:hetero}, unequal spokes, under~\eqref{eq:heterobox} & atom on every spoke satisfying~\eqref{eq:heterocond}; the runs lie outside~\eqref{eq:heterobox} & not tested & not tested \\
\addlinespace[2pt]
Proposition~\ref{prop:commondir} and Lemma~\ref{lem:indicators}, a common direction & indicator covariance of the lemma matched; the proposition's variance not computed & not tested; indicator correlation $0.25$ seen under a common price move & the excess of positive dependence is carried by the shared balance at same-range rebalances; the contribution of a common direction is not resolved \\
\addlinespace[2pt]
Position size (Section~\ref{sec:possize}; supplement, Proposition~S1) & contraction rate within $1.5\%$ for $N \geq 16$ & not established & not tested \\
\bottomrule
\end{tabular}
\end{table}

\subsection{What it means}\label{sec:means}

For each event, the Master Equation determines the new entries, and thus reconciles the ledger, from the withdrawn amounts, standing entries, swap legs, new range and mint price (Table~\ref{tab:cross_version}). A failure to reconcile points to an omitted or mismeasured flow, an incorrect reconstruction of the ledger or a departure from the assumed accounting. The non-closing events in the early versions are consistent with rewards or ledger movements that the event data do not record.

Forecasting the leftover before the transaction also requires the contract's swap rule, since the swap sets both the amounts offered to the mint and the price at which the mint occurs. Theorem~\ref{thm:binding} is an identity at the mint and not a forecast. With the swap's realised amounts and the pre-swap price, it names the token left no better than a guess from each pool's own mix (Section~\ref{sec:results}). The contract's swap rule does carry information before the swap. Its first-order form gives the sign of the execution residual on $98.2\%$ of measured events (Appendix~\ref{app:chaosproof}).

Evaluating positions separately requires entitlement accounting. The mint takes whatever the ledger holds, so a position can deploy value owed to another. Over the window studied, the positions of the ten portfolios did so for \$50{,}235 net. For one owner, these transfers affect only the attribution of value between positions. If a shared ledger served different owners, they would also move wealth between owners, as the prior paper notes\cite{ryan2026siphon}. In either case, they have to be tracked through entitlements, since the change in a position's own value does not measure them.

Balances under a shared ledger need not become independent as the portfolio grows. In the model they do, in the limit and more nearly so as the number of pools grows. That is a property of the model under its assumptions, and two features break it. First, a common price move aligns the rebalance directions of several pools, so the spokes of the simulated contract do not decouple. Second, at a rebalance that keeps its range, the token left is set by the standing balances. A balance left by another pool in a shared token can therefore change which token is left and carry dependence between pools.

The chain data show dependence in the zero indicators in 14 of 59 pairs that never share a pool, 13 of them positive. Dependence is less clear in the balances, appearing in 11 of 59 pairs. The replays attribute the excess of positive dependence to balances passed between pools at same-range rebalances. They do not account for every dependent pair; the remaining dependence, of both signs, is unexplained.
\subsection{Limitations}\label{sec:limitations}

The empirical anchor is one operator on one chain. The cross-version replication of Section~\ref{sec:cross_version} shows that the identity is not a peculiarity of any one contract version. All four versions were operated by the same operator on Aerodrome Slipstream, and they are the population on which the evidence of this paper rests. The prior paper reports a second deployment of the architecture, run by another operator, on which the ledger identity closes on every transaction it examined\cite{ryan2026siphon}. There, the mint of Theorem~\ref{thm:binding} on the amounts held after the swap reproduces the leftover, in a check not specified in advance. The test that was declared in advance, of a rule for the amounts a separate swap leaves, did not meet its thresholds. No figure from that deployment is used here.

Propositions~\ref{thm:chaos2A} to~\ref{prop:commondir} are statements about a model and not about the deployed contract. Three departures of the contract were found to matter. First, directions are not independent when the price of a shared token moves. Under such a move, the spokes of the simulated contract do not decouple at any $N$ tested (Section~\ref{sec:model_contract}). Second, the model has no same-range rebalance. That is where the agreement of the atom weakens on chain and where the shared balance carries the excess of positive dependence between entries (Section~\ref{sec:model_chain}). Third, the contract's execution residual is not inside the bound of (A1$'$). Taken by direction from the V9 rebalances with a measured swap, without the clip and with the skipped swaps at $\varepsilon = -1$, the execution residual can change the imbalance's sign on about $1\%$ of events. On those events, the side of the leftover depends on the state, and the atom of the simulated chain falls short of $q$ (Section~\ref{sec:chaos_numerical}).

The chain with $L$ held fixed does not conserve value. Against the update~\eqref{eq:Lupdate}, in which the mint sets the size, each of its events adds $|A_\eta|\,|\varepsilon| - H - B_k$ to the total value. The mean of this discrepancy is zero under the stationary law, but the discrepancy is not small event by event. In the simulations of Section~\ref{sec:chaos_numerical}, its standard deviation is $0.05\,Lw$ to $0.26\,Lw$ across the four laws. When both the hub and the selected spoke are empty, which occurs on between $18$ and $25\%$ of events at $N = 1024$, it is $Lw\,|\varepsilon|$. The fixed-$L$ results are statements about balances measured against a position of fixed size. Section~\ref{sec:possize} lets liquidity change at each mint.

The rate of Proposition~\ref{thm:chaos2A} is a statement about large $N$. Its bound for one spoke falls below the trivial bound $2 b^*$ only from $N = 7$ pools when $\varsigma = 0.3$ and from $N = 25$ when $\varsigma = 0.5$. The V9 portfolios hold between one and ten pools, and not all of their graphs are stars. At small $N$, the bound remains valid but is uninformative, and the exact covariance of the Remark after the proposition is the description of the model to use. The condition~\eqref{eq:heterobox} of Proposition~\ref{cor:hetero} is restrictive. With one bound $\kappa$, the largest position must be less than $(1 - \kappa)/\kappa$ times the smallest, and the V9 portfolios are not shown to satisfy it.

Three questions remain unresolved. The first is the dependence between entries that the shared balance does not account for. It is of both signs, and a price common to the pools at the moment of a same-range mint is bounded and not excluded. The second is how position sizes settle across pools of different depth under the contract, since the simulation did not meet the criterion declared for it. The third is the law of the slack side on chain, which is not compared with a closed form. The half-normal law is a property of a Gaussian residual, which the contract does not produce (Section~\ref{sec:model_contract}).

\subsection{Future work}\label{sec:future}

The most immediate extension is a kernel for the same-range rebalance. At such a rebalance, the contract skips the swap on most events and the standing balances set the side of the leftover. The side is the pool's own last side unless a balance from another pool intervenes (Section~\ref{sec:model_chain}). The property on which the proof rests, that a change to one entry stays in one entry, is therefore lost. The contraction~\eqref{eq:contraction} still holds, and the framework of Andreis, Dai Pra and Fischer\cite{andreis2018} for interacting systems with simultaneous jumps is a candidate approach. On the empirical side, the first candidate for the dependence that the shared balance does not carry is a direction common to the pools, whose sign for a pair of entries depends on where their tokens sit in their pools.

For the model, Appendix~\ref{app:chaosproof} and Section~7 of the supplement list what the proofs leave open. Beyond that list stands the joint stationary law on portfolio graphs other than a star, which is where the deployed portfolios sit. The simulations also suggest that the rate survives for unequal spokes beyond the condition~\eqref{eq:heterobox}.

For the contract, position-size dynamics need a dedicated test that starts from unequal own values and measures whether, and how fast, the differences contract. The most informative extension would be to repeat the measurements of Section~\ref{sec:empirical} in full on an independently operated deployment of the same architecture.\footnote{A companion paper\cite{ryan2026fingerprint} recovers rebalancing trigger fractions from public chain data across 81 distinct operators. It does not replace a replication of the ledger measurements of Section~\ref{sec:empirical}, which rest on one contract family.}

\section{Conclusion}\label{sec:conclusion}

The dust ledger of a multi-pool concentrated liquidity manager obeys an exact law at each rebalance. The Master Equation gives the jump of the ledger, the token the rebalance leaves and the conservation across pools that follows, extending the geometric residual of the prior paper\cite{ryan2026siphon}. Across four deployed versions of one manager, the equation closes to the wei on $150{,}357$ of $152{,}719$ recorded rebalances; the remaining discrepancies are consistent with rewards or ledger movements that the event data do not record. These identities hold for any portfolio whose leftovers are kept in such a ledger, and require no stochastic assumptions.

A model adds a description of the ledger's joint law. Under the assumptions of Section~\ref{sec:chaos}, the spokes of a hub-spoke portfolio decouple at rate $1/N$. Each has an atom at zero whose mass is the probability that direction and residual agree in sign. When the model is compared with an exact simulation of the contract and with the chain record, the rule determining which token is left, on which the atom rests, is supported wherever the model's kernel applies. Independence does not carry over in general. It is not found in the simulated contract where a common price move aligns the pools' rebalance directions, and on chain the zero indicators of the entries are dependent. At the rebalances that keep their range, which the model leaves out, the shared balance itself carries the excess of positive dependence. Extending the kernel to these events is a principal open problem.

The coupling follows from the V3 amount equations alone once a contract keeps leftovers in a balance shared across positions. Each mint can draw on value owed to other positions. The neutral accounting rule of the prior paper\cite{ryan2026siphon} caps each mint at the rebalancing position's own equity after swap costs. Under that rule, positions can draw on a shared balance without one deploying value owed to another.

\section*{Data and code availability}

The public repository holds the Foundry suite with its recorded output, the Lean~4 formalisation, a simulation of the model chain under its three synthetic laws, a supplement to this paper and a note of what changed from the versions of May and July 2026. It is on GitHub at \url{https://github.com/g4nc0r/hidden-microstructure}, with an archival snapshot on Zenodo at \url{https://doi.org/10.5281/zenodo.20115374}.

The supplement contains the tables behind the simulations of Section~\ref{sec:chaos_numerical}, Table~\ref{tab:modelift} by portfolio, the auxiliary indexers, the first set of replays, the contract-level tests by scenario, the coverage of the formalisation by statement, the proofs of Propositions~\ref{cor:hetero} and~\ref{prop:commondir} and Lemma~\ref{lem:indicators}, and the statement on position size with its proof.

The chain record of Section~\ref{sec:empirical} is derived from public contract and pool events on Base. Its derived files, the measured residual law, the contract simulations, the replay scripts and the test specifications fixed in advance are not part of the current deposit. The author intends to release them with those of the prior paper\cite{ryan2026siphon}.

\phantomsection\addcontentsline{toc}{section}{References}
\begin{thebibliography}{99}

\bibitem{ryan2026siphon}
Ryan, K.R. ``The Geometric Siphon: Existence, Equilibrium, and Directional Properties of the Residual in Concentrated Liquidity Portfolios.'' \emph{\acr{SSRN} preprint 6686798} (2026) \url{https://ssrn.com/abstract=6686798}.

\bibitem{ryan2026om}
Ryan, K.R. ``Operator \& Quantisation Microstructure: The Frequency and Value Channels of Concentrated Liquidity Dust Flows.'' \emph{\acr{SSRN} preprint 7166739} (2026) \url{https://ssrn.com/abstract=7166739}.

\bibitem{pi2026convexity}
P., I. ``Convexity and Topology in Cross-Pool Dust Redistribution: Structural Edge in Active \acr{CL} Management.'' \emph{\acr{SSRN} preprint 6509938} (2026) \url{https://ssrn.com/abstract=6509938}.

\bibitem{cartea2024}
Cartea, \'A., Drissi, F., Monga, M. ``Decentralized Finance and Automated Market Making: Predictable Loss and Optimal Liquidity Provision.'' \emph{\acr{SIAM} Journal on Financial Mathematics} \textit{15}(3), 931--959 (2024) \url{https://doi.org/10.1137/23M1602103}.

\bibitem{milionis2022b}
Milionis, J., Moallemi, C.\,C., Roughgarden, T., Zhang, A.\,L. ``Automated Market Making and Loss-Versus-Rebalancing.'' arXiv:2208.06046 (2022) \url{https://arxiv.org/abs/2208.06046}.

\bibitem{heimbach2022}
Heimbach, L., Schertenleib, E., Wattenhofer, R. ``Risks and Returns of Uniswap V3 Liquidity Providers.'' In \emph{Proceedings of the 4th \acr{ACM} Conference on Advances in Financial Technologies (\acr{AFT} 2022)}, pp.~89--101 (2022) \url{https://doi.org/10.1145/3558535.3559772}.

\bibitem{sznitman1991}
Sznitman, A.-S. ``Topics in Propagation of Chaos.'' In \emph{\'Ecole d'\'Et\'e de Probabilit\'es de Saint-Flour XIX, 1989}, Lecture Notes in Mathematics 1464, pp.~165--251, Springer (1991) \url{https://doi.org/10.1007/BFb0085169}.

\bibitem{graham1993}
Graham, C., M\'el\'eard, S. ``Propagation of chaos for a fully connected loss network with alternate routing.'' \emph{Stochastic Processes and their Applications} \textit{44}(1), 159--180 (1993) \url{https://doi.org/10.1016/0304-4149(93)90043-4}.

\bibitem{chaintron2022}
Chaintron, L.-P., Diez, A. ``Propagation of chaos: a review of models, methods and applications. I. Models and methods.'' \emph{Kinetic and Related Models} \textit{15}(6), 895--1015 (2022). arXiv:2203.00446 \url{https://arxiv.org/abs/2203.00446}.

\bibitem{kelly1979}
Kelly, F.\,P. \emph{Reversibility and Stochastic Networks}. Wiley, Chichester (1979). Reprinted by Cambridge University Press, 2011.

\bibitem{carmona2018}
Carmona, R., Delarue, F. \emph{Probabilistic Theory of Mean Field Games with Applications II: Mean Field Games with Common Noise and Master Equations}. Probability Theory and Stochastic Modelling, vol.~84, Springer (2018) \url{https://doi.org/10.1007/978-3-319-56436-4}.

\bibitem{cardaliaguet2019}
Cardaliaguet, P., Delarue, F., Lasry, J.-M., Lions, P.-L. \emph{The Master Equation and the Convergence Problem in Mean Field Games}. Annals of Mathematics Studies 201, Princeton University Press (2019) \url{https://doi.org/10.2307/j.ctvckq7qf}.

\bibitem{adams2021v3}
Adams, H., Zinsmeister, N., Salem, M., Keefer, R., Robinson, D. ``Uniswap v3 Core.'' \emph{Uniswap Labs Technical Whitepaper} (2021) \url{https://uniswap.org/whitepaper-v3.pdf}.

\bibitem{andreis2018}
Andreis, L., Dai Pra, P., Fischer, M. ``McKean--Vlasov limit for interacting systems with simultaneous jumps.'' \emph{Stochastic Analysis and Applications} \textit{36}(6), 960--995 (2018). arXiv:1704.01052 \url{https://arxiv.org/abs/1704.01052}.

\bibitem{ryan2026fingerprint}
Ryan, K.R. ``Operator Fingerprinting: Recovering and Defending the Hidden Rebalancing Policy of Concentrated Liquidity Managers.'' \emph{\acr{SSRN} preprint 7202340} (2026) \url{https://ssrn.com/abstract=7202340}.

\end{thebibliography}

\appendix
\titleformat{\section}[block]{\centering\titleface\fontsize{17pt}{20pt}\selectfont}{Appendix \thesection}{0.7em}{}

\section{Notation}\label{app:notation}

The notation for a single rebalance follows the prior paper\cite{ryan2026siphon}, with two exceptions. Here $D$ denotes the value of a standing balance; in the prior paper it denotes the dust flow of an event. Here $\sigma_x$ and $\sigma_y$ denote the swap legs; the prior paper uses $\sigma$ for volatility. Section~\ref{sec:chaos} introduces notation specific to this paper. The table lists symbols used in more than one place; symbols used only in a proof or paragraph are defined where they appear.

\begingroup
\renewcommand{\arraystretch}{1.0}%
\setlength{\extrarowheight}{1pt}%
\setlength{\aboverulesep}{1pt}%
\setlength{\belowrulesep}{1pt}%
\setlength{\abovecaptionskip}{4pt}%
\setlength{\belowcaptionskip}{2pt}%
\begin{longtable}{@{}lL{10.4cm}@{}}
\caption{Notation.}\label{tab:notation}\\
\toprule
\thd{Symbol} & \thd{Meaning}  \\
\midrule
\endfirsthead
\multicolumn{2}{l}{\textit{Table~\ref{tab:notation} (continued).}}\\
\toprule
\thd{Symbol} & \thd{Meaning}  \\
\midrule
\endhead
\midrule
\multicolumn{2}{r}{\textit{Continued on next page.}}\\
\endfoot
\bottomrule
\endlastfoot
\multicolumn{2}{@{}l}{{\groupface\itshape{Portfolio and events}}}\\*
$T_i$ & Tokens of the portfolio. \\
$T_a$, $T_b$ & The tokens of one pool, at the token0 and token1 positions. \\
$P_k$ & The pools, with one position each. \\
$\tau^\pm$ & Just after and just before a rebalance event. \\
\addlinespace[5pt]
\multicolumn{2}{@{}l}{{\groupface\itshape{Pool geometry and the mint}}}\\*
$s$ & Sqrt price at the rebalance. \\
$s_a$, $s_b$; $s_a'$, $s_b'$ & Bounds of the old range, and of the new. \\
$s_0$ & Centre of the old range in tick space, $s_a s_b = s_0^2$. \\
$\ell$, $h$ & Logarithmic displacement $\ln(s/s_0)$ and half-width $\ln(s_b/s_0)$. \\
$w$ & Width of a range, $s_b - s_a$. \\
$L$, $L_{\symup{new}}$ & Position liquidity before and after the rebalance. \\
$x_w$, $y_w$ & Amounts withdrawn, eq.~\eqref{eq:oldamount}. \\
$a_x$, $a_y$ & V3 amounts per unit of liquidity (Appendix~\ref{app:v3formulas}). \\
$g_x$, $g_y$ & The same at the new range. \\
$\hat x$, $\hat y$ & Mint inputs, eq.~\eqref{eq:avail}. \\
$\hat W$ & Value of the mint inputs in units of $T_b$. \\
$\hat\omega_a$, $\omega'_a$ & Share of $T_a$ in the mint inputs, and in what the new range requires. \\
$V$ & Withdrawn value. \\
$\alpha(\ell)$ & Share of $V$ that a swap-free rebalance leaves from an empty balance (Corollary~\ref{cor:swapfree}). \\
\addlinespace[5pt]
\multicolumn{2}{@{}l}{{\groupface\itshape{Ledger, swap and transfers}}}\\*
$D_{\symup{pool},i}$ & Ledger entry of token $T_i$. \\
$x_{\symup{dust}}$, $y_{\symup{dust}}$ & Entries of the pool's two tokens before the rebalance, primed after it, eq.~\eqref{eq:newdust}. \\
$D$ & Value of a standing balance. \\
$\sigma_x$, $\sigma_y$ & Swap legs, positive for what the contract pays away. \\
$\varepsilon$ & Execution residual, the swap's amount divided by the exact amount, less one. \\
$\eta$ & Direction of an event, $+1$ or $-1$. In the frame of the pool it is $+1$ when the swap buys token0. In the frame of one ledger entry it is positive when the swap sells that entry's token. \\
$G_k$, $U_k$ & Entitlement and equity of position $k$. \\
$\Xi_j$, $\xi_k$ & Transfer at a rebalance of position $j$, and net transfer to position $k$ (Lemma~\ref{lem:partition}). \\
$\Delta V_k$ & Change of a position's value at its own rebalance. \\
\addlinespace[5pt]
\multicolumn{2}{@{}l}{{\groupface\itshape{The model (Section~\ref{sec:chaos})}}}\\*
$N$ & Number of spokes. \\
$H$, $B_k$ & Ledger entries of the hub and of spoke $k$. \\
$p_+$, $p_-$ & Probabilities of the two directions. \\
$\nu_+$, $\nu_-$ & Laws of the residual in each direction. \\
$K$ & Bound on the size of the residual. \\
$c_0$ & Positive scale factor of the leftover, relative to the size of the residual times the imbalance. One for the kernel as derived. \\
$\kappa$ & $c_0 K$, less than one half. \\
$\kappa_1$ & $c_0\,\mathbb{E}|\varepsilon|$. \\
$\varsigma$ & $c_0^{1/2}\,\mathbb{E}\bigl[|\varepsilon|^{1/2}\bigr]$. \\
$\sigma_{\symup{slip}}$ & Standard deviation of the Gaussian residual law before it is conditioned. \\
$\sigma_{\symup{hn}}$, $\mathrm{HN}(\sigma_{\symup{hn}})$ & Scale of the half-normal law, and that law (Remark after Proposition~\ref{thm:chaos2B}). \\
$A_\eta$ & Imbalance the swap is asked to remove, eq.~\eqref{eq:imbalance}. \\
$b^*$ & Bound of the invariant set, $\kappa Lw/(1 - \kappa)$, or~\eqref{eq:heterobox} for unequal spokes. \\
$X^N$, $\pi^N$ & The balances $(H, B_1, \ldots, B_N)$ and their stationary law. \\
$\pi^N_{(m)}$, $\pi^N_{(0,m)}$ & Marginal of $\pi^N$ on $m$ spokes, and on the hub and $m$ spokes. \\
$\mu$, $\pi_H$ & Laws of a spoke and of the hub in the limit, eq.~\eqref{eq:McKVfixed}. \\
$\bar b$, $\bar h$ & Means of $\mu$ and $\pi_H$. \\
$q$ & Probability mass at zero of a spoke's law. \\
$\rho_+$ & Leading-order law of a spoke's slack side (Proposition~\ref{thm:chaos2B}). \\
$W_1$ & Wasserstein distance for the sum of absolute differences. \\
$\chi_\pm$ & $c_0\,\mathbb{E}\bigl[(\eta\varepsilon)^\pm\eta\bigr]$. \\
$\gamma_\pm$ & $c_0\,\mathbb{E}\bigl[(\eta\varepsilon)^\pm\bigr]$, with $\gamma_+$ the mean leftover that lands on the hub. \\
$\varpi_k$ & Probability that an event is of pool $k$. \\
$q_k$, $\mu_k$ & Mass at zero and law of spoke $k$ (Proposition~\ref{cor:hetero}). \\
$\varsigma^*$, $\varpi^*$ & A star marks the largest value over the pools, or over the states of an environment. \\
$e_t$ & Environment state at event $t$. \\
$\rho_t$, $\pi_H^t$, $\mu^t$ & Laws of the balances, of the hub and of a spoke at time $t$ given the environment (Proposition~\ref{prop:commondir}). \\
$I_a(t)$ & Indicator that the entry of token $a$ is zero after event $t$ (Lemma~\ref{lem:indicators}). \\
$Z_k$ & Own value $L_k w + B_k$ of position $k$ (Section~\ref{sec:possize}). \\
\end{longtable}
\endgroup

\section{V3 amount formulas}\label{app:v3formulas}

The V3 amount functions for a position over $[s_a, s_b]$ at sqrt price $s$ are
\begin{align}
a_x(s, s_a, s_b) &= \begin{cases}
1/s_a - 1/s_b & \text{if } s \leq s_a, \\
1/s - 1/s_b & \text{if } s \in [s_a, s_b], \\
0 & \text{if } s \geq s_b,
\end{cases} \\[6pt]
a_y(s, s_a, s_b) &= \begin{cases}
0 & \text{if } s \leq s_a, \\
s - s_a & \text{if } s \in [s_a, s_b], \\
s_b - s_a & \text{if } s \geq s_b.
\end{cases}
\end{align}
These are normalised to $L = 1$. For liquidity $L$, the corresponding token amounts are $L\,a_x(s, s_a, s_b)$ and $L\,a_y(s, s_a, s_b)$. As $s$ increases, $a_x$ is non-increasing and $a_y$ is non-decreasing. At or beyond either bound of the range, the position holds one token only.

\section{Proof of Theorem~\ref{thm:binding}}\label{app:masterproof}

The mint chooses $L_{\symup{new}}$ to maximise deployed liquidity subject to the available token balances. This appendix solves the resulting linear programme with the price inside, below and above the new range.

\subsection*{The mint as a linear programme}

Let
\[
g_x = a_x(s, s_a', s_b'), \qquad g_y = a_y(s, s_a', s_b').
\]
By Appendix~\ref{app:v3formulas}, both satisfy $g_x \geq 0$ and $g_y \geq 0$. The mint problem is
\begin{equation}\label{eq:lp}
L_{\symup{new}} = \max\bigl\{L \geq 0 : L\,g_x \leq \hat x \text{ and } L\,g_y \leq \hat y\bigr\}.
\end{equation}
This is a one-dimensional linear programme in $L$. Since the objective is increasing in $L$, the optimum is attained when at least one constraint binds. The two candidate binding values are
\begin{equation}
L_x^* = \hat x / g_x \text{ (when } g_x > 0\text{)}, \qquad L_y^* = \hat y / g_y \text{ (when } g_y > 0\text{)}.
\end{equation}
Hence
\[
L_{\symup{new}} = \min(L_x^*, L_y^*),
\]
with the binding side determined by whichever candidate is smaller. A token whose requirement is zero imposes no bound, and its candidate is left out of the minimum.

\subsection*{Sub-case (i): price inside the range}

Suppose $s \in (s_a', s_b')$. Then
\[
g_x = 1/s - 1/s_b' > 0, \qquad g_y = s - s_a' > 0,
\]
so both candidate liquidities are finite. If $L_x^* \leq L_y^*$, the $x$-constraint binds, so
\begin{equation}
L_{\symup{new}} = \hat x / g_x, \quad x_{\symup{dust}}' = 0, \quad y_{\symup{dust}}' = \hat y - L_{\symup{new}}\,g_y = \hat y - \hat x\,(g_y / g_x).
\end{equation}
If $L_y^* < L_x^*$, the $y$-constraint binds, so
\begin{equation}
L_{\symup{new}} = \hat y / g_y, \quad x_{\symup{dust}}' = \hat x - L_{\symup{new}}\,g_x, \quad y_{\symup{dust}}' = 0.
\end{equation}
In the swap-free regime $\hat x = x_w + x_{\symup{dust}}$ and $\hat y = y_w + y_{\symup{dust}}$, so the binding side is determined by the supply ratio $\hat x / \hat y$ versus $g_x / g_y$, with $x_{\symup{dust}}' \cdot y_{\symup{dust}}' = 0$ in either case.

\subsection*{Sub-cases (ii) and (iii): price at or outside a bound}

Sub-cases (ii) and (iii) are $s \leq s_a'$ and $s \geq s_b'$. They lie outside the hypothesis of Theorem~\ref{thm:binding} and are given for completeness.

If $s \leq s_a'$, then
\[
g_x = 1/s_a' - 1/s_b' > 0, \qquad g_y = 0.
\]
Only the $x$-constraint is active, so
\[
L_{\symup{new}} = \hat x / g_x, \quad x_{\symup{dust}}' = 0, \quad y_{\symup{dust}}' = \hat y.
\]
The new position is therefore single-sided in $T_a$, the price being at or below the lower bound of the range.

If $s \geq s_b'$, the argument is symmetric. In that case
\[
g_x = 0, \qquad g_y = s_b' - s_a' > 0,
\]
so only the $y$-constraint is active and
\[
L_{\symup{new}} = \hat y / g_y, \quad x_{\symup{dust}}' = \hat x, \quad y_{\symup{dust}}' = 0.
\]

\subsection*{Swap-corrected inputs}

With a non-zero swap correction $(\sigma_x, \sigma_y)$ as in Section~\ref{sec:swap}, the inputs become
\[
\hat x = x_w + x_{\symup{dust}} - \sigma_x, \qquad \hat y = y_w + y_{\symup{dust}} - \sigma_y.
\]
The \acr{LP} structure in~\eqref{eq:lp} is unchanged, so the same binding-side and leftover formulas apply to the corrected inputs. The ratio-preserving locus of the prior paper (the corollary to its Theorem~1) becomes
\[
\hat x / \hat y = g_x / g_y,
\]
which is the ratio the new range requires. An exact swap brings the inputs to it, and a swap that misses it leaves one token unused, as before. \proofqed

\section{Proofs of Proposition~\ref{thm:connector} and Corollaries~\ref{cor:swapfree} and~\ref{cor:swapcorrected}}\label{app:connectorproof}

The proposition and both corollaries follow from Theorem~\ref{thm:binding}. All values below are measured in units of $T_b$ at the mint price $s$. The symbols $\hat W$, $\hat\omega_a$ and $\omega'_a$ are defined in Section~\ref{sec:connectorrule}.

\subsection*{Share form and the swap-free rule}

\emph{Proposition~\ref{thm:connector}.} By Theorem~\ref{thm:binding}, $L_{\symup{new}} = \min(\hat x/g_x,\ \hat y/g_y)$, and
\begin{equation}\label{eq:bindingratio}
\frac{\hat x}{g_x} \leq \frac{\hat y}{g_y}
\iff s^2\hat x\,g_y \leq \hat y\,s^2 g_x
\iff \frac{\hat\omega_a}{1 - \hat\omega_a} \leq \frac{\omega'_a}{1 - \omega'_a}
\iff \hat\omega_a \leq \omega'_a,
\end{equation}
with $\hat\omega_a = 1$ when $\hat y = 0$. If $\hat\omega_a \leq \omega'_a$, then $x_{\symup{dust}}' = 0$ and $y_{\symup{dust}}' = \hat y - \hat x\,g_y/g_x$, which in value is
\[
\hat W\,(1 - \hat\omega_a) - \hat W\,\hat\omega_a\,\frac{1 - \omega'_a}{\omega'_a} = \hat W\,\frac{\omega'_a - \hat\omega_a}{\omega'_a}.
\]
The other case follows by exchanging the two tokens. The expressions agree at $\hat\omega_a = \omega'_a$, where both leftovers vanish.

\emph{Corollary~\ref{cor:swapfree}(i).} For a new range with $s_a' s_b' = s^2$,
\[
s^2 g_x = s - \frac{s^2}{s_b'} = s - s_a' = g_y,
\]
so $\omega'_a = 1/2$. With no swap and no standing balance, $\hat\omega_a$ is the share $\omega$ of $T_a$ in the withdrawn position. For $s \in [s_a, s_b]$ and $s_a s_b = s_0^2$, the V3 amount functions give
\[
s^2 a_x(s, s_a, s_b) - a_y(s, s_a, s_b) = s_a - \frac{s^2}{s_b} = \frac{s_0^2 - s^2}{s_b},
\]
which has the sign of $s_0 - s$. For $s \geq s_b$ the position holds $T_b$ alone and $\omega = 0$, and for $s \leq s_a$ it holds $T_a$ alone and $\omega = 1$. Hence $\omega > 1/2$ for $s < s_0$, $\omega = 1/2$ at $s_0$ and $\omega < 1/2$ for $s > s_0$, and Proposition~\ref{thm:connector} with $\omega'_a = 1/2$ gives the side.

It remains to compute the value of the leftover. By Proposition~\ref{thm:connector} it is $V\,|1 - 2\omega|$, which outside the old range is $V$. Inside it, $V\,|1 - 2\omega| = L\,|s_0^2 - s^2|/s_b$ and $V = L\,(2s - s^2/s_b - s_0^2/s_b)$, and dividing numerator and denominator by $s\,s_0$ gives
\[
\alpha(\ell) = \frac{|s_0/s - s/s_0|}{2 s_b/s_0 - s/s_0 - s_0/s} = \frac{\sinh|\ell|}{e^{h} - \cosh\ell}.
\]

\emph{Corollary~\ref{cor:swapfree}(ii).} Let $\omega_d$ be the share of $T_a$ in the standing balance. Then
\[
\hat\omega_a - 1/2 = \frac{V\,(\omega - 1/2) + D\,(\omega_d - 1/2)}{V + D},
\]
with $|\omega_d - 1/2| \leq 1/2$ and, by part~(i), $|\omega - 1/2| = \alpha(\ell)/2$. If $D < \alpha(\ell)\,V$ the numerator has the sign of $\omega - 1/2$, and Proposition~\ref{thm:connector} gives the side of part~(i). If the balance is held wholly in the token that the mint of part~(i) uses up then $D\,|\omega_d - 1/2| = D/2$ with the opposite sign, so the numerator vanishes at $D = \alpha(\ell)\,V$ and changes sign beyond it.

\subsection*{Swap-corrected rebalances}

\emph{Corollary~\ref{cor:swapcorrected}.} The swap is at the price $s$, so $\hat W = V$. The bound $\varepsilon \leq 1/\alpha(\ell)$ keeps the amount of the sold token offered to the mint non-negative, as Proposition~\ref{thm:connector} requires. The swap raises the supplied share of the token it buys by $Q/V$. Before the swap that share is short of $1/2$ by $\alpha(\ell)/2$, so after it the supplied share differs from $1/2$ by
\[
\Bigl|\frac{\alpha(\ell)}{2} - \frac{Q}{V}\Bigr| = \frac{\alpha(\ell)}{2}\,|\varepsilon|,
\]
on the side of Corollary~\ref{cor:swapfree}(i) if $Q < \tfrac{1}{2}\alpha(\ell)V$ and on the other side if $Q > \tfrac{1}{2}\alpha(\ell)V$. By Proposition~\ref{thm:connector} with $\omega'_a = 1/2$ the leftover is worth $2V$ times that difference, which is $\alpha(\ell)\,V\,|\varepsilon|$. \proofqed

\paragraph{Remarks.} For ranges of equal width in sqrt price in place of tick space, the side statement of Corollary~\ref{cor:swapfree}(i) holds with $s_0$ replaced by the arithmetic midpoint of the old range. If the new range is not centred on the mint price, $\omega'_a \neq 1/2$ and the neutral price is the one at which the withdrawn share equals $\omega'_a$, so a small displacement can leave either token. For $\varepsilon \leq 0$, Corollary~\ref{cor:swapcorrected} is the displacement identity of the prior paper\cite{ryan2026siphon}.

Agreement in sign at every event does not in general fix the sign of a sample correlation. Across a sample of swap-free rebalances that satisfy Corollary~\ref{cor:swapfree}(i), the signed leftover has the sign of the displacement at every event. That fixes the sign of a sample covariance, and of a rank correlation, when the displacement takes two values of opposite sign. When it takes many values it fixes neither, since both then depend on the ordering of the leftovers within each direction.

\section{Proofs of Theorem~\ref{thm:extinction} and Lemma~\ref{lem:partition}}\label{app:extinctionproof}

Theorem~\ref{thm:extinction} follows from per-token conservation. At each rebalance, every position is valued at the sqrt price of its own pool, the event price. Lemma~\ref{lem:partition} shows that transfers between positions in a ledger sum to zero at every rebalance. The division into net receivers and net givers then follows.

\subsection*{Proof of Theorem~\ref{thm:extinction}}

Consider a rebalance of pool $P_k$ at time $\tau$, at the event price $s = s(\tau)$, in the swap-free regime $\sigma_x = \sigma_y = 0$. If a token $T_i \notin P_k$, then the dust-ledger entry $D_{\symup{pool},i}$ is unchanged by~\eqref{eq:master}, and no position outside $P_k$ moves any amount of $T_i$. Hence the per-token total $M_i$ defined in~\eqref{eq:masscons} is unchanged.

Now suppose $T_i \in P_k$, and write $T_i = T_a$ without loss of generality. By the V3 amount equations and Theorem~\ref{thm:binding}, position $k$'s amount of $T_a$ jumps from
\[
L^{(k)} a_x(s, s_a, s_b) = x_w
\]
to
\[
L_{\symup{new}}^{(k)}\,g_x = \hat x - x_{\symup{dust}}',
\]
while the corresponding dust-ledger entry jumps from $x_{\symup{dust}}$ to $x_{\symup{dust}}'$. In the swap-free regime $\hat x = x_w + x_{\symup{dust}}$, so
\[
L_{\symup{new}}^{(k)}\,g_x + x_{\symup{dust}}' = (\hat x - x_{\symup{dust}}') + x_{\symup{dust}}' = \hat x = x_w + x_{\symup{dust}} = L^{(k)} a_x(s, s_a, s_b) + x_{\symup{dust}}.
\]
Thus the total allocation of token $T_a$ across the moving position and the dust ledger is preserved across the rebalance at price $s$.

All other positions containing $T_a$ are unchanged by the event and are valued at the price of their own pool. Summing over positions and adding the dust ledger therefore gives
\[
M_{T_a}(\tau^+) = M_{T_a}(\tau^-).
\]
The same argument applies to the other token in the rebalanced pool. Applying the argument to each token gives per-token conservation at the event, and weighting the conserved totals by the fixed reference prices $p_i$ gives the invariance of $V_{\symup{total}}$.

With $m = 1$ the argument is the conservation step of the prior paper's retained-excursion proposition: the lone position's leftover stands in the ledger and is offered to its next mint. \proofqed

\subsection*{Proof of Lemma~\ref{lem:partition}}

The entitlements, the equity, the transfer $\Xi_j$ and the net transfers $\xi_k$ are as in Section~\ref{sec:extinction}.

\begin{proof}
Write $D$ and $D^{\symup{new}}$ for the value of the ledger before and after the event, at the fixed reference prices. By \eqref{eq:avail} and \eqref{eq:newdust}, $L_{\symup{new}}\,g_x + x_{\symup{dust}}' = x_w + x_{\symup{dust}} - \sigma_x$, and likewise for $T_b$, while ledger entries of tokens outside the pool do not move. Multiplying by the prices and adding,
\[
V_j^{\symup{new}} + D^{\symup{new}} = V_j + D - \bar\sigma_j .
\]


(i) If $\Xi_j = 0$ then $G_j^{\symup{new}} = U_j - \bar\sigma_j - V_j^{\symup{new}} \geq 0$ and no other entitlement moves, so the entitlements sum to $D - G_j + G_j^{\symup{new}} = D^{\symup{new}}$. If $\Xi_j > 0$ then $G_j^{\symup{new}} = 0$ and the other entitlements sum to $D - G_j - \Xi_j = D^{\symup{new}}$, which is non-negative because the entries of the ledger are (Theorem~\ref{thm:binding}), so the proportional reduction is well defined. In both cases $\sum_k \xi_k = \Xi_j - \Xi_j = 0$.

(ii) follows by summing (i) over the events of the window.

(iii) With one position $G_1 = D$, so $U_1 - \bar\sigma_1 = V_1^{\symup{new}} + D^{\symup{new}} \geq V_1^{\symup{new}}$ and $\Xi_1 = 0$.

(iv) In the first case $V_j^{\symup{new}} = U_j - \bar\sigma_j - G_j^{\symup{new}}$ and $\xi_j = 0$. In the second $V_j^{\symup{new}} = U_j - \bar\sigma_j + \xi_j$ and $G_j^{\symup{new}} = 0$. Subtracting $V_j = U_j - G_j$ gives the identity in both.
\end{proof}

The transfer accounting alone determines the division into net receivers and net givers. A window has net receivers and net givers exactly when its cumulative net transfers are not all zero. This requires at least one mint in the window to have exceeded the equity of its position. By (iv), a position's value change at its own rebalances differs from the transfer it receives by its swap costs and the change in its own entitlement. A position that leaves its own leftover in the ledger and takes it back later therefore has value changes of both signs and no transfer.

\section{Proofs for Section~\ref{sec:chaos}}\label{app:chaosproof}

\paragraph{Plan of the appendix.}
This appendix compares the model's assumptions with the execution residual measured on chain. It then derives the kernel~\eqref{eq:Khom} from Theorem~\ref{thm:binding}, establishes the invariant set and stationary law (Lemma~\ref{lem:stationary}), and proves Propositions~\ref{thm:chaos2A} and~\ref{thm:chaos2B} under (A1$'$)--(A3). Section~7 of the supplement gives the proofs of Propositions~\ref{cor:hetero} and~\ref{prop:commondir} as modifications to the proof of Proposition~\ref{thm:chaos2A}. It also proves Lemma~\ref{lem:indicators}.

\paragraph{The assumptions against chain data.}
The V9 chain data allow the model's assumptions to be compared with the deployed contract (Table~\ref{tab:residual_chain}). The table covers 74{,}490 of the 82{,}505 range changes in automatic mode from 21~March to 14~October 2026. Events are grouped by the token the corrective swap buys or, when the contract skips the swap, would have bought. The 7{,}560 whose swap executed and could not be measured, and the 455 on which the swap was skipped and the mint was single-sided, are left out. Moments and the first-order law are taken on the events whose residual is measured.

(i)~The model routes the swap at the spot price, without fee or price impact, so that $\sigma_y^{\symup{raw}} = -s^2\sigma_x^{\symup{raw}}$ between its two legs in raw token units. This reduces the swap correction to the single residual $\varepsilon$, measured here against the input that would have left nothing.

(ii)~The residual's law depends on direction, as (A1$'$) allows. When the swap buys token0 the residual is negative on $2.1\%$ of measured events, and when it buys token1 on $53.7\%$. Counting the swaps the contract skipped, at $\varepsilon = -1$, the shares are $4.8\%$ and $56.0\%$.

The residual is not drawn independently of the state. To first order it is the position's liquidity as a share of the pool's other in-range liquidity\cite{ryan2026siphon}, less half the fee, plus a floor-tick term whose sign follows the direction\cite{ryan2026om}. On the chain record used here this law has a median error of $6 \times 10^{-4}$ and the right sign on $98.2\%$ of measured events. The measured residual has median absolute size $0.045$ and second moment about zero $0.069$, or $0.072$ and $0.108$ in the ten portfolios of Table~\ref{tab:donor_absorber}, whose positions are a larger share of their pools. Of measured events, $1\%$ exceed one, so the bound $c_0 K < 1/2$ of (A1$'$) does not hold for the contract's tail.

(iii)~The standing balance has a median value equal to $1.8\%$ of the position's value ($2.8\%$ in the ten portfolios of Table~\ref{tab:donor_absorber}). The leading-order step neglects the balance beside the position's own mismatch, and the step's relative remainder exceeds $0.1$ on $29\%$ of range changes. On $2.4\%$ of them, the standing balance determines the side by reversing the direction of the corrective swap or bringing the mismatch below the contract's threshold. Same-range rebalances, $42\%$ of all rebalances, lie outside the kernel.

(iv)~Same-range rebalances are where the contract departs furthest from (A1$'$). On the $36{,}318$ same-range rebalances of portfolios A to J in the fully covered window, the contract skips the swap on nine events in ten. On $83\%$ of the $36{,}318$, the withdrawn amounts fit the range to within $10^{-9}$ of the withdrawn value. The token left is then set by the standing balances (Section~\ref{sec:model_chain}).
\begin{table}[!htb]
\centering
\caption{The residual on V9 range changes in automatic mode, measured or skipped, by the token the swap buys.}
\label{tab:residual_chain}
\begin{tabular}{@{}lrrr@{}}
\toprule
 & \thd{Buys token}\scn{0} & \thd{Buys token}\scn{1} & \thd{All}  \\
\midrule
Events                                         & 37{,}172 & 37{,}318 & 74{,}490 \\
Swap skipped                                   & 0.027 & 0.049 & 0.038 \\
$\varepsilon < 0$, skipped included            & 0.048 & 0.560 & 0.304 \\
$\varepsilon < 0$, measured                    & 0.021 & 0.537 & 0.277 \\
Median $|\varepsilon|$                         & 0.054 & 0.037 & 0.045 \\
90th and 99th percentiles of $|\varepsilon|$   & 0.32, 1.22 & 0.23, 0.74 & 0.27, 1.02 \\
$\mathbb{E}[\varepsilon]$, $\mathbb{E}[\varepsilon^2]$ & 0.134, 0.098 & 0.004, 0.040 & 0.070, 0.069 \\
First-order law, median error                  & $6.0 \times 10^{-4}$ & $5.5 \times 10^{-4}$ & $5.7 \times 10^{-4}$ \\
First-order law, sign agreement                & 0.983 & 0.980 & 0.982 \\
Median balance over position, $D/V$            & 0.017 & 0.019 & 0.018 \\
Remainder above 0.1                            & 0.27 & 0.31 & 0.29 \\
Balance decides the side                       & 0.018 & 0.030 & 0.024 \\
\bottomrule
\end{tabular}
\end{table}
\subsection*{The finite-$N$ Markov chain}

The state of the homogeneous hub-spoke system is $X^N = (H, B_1, \ldots, B_N) \in [0, b^*]^{N+1}$, where $H = D_{\symup{pool},T_0}$ and $B_k = D_{\symup{pool},T_k}$. The firing pool of each event is uniform on the $N$ pools and independent of the past, as it is when the pools fire at one rate on independent clocks without memory. When $P_k$ fires, draw $(\eta, \varepsilon)$ as in (A1$'$), then update $(H, B_k)$ via the kernel~\eqref{eq:Khom} and leave $B_j$ ($j \neq k$) unchanged. The full chain is exchangeable in its spoke coordinates.

\subsection*{Per-event jump on a homogeneous hub-spoke pool}

Consider a rebalance of pool $P_k = (T_0, T_k)$ in the homogeneous hub-spoke topology, with the hub token at the token0 position and the pre-event balances $H$ and $B_k$. Amounts are written in units in which the pool price is one, $s = 1 + O(w)$, with $w = s_b - s_a$ the width of the range. Take first an exit at the upper bound, $\eta = +1$. The position withdraws $x_w = 0$ and
\[
y_w = L\,(s_b - s_a) = Lw .
\]
The new range is centred on $s = s_b$ with the same width, giving per-unit-$L$ token amounts
\[
g_x = \frac{w}{2}\,\bigl(1 + O(w)\bigr), \qquad g_y = \frac{w}{2}\,\bigl(1 + O(w)\bigr).
\]
The available mint inputs are $H$ of the hub token and $Lw + B_k$ of the spoke token. With the swap routed at the spot price, so that $\sigma_y = -s^2\sigma_x$, the exact swap buys hub token,
\[
\sigma_x^* = -\frac{Lw + B_k - H}{2} + O(w^2).
\]
Take next an exit at the lower bound, $\eta = -1$. The position withdraws $y_w = 0$ and hub token worth $Lw\,s_a/s_b$, which is $Lw$ to the same order. The available mint inputs are $Lw + H$ of the hub token and $B_k$ of the spoke token, and the exact swap sells hub token, $\sigma_x^* = +(Lw + H - B_k)/2 + O(w^2)$. In both cases
\[
\sigma_x^* = -\tfrac{1}{2}\eta A_\eta + O(w^2), \qquad A_\eta = Lw + \eta\,(B_k - H),
\]
and the post-swap mint inputs lie at the \acr{LP}-locus
\[
\hat x = \hat y = \frac{H + Lw + B_k}{2}
\]
to leading order. The imbalance depends on the direction, being $Lw + B_k - H$ in one and $Lw + H - B_k$ in the other.

\subsection*{The execution residual and the slack-side leftover}

The model routes the swap at the spot price, without fee or price impact, so that
\[
\sigma_y^{\symup{raw}} = -s^2\,\sigma_x^{\symup{raw}}
\]
between the swap's two legs in raw token units, with $s$ there the raw sqrt price. Under this assumption the error of the swap is a single residual $\varepsilon$. Write
\[
\sigma_x = \sigma_x^*(1 + \varepsilon), \qquad \sigma_y = \sigma_y^*(1 + \varepsilon), \qquad \varepsilon \sim \nu_\eta.
\]
To leading order the post-swap mint inputs are then
\begin{align*}
\hat x  &=  \tfrac{1}{2}(H + Lw + B_k) + \tfrac{1}{2}\,\eta A_\eta \varepsilon, \\
\hat y  &=  \tfrac{1}{2}(H + Lw + B_k) - \tfrac{1}{2}\,\eta A_\eta \varepsilon,
\end{align*}
both in the common unit, in which the swap gives up as much of one token as it buys of the other, so that
\begin{equation}\label{eq:hatdiff}
\hat x - \hat y  =  \eta A_\eta \varepsilon ,
\end{equation}
which is $c_0 \eta A_\eta \varepsilon$ with $c_0 = 1$.
By Theorem~\ref{thm:binding}, with $g_x$ and $g_y$ equal in common units to leading order, the token supplied in excess is the one left. The hub token is token0, so the leftover is on the hub when $\eta A_\eta \varepsilon > 0$ and on the spoke when $\eta A_\eta \varepsilon < 0$. The spoke binds, that is $B_k' = 0$, if and only if $\eta A_\eta \varepsilon \geq 0$. The slack-side leftover is
\begin{equation}\label{eq:leftovermag}
\text{leftover}  =  |\hat x - \hat y|\,\bigl(1 + O(w)\bigr)  =  c_0\,|A_\eta|\,|\varepsilon|\,\bigl(1 + O(w)\bigr).
\end{equation}
For $A_\eta > 0$ and an empty ledger this is Corollary~\ref{cor:swapcorrected} at the bound of the old range, where $\alpha = 1$ and $A_\eta = Lw$. With a standing balance it is Proposition~\ref{thm:connector} for the same inputs, with the imbalance $A_\eta$ in place of $\alpha(\ell)\,V$. After $\eta = +1$ the swap buys hub token, and $\varepsilon > 0$ buys too much of it and leaves it. After $\eta = -1$ the swap buys the spoke token, and $\varepsilon > 0$ leaves that.

\subsection*{The $\ell^1$ identity, the invariant set and the stationary law}

Write $u^+ = \max(u, 0)$ and $u^- = \max(-u, 0)$. For one mark $(\eta, \varepsilon)$ the kernel~\eqref{eq:Khom} reads $H' = c_0\,(\eta A_\eta\varepsilon)^+$ and $B_k' = c_0\,(\eta A_\eta\varepsilon)^-$. Since $|u^+ - v^+| + |u^- - v^-| = |u - v|$ for real $u$ and $v$, two states under one mark satisfy
\begin{equation}\label{eq:l1}
\begin{aligned}
|H_1' - H_2'| + |B_{k,1}' - B_{k,2}'| &= c_0\,|\varepsilon|\,\bigl|A_{\eta,1} - A_{\eta,2}\bigr| \\
&\leq c_0\,|\varepsilon|\,\bigl(|H_1 - H_2| + |B_{k,1} - B_{k,2}|\bigr),
\end{aligned}
\end{equation}
which is the Lipschitz property that arguments of this kind usually assume. The first step is an equality also when the two copies leave their leftover on different tokens. No total-variation form of~\eqref{eq:l1} holds, because if a residual law has an atom, two states at any positive distance send it to different points.

\emph{Invariant set.} If $H \leq b^*$ and $B_k \leq b^*$ then $|A_\eta| \leq Lw + b^*$ and the leftover is at most $\kappa\,(Lw + b^*) = b^*$, by the definition of $b^*$. So the set on which $H$ and every $B_j$ are at most $b^*$ is invariant, and on it $A_\eta \geq Lw - b^*$, which is positive exactly when $\kappa < 1/2$. The bound is sharp for a residual law that charges its bound. For $\kappa > 1/2$ and such a law, events of one pool alternating in direction with $\varepsilon = K$ carry the hub balance towards $b^* > Lw$, and a spoke with an empty balance then has $A_+ < 0$. The atom of Proposition~\ref{thm:chaos2B}, the moment identities and the property that a perturbation of one balance remains in one balance all rest on this positivity.

\emph{Contraction of the full chain.} Drive two copies of $X^N$ by the same firing pools and marks and let $\|X_1 - X_2\|_1 = |H_1 - H_2| + \sum_j |B_{j,1} - B_{j,2}|$. By~\eqref{eq:l1}, an event of pool $k$ gives $\|X_1' - X_2'\|_1 \leq \|X_1 - X_2\|_1 - (1 - c_0|\varepsilon|)\,\bigl(|H_1 - H_2| + |B_{k,1} - B_{k,2}|\bigr)$, and since $k$ is uniform and the mark is independent of the state,
\begin{equation}\label{eq:contraction}
\mathbb{E}\bigl[\|X_1' - X_2'\|_1 \mid X_1, X_2\bigr]  \leq  \Bigl(1 - \frac{1 - \kappa_1}{N}\Bigr)\,\|X_1 - X_2\|_1 .
\end{equation}
On the invariant set this contraction alone gives Lemma~\ref{lem:stationary}, as the proof below shows.

\emph{Remark (residual laws without bounded support).} For a residual law without bounded support the same contraction holds whenever $\kappa_1 < 1$, and the squared norm $\|X\|_2^2 = H^2 + \sum_k B_k^2$ then satisfies, on the firing pool $k$ and with $m_2 := c_0^2\,\mathbb{E}[\varepsilon^2]$,
\begin{equation}\label{eq:fosterdrift}
\begin{aligned}
\mathbb{E}\bigl[\|X'\|_2^2 - \|X\|_2^2 \mid X, k\bigr] &= c_0^2\,\mathbb{E}\bigl[\varepsilon^2 A_\eta^2\bigr] - H^2 - B_k^2 \\
&\leq -(1 - 2m_2)\,(H^2 + B_k^2) + 2m_2\,(Lw)^2 ,
\end{aligned}
\end{equation}
using $A_\eta^2 \leq 2 (B_k - H)^2 + 2 (Lw)^2$ and $(B_k - H)^2 \leq B_k^2 + H^2$. This is a drift towards a compact set when $m_2 < 1/2$, and $\mathbb{E}[\varepsilon^2]$ here is the second moment about zero and not the variance.

\begin{proof}[Proof of Lemma~\ref{lem:stationary}]
The laws on $[0, b^*]^{N+1}$ with the distance $W_1$ for the $\ell^1$ norm form a complete metric space. Coupling two initial laws optimally and driving both copies by the same firing pools and marks, \eqref{eq:contraction} shows that one event maps this space to itself with Lipschitz constant $1 - (1 - \kappa_1)/N < 1$. The fixed point of that map is the unique stationary law, and the iterates converge to it at the stated rate. Tightness follows from the bounded support.
\end{proof}

\subsection*{Decoupling at rate $1/N$}

Between two consecutive events of one spoke the hub is rewritten about $N$ times by other pools. Consider a Markov chain $\bar B$ on $[0, b^*]$ defined as follows. At each step, draw $\bar H \sim \pi_H$ independently of the past state, draw $(\eta, \varepsilon)$ as before, and apply~\eqref{eq:Khom} with $(\bar B, \bar H)$ in place of $(B_k, H)$. Define likewise a hub chain $\bar H$ fed by independent draws $\bar B \sim \mu$. The pair $(\mu, \pi_H)$ is defined by the McKean--Vlasov system
\begin{equation}\label{eq:McKVfixed}
\mu  =  \operatorname{Stat}\bigl(\widetilde{\mathcal{K}}_w \mid \bar H \sim \pi_H\bigr), \qquad \pi_H  =  \operatorname{Stat}\bigl(\widetilde{\mathcal{K}}_w \mid \bar B \sim \mu\bigr),
\end{equation}
where $\operatorname{Stat}(\cdot)$ takes a kernel parametrised by an environment law to its stationary law. Equivalently, $\mu$ and $\pi_H$ are the laws of $B_k'$ and $H'$ when $(H, B_k) \sim \pi_H \otimes \mu$ and the mark is independent of both. Existence and uniqueness follow from~\eqref{eq:l1}. For two pairs of laws, coupling each coordinate optimally and applying one mark gives $\mathbb{E}|H_1' - H_2'| + \mathbb{E}|B_1' - B_2'| \leq \kappa_1\,\bigl(W_1(\pi_{H,1}, \pi_{H,2}) + W_1(\mu_1, \mu_2)\bigr)$. The map that takes $(\pi_H, \mu)$ to the laws of $(H', B_k')$ is therefore a contraction with factor $\kappa_1 < 1$ in the sum of the two distances. The contraction comes from $\kappa_1 = c_0\,\mathbb{E}|\varepsilon| < 1$. The reset of a balance to zero does not contract by itself, since on the other branch a difference is multiplied by $c_0|\varepsilon|$.

Write $\Gamma$ for the joint law of $(H', B_k')$ in the statement above. Its marginals are $\pi_H$ and $\mu$, and it is not their product, since one of the two balances is zero after every event. Write $\bar\pi^N := \pi_H \otimes \mu^{\otimes N}$. The proof of Proposition~\ref{thm:chaos2A} does not follow the history of a tagged spoke. It compares $\pi^N$ with $\bar\pi^N$ along the chain started from $\bar\pi^N$, and it rests on two properties of the kernel on the invariant set.

\emph{One register.} Call the coordinates of $X^N$ registers, numbered $0$ for the hub and $k$ for spoke $k$. Let $\zeta := c_0 \eta \varepsilon \in [-\kappa, \kappa]$. On the invariant set $A_\eta > 0$, and~\eqref{eq:Khom} reads $H' = \zeta^+ A_\eta$ and $B_k' = \zeta^- A_\eta$, an affine map of $(H, B_k)$ whose coefficients depend on the mark alone. Suppose two states in the invariant set differ by $a$ in the hub register only. After an event of pool $k$, their difference is $-\eta\,|\zeta|\,a$. It lies in the hub register if $\zeta > 0$ and in register $k$ if $\zeta < 0$; if $\zeta = 0$, the states coincide. If they differ by $a$ in register $k$ only, the corresponding difference is $+\eta\,|\zeta|\,a$. An event leaves a difference in another spoke's register unchanged. Both states remain in the invariant set, so the statement iterates. A perturbation of one register is therefore, at every later time, a perturbation of exactly one register, its position, multiplied by a factor whose absolute value is the product of $|\zeta|$ over the events that touched it. An event always touches a perturbation on the hub, and touches one on spoke $j$ when pool $j$ fires. Position and factor depend on the sequence of events and not on the state. This is where $\kappa < 1/2$ is used. Where $A_\eta$ can change sign the destination of the leftover depends on the state, and a perturbation of one register can reach two.

\emph{A defect of second order.} Let $X_0 \sim \bar\pi^N$ and let pool $k$ fire. The pair $(H, B_k)$ then has law $\Gamma$ and the other balances are unchanged, so the event changes only the dependence between two registers whose marginal laws it leaves as they were. Let $(a, c)$ and $(a', c')$ be independent with law $\Gamma$, and write $x[a, c]$ for the state with hub balance $a$, balance $c$ on spoke $k$ and the remaining balances drawn independently from $\mu$, the same draws in each of the four states below. Then $x[a, c]$ and $x[a', c']$ have the law of the state after the event, $x[a, c']$ and $x[a', c]$ have law $\bar\pi^N$, and for every bounded $g$
\begin{equation}\label{eq:mixeddefect}
\mathbb{E}\,g(X_1) - \mathbb{E}\,g(X_0)  =  \tfrac{1}{2}\,\mathbb{E}\bigl[\,g(x[a, c]) + g(x[a', c']) - g(x[a, c']) - g(x[a', c])\,\bigr] .
\end{equation}
Fix $m$, call the hub and spokes $1, \ldots, m$ the observed registers, and let $f$ be a function of the observed balances that is 1-Lipschitz for the $\ell^1$ norm. Take $g(x)$ to be the mean of $f$ after $n$ further events from $x$. The four states in~\eqref{eq:mixeddefect} differ in the hub register and register $k$. Driven by the same events, they differ afterwards only at the positions of those two perturbations. There the differences are $a - a'$ and $c - c'$ multiplied by the factors $\psi^1_n$ and $\psi^2_n$, respectively. If either position is not observed, $f$ takes equal values in pairs and the bracket vanishes. Otherwise the bracket is at most $2\,|\psi^1_n|\,|a - a'|$ and at most $2\,|\psi^2_n|\,|c - c'|$ in absolute value, hence at most $2 b^*\,|\psi^1_n\,\psi^2_n|^{1/2}$.

\begin{table}[!htb]
\centering
\caption{The eight configurations of two perturbations.}
\label{tab:configurations}
\begin{tabular}{@{}llll@{}}
\toprule
 & \thd{Code} & \thd{Where the two perturbations are} & \thd{Both observed}  \\
\midrule
Apart & $\symup{hu}$ & the hub and an unobserved spoke & no \\
 & $\symup{ht}$ & the hub and an observed spoke & yes \\
 & $\symup{uu}$ & two unobserved spokes & no \\
 & $\symup{ut}$ & an unobserved spoke and an observed one & no \\
 & $\symup{tt}$ & two observed spokes & yes \\
\midrule
Together & $\symup{H}$ & the hub & yes \\
 & $\symup{U}$ & an unobserved spoke & no \\
 & $\symup{T}$ & an observed spoke & yes \\
\bottomrule
\end{tabular}
\end{table}

\emph{The two-perturbation sum.} Run two perturbations along one sequence of events, the first from the hub and the second from a spoke chosen uniformly, with positions $r^1_n, r^2_n$ and factors $\psi^1_n, \psi^2_n$ after $n$ events, and let
\begin{equation}\label{eq:twosum}
\mathcal{S}_{N,m}  :=  \mathbb{E}\sum_{n \geq 0}\,\bigl|\psi^1_n\,\psi^2_n\bigr|^{1/2}\;\mathbf{1}\bigl\{r^1_n \leq m,\ r^2_n \leq m\bigr\} .
\end{equation}


The sum collects, with their weights, the events in which both perturbations lie in observed registers. Two perturbations in one register are touched by the same events and stay together from then on. Up to the symmetry between pools the pair is in one of the eight configurations of Table~\ref{tab:configurations}, five in which the two are apart and three in which they are together.
One event multiplies $|\psi^1\psi^2|^{1/2}$ by $|\zeta|$ if it touches both perturbations, by $|\zeta|^{1/2}$ if it touches one, and by $1$ if it touches neither. Let $\gamma_\pm := \mathbb{E}[\zeta^\pm]$ and $\varsigma_\pm := \mathbb{E}[(\zeta^\pm)^{1/2}]$, so that $\gamma_+ + \gamma_- = \kappa_1$ and $\varsigma_+ + \varsigma_- = \varsigma$, with $\kappa_1 \leq \varsigma \leq \sqrt{\kappa}$. For a function $\varphi$ of the configuration, the mean of the multiplier times the value of $\varphi$ after one event is
\begingroup\allowdisplaybreaks
\begin{align*}
(\mathcal{Q}\varphi)_{\symup{H}} &= \gamma_+\,\varphi_{\symup{H}} + \tfrac{m}{N}\,\gamma_-\,\varphi_{\symup{T}} + \tfrac{N-m}{N}\,\gamma_-\,\varphi_{\symup{U}}, \\
(\mathcal{Q}\varphi)_{\symup{U}} &= \bigl(1 - \tfrac{1}{N}\bigr)\varphi_{\symup{U}} + \tfrac{1}{N}\bigl(\gamma_+\,\varphi_{\symup{H}} + \gamma_-\,\varphi_{\symup{U}}\bigr), \\
(\mathcal{Q}\varphi)_{\symup{T}} &= \bigl(1 - \tfrac{1}{N}\bigr)\varphi_{\symup{T}} + \tfrac{1}{N}\bigl(\gamma_+\,\varphi_{\symup{H}} + \gamma_-\,\varphi_{\symup{T}}\bigr), \\
(\mathcal{Q}\varphi)_{\symup{hu}} &= \tfrac{1}{N}\bigl(\gamma_+\,\varphi_{\symup{H}} + \gamma_-\,\varphi_{\symup{U}}\bigr) + \tfrac{N-1}{N}\,\varsigma_+\,\varphi_{\symup{hu}} + \tfrac{m}{N}\,\varsigma_-\,\varphi_{\symup{ut}} + \tfrac{N-m-1}{N}\,\varsigma_-\,\varphi_{\symup{uu}}, \\
(\mathcal{Q}\varphi)_{\symup{ht}} &= \tfrac{1}{N}\bigl(\gamma_+\,\varphi_{\symup{H}} + \gamma_-\,\varphi_{\symup{T}}\bigr) + \tfrac{N-1}{N}\,\varsigma_+\,\varphi_{\symup{ht}} + \tfrac{m-1}{N}\,\varsigma_-\,\varphi_{\symup{tt}} + \tfrac{N-m}{N}\,\varsigma_-\,\varphi_{\symup{ut}}, \\
(\mathcal{Q}\varphi)_{\symup{uu}} &= \bigl(1 - \tfrac{2}{N}\bigr)\varphi_{\symup{uu}} + \tfrac{2}{N}\bigl(\varsigma_+\,\varphi_{\symup{hu}} + \varsigma_-\,\varphi_{\symup{uu}}\bigr), \\
(\mathcal{Q}\varphi)_{\symup{ut}} &= \bigl(1 - \tfrac{2}{N}\bigr)\varphi_{\symup{ut}} + \tfrac{1}{N}\bigl(\varsigma_+\,\varphi_{\symup{ht}} + \varsigma_+\,\varphi_{\symup{hu}} + 2 \varsigma_-\,\varphi_{\symup{ut}}\bigr), \\
(\mathcal{Q}\varphi)_{\symup{tt}} &= \bigl(1 - \tfrac{2}{N}\bigr)\varphi_{\symup{tt}} + \tfrac{2}{N}\bigl(\varsigma_+\,\varphi_{\symup{ht}} + \varsigma_-\,\varphi_{\symup{tt}}\bigr).
\end{align*}
\endgroup
In each line the event either touches no perturbation, or touches both with the pair moving to the hub when $\zeta > 0$ and to the firing spoke when $\zeta < 0$, or touches one, which moves in the same way. If $\varphi \geq 0$ satisfies $\varphi \geq \mathbf{1}_{\symup{obs}} + \mathcal{Q}\varphi$, with $\mathbf{1}_{\symup{obs}}$ the indicator of the four observed configurations, then by induction on the number of events the mean of the sum in~\eqref{eq:twosum} started from a configuration is at most the value of $\varphi$ there. Such a function is
\begin{gather*}
\varphi_{\symup{H}} = \varphi_{\symup{U}} = \Lambda, \qquad \varphi_{\symup{T}} = \Lambda + \frac{N}{1 - \gamma_-}, \qquad \varphi_{\symup{hu}} = \varphi_{\symup{uu}} = \min(\Lambda_0, \Lambda_1), \\
\varphi_{\symup{ht}} = \varphi_{\symup{ut}} = \Lambda_1, \qquad \varphi_{\symup{tt}} = \Lambda_1 + \frac{N}{2 (1 - \varsigma_-)},
\end{gather*}
with
\begin{gather*}
\Lambda := \frac{1}{1 - \kappa_1}\Bigl(1 + \frac{\gamma_-\,m}{1 - \gamma_-}\Bigr), \qquad
\Lambda_0 := \frac{\kappa_1 \Lambda + \varsigma_-\,m\,\Lambda_1}{(1 - \varsigma)\,N}, \\
\Lambda_1 := \frac{1}{1 - \varsigma}\Bigl(\frac{1}{1 - \gamma_-} + \frac{\kappa_1 \Lambda}{N} + \frac{\varsigma_-\,(m - 1)}{2 (1 - \varsigma_-)}\Bigr) .
\end{gather*}
In configuration $\symup{H}$ the inequality is an equality, by the choice of $\Lambda$. In $\symup{ht}$ it holds by the choice of $\Lambda_1$, after $(N-1)\varsigma_+/N \leq \varsigma_+$ and $(m-1)/N + (N-m)/N \leq 1$. In $\symup{hu}$ it holds by the choice of $\Lambda_0$ when the minimum is $\Lambda_0$, and by $(1 - \varsigma)\Lambda_1 \geq \kappa_1\Lambda/N$ when the minimum is $\Lambda_1$. In the other five it follows from $\gamma_+ \leq 1 - \gamma_-$, $\varsigma_+ \leq 1 - \varsigma_-$ and $\min(\Lambda_0, \Lambda_1) \leq \Lambda_1$. The second perturbation starts on an observed spoke with probability $m/N$, so
\begin{equation}\label{eq:twosumbound}
\mathcal{S}_{N,m}  \leq  \frac{m}{N}\,\Lambda_1 + \Lambda_0  =  \frac{\kappa_1 \Lambda + (1 - \varsigma_+)\,m\,\Lambda_1}{(1 - \varsigma)\,N}  \leq  \frac{m\,(m + 5)}{2 (1 - \varsigma)^4\,N} ,
\end{equation}
the last step by $\Lambda \leq m/(1 - \varsigma)^2$ and $\Lambda_1 \leq (m + 3)/\bigl(2 (1 - \varsigma)^3\bigr)$, which follow from $\kappa_1 \leq \varsigma$ and $m \leq N$.

\begin{proof}[Proof of Proposition~\ref{thm:chaos2A}]
Let $f$ be as above. By Lemma~\ref{lem:stationary} the mean of $f$ after $t$ events from $\bar\pi^N$ tends to $\pi^N f$ as $t \to \infty$. Its difference from $\bar\pi^N f$ is the sum over $n < t$ of the change made by inserting one event before $n$ further ones, which is~\eqref{eq:mixeddefect} with the firing pool $k$ uniform. By the bound that follows~\eqref{eq:mixeddefect}, with the two perturbations started from the hub and from spoke $k$, the term for $n$ is at most $b^*$ times the summand of~\eqref{eq:twosum} in absolute value. Hence $|\pi^N f - \bar\pi^N f| \leq b^*\,\mathcal{S}_{N,m}$ for every such $f$, which by the duality of Kantorovich and Rubinstein is the bound on $W_1\bigl(\pi^N_{(0,m)}, \pi_H \otimes \mu^{\otimes m}\bigr)$ once~\eqref{eq:twosumbound} is used. Projecting a coupling onto the spoke coordinates gives the bound for $\pi^N_{(m)}$, and $W_1$ metrises weak convergence on a compact set.

For the moments, let $X \sim \pi^N$ and let $X'$ be the state after one event, which has the same law. The kernel is affine with coefficients independent of the state, so with $\chi_\pm$ as in the Remark after the proposition the means satisfy $\bar h = \gamma_+ Lw + \chi_+\,(\bar b - \bar h)$ and $\bar b = \gamma_- Lw + \chi_-\,(\bar b - \bar h)$ at every $N$. The means of $\pi_H$ and $\mu$ satisfy the same two equations, whose determinant $1 + \chi_+ - \chi_-$ is positive. Next, $H'B_1' = 0$ when pool $1$ fires and $H'B_1' = \zeta^+ A_\eta\,B_1$ when another pool $k$ does, and $B_1'B_2$ changes only when pool $1$ fires. With $x_3 := \mathbb{E}[H B_1]$ and $x_4 := \mathbb{E}[B_1 B_2]$, stationarity and the exchangeability of the spokes give
\[
x_3 = \tfrac{N-1}{N}\bigl(\gamma_+ Lw\,\bar b + \chi_+\,(x_4 - x_3)\bigr), \qquad x_4 = \gamma_- Lw\,\bar b + \chi_-\,(x_4 - x_3) .
\]
At $N = 1$ only the first equation is present and reads $x_3 = 0$. The products $\bar h\,\bar b$ and $\bar b^2$ solve the same system without the factor $(N-1)/N$. Subtracting the two systems gives the covariance of the Remark, together with
\[
\operatorname{Cov}_{\pi^N}(B_1, B_2) = \frac{\bar h\,\bar b}{N}\cdot\frac{\chi_-}{1 - \chi_- + (1 - 1/N)\,\chi_+} .
\]
Since $(h, b) \mapsto h\,b/b^*$ is 1-Lipschitz on $[0, b^*]^2$, the distance $W_1\bigl(\pi^N_{(0,1)}, \pi_H \otimes \mu\bigr)$ is at least $\lvert\operatorname{Cov}_{\pi^N}(H, B_1)\rvert/b^* \geq \bar h\,\bar b/(3 b^*N)$. The same computation for $\mathbb{E}[B_1^2]$ gives
\[
\mathbb{E}_{\pi^N}[B_1^2] - \int b^2\,\mu(db) = -\frac{2 \mathbb{E}[(\zeta^-)^2]}{1 - \mathbb{E}[\zeta^2]}\,\operatorname{Cov}_{\pi^N}(H, B_1) ,
\]
so when $\bar h\,\bar b > 0$ the law of a single spoke is also at a distance of order $1/N$ from $\mu$, and no faster rate holds for either. The means are positive exactly when the corresponding sign of $\zeta$ has positive probability. If leftovers fall on one token only, one coordinate of $\Gamma$ is almost surely zero, so $\Gamma$ is the product of its marginals. An event of any pool then carries $\bar\pi^N$ to itself, and $\pi^N = \bar\pi^N$ by the uniqueness in Lemma~\ref{lem:stationary}.
\end{proof}

\paragraph{The hub balance at a spoke's own event.} The firing pool is independent of the state, so in stationarity the hub balance a spoke meets at its own event and that spoke's standing balance have joint law $\pi^N_{(0,1)}$. By Proposition~\ref{thm:chaos2A} with $m = 1$ this law is within $3 b^*/\bigl((1 - \varsigma)^4 N\bigr)$ of $\pi_H \otimes \mu$ in $W_1$. The covariance in the Remark shows that the two are negatively correlated at exactly that order, because with probability $1/N$ the spoke fired at the previous event and then one of the two balances is zero. This is a statement about the pair at one time. The proof does not use, and does not establish, closeness of the conditional law of the hub balance given the whole history of the spoke. The supplement measures it given part of that history in the simulated chain (its Table~S2).
\subsection*{The spoke marginal, exact atom and slack-side approximation}

\begin{proof}[Proof of Proposition~\ref{thm:chaos2B}]
After an event of pool $P_k$ the spoke balance is $c_0\,(\eta A_\eta\varepsilon)^-$, and $A_\eta > 0$ on the invariant set, so it is zero exactly when $\eta \varepsilon \geq 0$. The mark is independent of the state, so this has probability $q$ at every event of the pool, and the spoke balance does not change between events of its own pool. The stationary law of $B_k$ therefore has an atom at zero of mass exactly $q$, for every $N$. On $\eta \varepsilon < 0$ the balance is $c_0 A_\eta|\varepsilon|$ with $|A_\eta - Lw| = |B_k - H| \leq b^*$, so it differs from $c_0 Lw|\varepsilon|$ by at most $c_0 b^*\,|\varepsilon|$ under the same mark, which is the bound on $W_1(\rho^N_+, \rho_+)$. The limit chain $\bar B$ has the same kernel, with the hub balance drawn from $\pi_H$ on $[0, b^*]$, and the same argument applies.
\end{proof}


\subsection*{Open: beyond the invariant set}

Proposition~\ref{thm:chaos2A} holds under (A1$'$)--(A3) with no further assumption, and its proof shows where each part of the model is used. The bound $\kappa < 1/2$ enters through the one-register property alone. For a residual law with $\kappa_1 < 1$ whose bound exceeds one half, or which is unbounded, the chain still contracts in $W_1$ by~\eqref{eq:contraction} and has a unique stationary law. The imbalance can change sign, however, so the destination of a leftover depends on the state and the argument above does not apply. The equality of the spokes and the uniform firing enter only through the single fixed point~\eqref{eq:McKVfixed} and the count of configurations behind~\eqref{eq:twosumbound}. Proposition~\ref{cor:hetero} removes both, at the price of the condition~\eqref{eq:heterobox} on the sizes. That condition is the form the bound on the residual takes for unequal spokes, and a spread of sizes wider than it allows is open for the same reason as a bound above one half.

Five other questions remain. The statement is in $W_1$, and no total-variation form is available, for the reason given after~\eqref{eq:l1}. Neither the constant in~\eqref{eq:twosumbound} nor its quadratic growth in $m$ is claimed to be sharp, and solving the linear system for the sum $\mathcal{S}_{N,m}$ itself, in at most eight unknowns, gives the sharper bound. A residual that depends on the state, or an environment that reacts to it, removes the independence on which~\eqref{eq:mixeddefect} rests, and none of the propositions says anything about them. Proposition~\ref{prop:commondir} is for equal spokes, and unequal spokes under a common environment are not treated. For an environment that persists for a number of events of the order of $N$, the size of the unconditional covariance of two balances is not determined here.

The count in~\eqref{eq:twosum} is the forward form of the interaction-graph bound of Graham and M\'el\'eard\cite{graham1993} for loss networks with alternate routing. In that setting, two tagged components are dependent to the extent that their interaction histories meet. Here, two perturbations matter to the extent that both reach the observed registers. An $\ell^1$ Lipschitz property of the kind that Andreis, Dai Pra and Fischer\cite{andreis2018} assume for systems with simultaneous jumps holds here as~\eqref{eq:l1}, and their route is a candidate for the case in which the imbalance changes sign, where a contraction is available and the one-register property is not.
\section{Methodology of the empirical validation}\label{app:methodology}

\subsection*{Contracts and schemas}

Four \acr{PM} contracts in the framework's class are studied. \acr{PM}~V6 was active from 1 to 5 March 2026, V7 from 6 to 13 March, V8 from 13 to 21 March, and V9 from 21 March to 14 October 2026. Table~\ref{tab:cross_version} gives the dates of their rebalances, all of which fall within these periods. V6 and V7 share a 2-tuple $(\text{depositor}, \text{token})$ dust mapping. V8 added the portfolio dimension, giving the 3-tuple $(\text{depositor}, \text{portfolioId}, \text{token})$ inherited by V9. The changes from V8 to V9 were internal and did not modify event signatures.

In the 49 days from deployment to block 45{,}746{,}143 on 8 May 2026, the dataset cutoff in Section~\ref{sec:empirical}, the V9 event lake holds 126{,}339 \texttt{Rebalanced} events and 143{,}991 \texttt{DustCredited} events. Over its full life, through 14 October 2026, V9 made 280{,}097 rebalances in 70 pools. V6, V7 and V8 are read from their pool and \texttt{DustCredited} events.

\subsection*{Conservation check and mint reconstruction}

\emph{Conservation check.} For each rebalance the conservation form takes its terms from the events of the transaction.
\begin{center}
\begin{tabular}{@{}l@{\quad}>{\raggedright\arraybackslash}p{0.68\linewidth}@{}}
\toprule
\thd{Quantity} & \thd{Source}  \\
\midrule
$(x_w, y_w)$ & its Pool.Burn events, with any fee collected in its Pool.Collect events \\
$(\sigma_x, \sigma_y)$ & its Pool.Swap event on the rebalance pool, or the contract's \texttt{ExternalSwapExecuted} event when the swap is routed elsewhere \\
$(x_{\symup{mint}}, y_{\symup{mint}})$ & its Pool.Mint event \\
Standing balance & the ledger replayed in chain order from \texttt{DustCredited} and \texttt{DustWithdrawn} events, per (depositor, portfolio, token) \\
\bottomrule
\end{tabular}
\end{center}
\texttt{DustCredited} carries the new balance and not an increment. A rebalance sets the portfolio's balance in the pool's two tokens to what it credits, and a token that receives no credit is set to zero. The expression
\[
(x_{\symup{dust}}', y_{\symup{dust}}') = (x_w + x_{\symup{dust}} - \sigma_x - x_{\symup{mint}},\; y_w + y_{\symup{dust}} - \sigma_y - y_{\symup{mint}})
\]
is compared with these amounts in integer token units. The replayed ledger matches every \texttt{DustWithdrawn} amount. For 600 sampled token balances, it also matches an archive read of \texttt{dustBalance} at the preceding block. On V6 to V8 the event data hold no Collect events, and a rebalance closes when the fee leg the identity implies is zero.

\emph{Mint reconstruction.} The steps of Table~\ref{tab:modelift} solve the mint in real arithmetic. For new range $[s_a', s_b']$ and sqrt price $s$, the inputs required per unit of liquidity are $(g_x, g_y) = (1/s - 1/s_b',\; s - s_a')$ inside the range, $(1/s_a' - 1/s_b',\; 0)$ below it, and $(0,\; s_b' - s_a')$ above it. Given inputs $(\hat x, \hat y)$, the mint creates $L_{\symup{new}} = \min(\hat x/g_x,\; \hat y/g_y)$ units of liquidity and leaves $(\hat x - L_{\symup{new}}\,g_x,\; \hat y - L_{\symup{new}}\,g_y)$.

\emph{Prices.} The pre-swap price is read from the burn when the old position was in range, from the pool's last swap before the transaction otherwise, and from the mint on swap-free events. When two estimates are available, their difference has a median of about one part in $10^{12}$. The post-swap price is the \texttt{sqrtPriceX96} of the rebalance's Pool.Swap event. The leftover's token is the token of larger value at the pre-swap price.

\emph{Rounding.} For swap-free events, steps~(ii) to~(iv) coincide because no swap moves the price. These steps identify the wrong leftover token on 5.1\% of range changes and 2.7\% of same-range events. All but three of those misses are leftovers smaller than $10^{-12}$ of the value withdrawn and held, which is the scale of the pool's integer rounding. A reconstruction in real arithmetic does not reliably tell which token keeps a balance at that scale. On the $117{,}414$ events whose leftover exceeds $10^{-9}$ of that value, step~(iv) names the token on all but one. On 477 of the swap-free range changes the withdrawal fits the new range exactly and leaves no residual of its own.
\subsection*{Data sources}

The empirical validation in Section~\ref{sec:empirical} uses four event sources. The first is the \acr{PM}~V9 event lake, which contains all contract events on Base from deployment through the contract's last rebalance on 14 October 2026. The second is the set of pool Burn, Collect, Swap and Mint events for every \acr{PM}~V9 rebalance transaction through 14 May 2026. The third is the pool events for every \acr{PM}~V6, V7 and V8 rebalance. For V9 after 14 May, the fourth is the position-manager events from the 17 pools covered by a position-level dataset. No field of the off-chain rebalancing engine's log enters Tables~\ref{tab:cross_version}, \ref{tab:results} or~\ref{tab:modelift}, apart from the swap leg of 504 rebalances in the whole-life row of Table~\ref{tab:cross_version}. Across the four contracts 90{,}272 of the 152{,}719 rebalances carry a swap, which is 59\%, and 96\% of those that change the range do.

The tables are built per rebalance. The integer terms of the identity are assembled from those events. The ledger is replayed per (depositor, portfolio, token) in chain order, with each rebalance's credits applied before the next event. Each rebalance is labelled a range change or a same-range rebalance by comparing the ticks of the burned and the minted position.

Two auxiliary indexers support the validation. One scans the \acr{PM}~V9 transactions to block 44{,}277{,}369 for administrative calls and finds none that changes the ledger outside the events used here. The only administrative calls, $27$ in number, are deposits already recorded as \texttt{DustCredited} events. The other records the 73{,}350 Pool.Swap events issued by \acr{PM}~V9 across its 52 pools to block 45{,}746{,}143, one for each rebalance that swapped. These provide the swap-correction inputs $(\sigma_x, \sigma_y)$.



\subsection*{Simulations and replays}

Sections~\ref{sec:model_tests} and~\ref{sec:model_chain} rest on three computations. The model chain of Section~\ref{sec:chaos_numerical} applies the kernel~\eqref{eq:Khom} in double precision and is checked against an exact rational implementation of the same kernel. The contract of Section~\ref{sec:model_contract} is the exact integer simulation of the prior paper\cite{ryan2026siphon}, which reproduces the arithmetic of the pool and of the manager. It is extended here to 63 spokes and checked bit for bit in its integer state against a reference implementation.

The replays of Section~\ref{sec:model_chain} pass each portfolio's recorded events through the mint of Theorem~\ref{thm:binding} in exact integer arithmetic. With every input as recorded, they reproduce the liquidity, mint amounts, balances and zero indicators for all $84{,}615$ rebalances of portfolios A to J in the fully covered window. The mint price is not recorded for $34{,}244$ events and is inferred from the mint itself. For $24{,}033$ of these, the price is solved as the value inside the new range at which the mint returns the recorded amounts exactly.

The second set of replays carries the ledger through every same-range rebalance whose swap was skipped. In these replays, the principal is the burn, at the price of the event, of the liquidity the pool last minted.

\subsection*{What was specified in advance}

The table records when each analysis was specified and what results had been seen by then.
\begin{center}
\begin{tabular}{@{}>{\raggedright\arraybackslash}p{0.36\linewidth}@{\quad}>{\raggedright\arraybackslash}p{0.58\linewidth}@{}}
\toprule
\thd{Analysis} & \thd{Specified}  \\
\midrule
Tables of Sections~\ref{sec:cross_version} to~\ref{sec:framework_predictions} & before they were computed \\
Simulated and replayed quantities, their estimators and decision rules & before the scored runs \\
First set of replays, the steps that shuffle the same-range outcomes & each step before it was run, after the first results of that set had been seen \\
Second set of replays & before any replay of this set was run on the record, after the first set had been seen and after a check on the record of what a same-range rebalance can change. That check shaped the declaration. The decision rule was fixed on simulated records \\
Descriptive additions & after first outputs, or outside the declared runs \\
\bottomrule
\end{tabular}
\end{center}
The descriptive additions are the account of the misses of step~(iv) and of the criterion, the remark in Section~\ref{sec:model_chain} on which entries differ, and the entry-by-entry comparison of the atom on range changes. The others are the longer run measuring dependence under the Gaussian law (supplement, Section~1), the averages of the second replay set for pairs dependent on chain and for the others, and the comparison of simulated means and covariance with their exact values.

\section{Foundry verification suite}\label{app:foundry}

The suite of Section~\ref{sec:foundry} has 24 tests in 8 contracts. Twelve are network-free. Twelve run on a fork of Base at block~\texttt{43{,}175{,}000}, six against each of two unmodified V3 position managers.\footnote{\raggedright On Base. Aerodrome Slipstream, \texttt{0x827922686190790b37229fd06084350E74485b72}. Uniswap V3, \texttt{0x03a520b32C04BF3bEEf7BEb72E919cf822Ed34f1}.}

\subsection*{Test--theorem mapping}

\begin{table}[H]
\centering
\caption{Foundry test mapping.}
\label{tab:foundry-mapping}
\small
\begin{tabular}{@{}lccL{5.5cm}@{}}
\toprule
\thd{Test contract} & \thd{Tests} & \thd{Network} & \thd{Statement}  \\
\midrule
\texttt{MasterEquationT1Mock} & 8 & none & Theorem~\ref{thm:binding} \\
\texttt{MasterEquationP1Mock} & 4 & none & Proposition~\ref{thm:connector}, Corollaries~\ref{cor:swapfree},~\ref{cor:swapcorrected} \\
\texttt{MasterEquationT1Fork} & 4 & Base fork & Theorem~\ref{thm:binding}, Slipstream \\
\texttt{MasterEquationP1Fork} & 1 & Base fork & Proposition~\ref{thm:connector}, Slipstream \\
\texttt{MasterEquationT2Fork} & 1 & Base fork & Theorem~\ref{thm:extinction}, Slipstream \\
\texttt{MasterEquationT1ForkUniV3} & 4 & Base fork & Theorem~\ref{thm:binding}, Uniswap V3 \\
\texttt{MasterEquationP1ForkUniV3} & 1 & Base fork & Proposition~\ref{thm:connector}, Uniswap V3 \\
\texttt{MasterEquationT2ForkUniV3} & 1 & Base fork & Theorem~\ref{thm:extinction}, Uniswap V3 \\
\bottomrule
\end{tabular}
\end{table}

The two reference contracts differ only in how they key pools. A test library implements the integer arithmetic of a V3 mint and burn. It computes the liquidity a position manager assigns to the offered amounts, the amounts the pool takes (rounded up), and the amounts returned by a burn (rounded down). The fork tests of Theorem~\ref{thm:binding} use the library to compute predictions and compare them with the contract results for equality. The network-free tests use the library as the pool.

\end{document}
